\documentclass[12pt,a4paper]{article}
\usepackage[UTF8]{ctex}
\usepackage{geometry}
\geometry{left=2.5cm, right=2.5cm, top=2.5cm, bottom=2.5cm}
\usepackage{fancyhdr}         
\usepackage{graphicx}
\usepackage{booktabs}
\usepackage{tabularx}     % 自动换行表格
\usepackage{amsmath} 
\usepackage{enumitem}
\usepackage{tikz}
\usepackage{mhchem}  
\usepackage{makecell} 
\usepackage{multirow} 

\usetikzlibrary{patterns}
\usetikzlibrary{arrows.meta, positioning, shapes.geometric}

\usepackage[dvipsnames]{xcolor}
\usepackage{forest}

\pagestyle{fancy}
\fancyhf{}
\fancyhead[C]{地球科学概论}
\fancyhead[L]{\leftmark}
\fancyhead[R]{\rightmark}
\fancyfoot[C]{\thepage}
\renewcommand{\headrulewidth}{0.4pt}
\renewcommand{\footrulewidth}{0pt}

\title{地球与空间科学技术概论(个人笔记)}
\author{编者: 王宇浩\\ \text{PB2507****} \\ \text{中国科学技术大学地空学院}}
\date{2025 年 12 月 24 日}

\begin{document}
	\maketitle
	\section*{前言}
	\textbf{考前} 本资料由笔者为准备地球与空间科学技术概论期末考试（12月25日）整理得到，耗费巨量时间精力。同时也是我第一次用\LaTeX 纯手敲这么多内容，也是很好的练习。
	
	每一章都对应一个上百页的PPT，由于“地球物理地震灾害与预测”一章课件文件过期，放在最后一章，是根据大纲拟写的,仅供参考。目录做的很细致是为了方便开卷考试的检索。
	
	在此留念我的第一个疯狂期末周（*＾-＾*）
	
	\textbf{考后} 努力没有白费，成功吃到了40\%的优秀率，不枉我花费大量时间。
	\tableofcontents
	\newpage
	\section{地球-独特的行星}
	\subsection{地球科学概论}
	\subsubsection{地幔对流}
	\textbf{定义}
	
	地幔对流是地球深部热能向地表传播的主要机制，是大陆漂移、造山运动、火山和地震活动等的内在动力.
	
	\textbf{意义}
	
	研究地幔对流的发生发展对理解当今岩石圈的热结构及其应力分布具有重要意义
	
	1)利用数值模拟的方法，建立了新的地幔岩石圈下小尺度地幔对流发生发展的规律
	
	2)研究了地幔对流与地表热流、地形等地球物理观测之间的关系
	
	3)首次详细研究了板块运动速度、软流圈的厚度、破碎带等因数对小尺度地幔对流的发生与发展的影响
	\subsubsection{对流地幔的化学演化}
	1)对流地幔内部存在各种尺度的空间和时间不均匀性
	
	2)这种不均匀分布可长期存在，其平均年龄大约1.8 Ga，与地球化学观测相当
	
	3)考虑榴辉岩相变的影响，可在核幔边界上方产生类似非洲下面的化学异常体；榴辉岩相变也为OIB和MORB源区的地球化学差异提供解释
	
	\subsubsection{火山源区岩浆系统的空间特征}
	
	方法：利用大地电磁成像
	
	1)首次获得了五大连池火山区尾山火山下方20公里深的三维高分辨率电阻率结构
	
	2)精细刻画尾山火山的地壳岩浆系统分布形态及估算了岩浆囊的部分熔融程度
	\subsection{地震与月震简介}
	\subsubsection{地震形成机理}
	\includegraphics[width=0.8\textwidth]{yupian1.png} 
	% image_name 是你的图片文件名（无需后缀，如 .jpg, .png, .pdf）
	% width=0.8\textwidth 表示图片宽度占行宽的 80%
	\begin{table}[htbp]
		\centering
		\caption{中英文对照表}
		\label{tab:geological_terms}
		\begin{tabular}{ll} % 两列左对齐
			\toprule
			英文原术语 (Original Term) & 中文翻译 (Chinese Translation) \\
			\midrule
			Japan trench & 日本海沟 \\
			Crust / Crustal EQs & 地壳 / 地壳地震 \\
			Moho & 莫霍面 \\
			Mantle wedge & 地幔楔 \\
			Partial melting / Mantle upwelling & 部分熔融 / 地幔上涌 \\
			Low-F EQs & 低频地震 \\
			Megathrust EQs & 巨型逆冲地震 \\
			M$>$8 Normal faulting & M$>$8 正断层 \\
			Intraslab EQs (Double-seismic zone) & 板内地震 (双地震带) \\
			Pacific plate & 太平洋板块 \\
			Subducting crust / Subducting mantle & 俯冲地壳 / 俯冲地幔 \\
			Fluids & 流体 \\
			\bottomrule
		\end{tabular}
	\end{table}
	
	\subsubsection{地震探测: 光纤测地震}
	\textbf{定义}
	研制分布式光纤声波/振动传感（DAS）系统，
	该设备使用现有通讯光缆可进行地震监测、
	地质灾害预测、地下结构成像、城市地下空
	间探测等。
	
	原理：从站房发射激光脉冲进入光纤，当地震波引起地面微震时，光纤会发生微小的形变。利用光在光纤中产生的“瑞利散射”回传信号，通过解调相位变化，就能精准还原光纤沿线每一米处的震动情况。
	\textbf{实例}
	1)DAS设备自2021年6月部署在合肥紫蓬山进
	行连续观测
	
	2)已成功监测到定远县2.3级地震（6月4日）、
	宣城2.7级地震（7月22日）、菲律宾6.6级
	地震（7月24日）、台湾宜兰县5.8级地震
	（8月5日）等一些区域和全球地震
	\subsubsection{月球科研站：光纤地震仪测月震?}
	在月球科研站（如中俄联合发起的ILRS）的建设规划中，利用光纤地震仪（特别是基于DAS分布式光纤传感技术）探测月震是当前深空探测领域的前沿热点。
	
	传统的登月地震仪（如阿波罗时代的被动地震实验包）是“点”式探测，而光纤地震仪可以将几公里甚至几十公里的光纤变成成千上万个“感应点”，这对于揭示月球内部结构具有革命性意义。
	
	1)\textbf{显著优势}
	
	极高的空间分辨率：月球地壳（月壳）对地震波的散射极强（因为月岩极其干燥且破碎），传统的几个孤立探测点很难精准定位震源。光纤形成的“密集台阵”可以像CT扫描一样清晰地刻画月球内部结构。
	
	超轻量化与易部署：发射几十公斤重的传统传感器成本极高，而光纤极轻、耐卷绕。未来可以通过月球车在站区周围铺设几公里的光纤网，实现大范围监控。
	
	抗恶劣环境：月球表面存在强辐射和极端温差（-170℃到120℃）。光纤本身是无源的（不需要供电），可以埋入月壤中，避开地表温差影响，且抗电磁干扰能力强。
	
	多功能一体化：光纤不仅能测地震，还能同时监测：
	
	科研站结构健康：监测月球基地建筑是否有裂缝或受损。
	
	陨石撞击监控：实时感知微陨石撞击月面的位置。
	
	资源勘探：通过主动人工地震手段，探测月壤下的月冰或空洞（熔岩管道）。
	
	2)\textbf{面临的挑战}
	
	耦合问题：在地球上，光纤通常埋在土壤里或用水泥固定。月球由于是高真空，月壤（Regolith）极其松散，如何确保光纤与月球表面“紧密贴合”以准确捕捉震动波，是工程难点。
	
	极端温差下的应变：月夜和月昼的剧烈温差会导致光纤热胀冷缩，这产生的信号扰动远大于月震信号，需要极其复杂的算法进行信号剥离。
	
	激光解调仪的寿命：光纤虽耐用，但负责发射和接收激光的“解调仪”含有精密光学元件，需在月球背面或地表下进行严密的辐射防护和热控。
	\subsection{战略性关键金属矿产资源}
	美国：2017年《美国的关键矿产资源—经济和环境地质及未来供应展望》27种；
	
	欧洲：2018年《关键原材料和循环经济》43种；
	
	中国：
	
	稀有金属（Li、Be、Rb、Cs、Nb、Ta、Zr、Hf、W等）
	
	稀土金属（La、Ce、Pr、Nd、Sm、Eu、Gd、Tb、Dy、Ho、Er、Tm、Yb、Lu、Sc、Y）
	
	稀散金属（Ga、Ge、Se、Cd、In、Te、Re、Tl）稀少稀贵金属（PEG、Cr、Co等）
	\subsection{地球系统科学}
	\textbf{内涵}
	
	电磁圈+ 大气圈+ 水圈+ 冰冻圈+ 生物圈+ 岩石圈
	
	\textbf{研究方法：交叉融合}
	
	地球物理学（空间物理、固体地球物理）;地质学（构造地质学、矿物岩石学矿床学、古生物学）;地球化学/天体化学;环境科学;数学、物理、化学、信息科学、生命科学等.
	
	\subsubsection{地球的独特性}
	
	\textbf{地球——人类已知的唯一一颗具有板块运动的行星}
	
	板块构造理论是二十世纪最伟大的科学理论之一.
	
	根据海岸线形态、岩石类型、地质构造、生物化石等.
	
	\textbf{地球——人类已知的唯一一颗具有液态水的行星}
	
	地球的温度;
	地球与太阳的距离;
	水是外动力地质作用的重要介质;
	地表水圈与地球内部水之间的循环;
	地球和月球的含水性存在本质差异.
	
	\textbf{地球——人类已知的唯一一颗孕育生命和文明的行星}
	
	约7亿年前新元古代雪球地球事件，引发寒武纪生命大爆发.
	\subsubsection{为什么地球如此独特?}
	地球之所以独特，是因为它是一个处于“动态平衡”的湿润星球。板块构造不仅是地质现象，更是地球生命支撑系统（Life Support System）的基础设施。
	
	\textbf{板块构造(plate tectonics)}
	
	其他行星（如金星或火星）通常被认为是“停滞盖层”（Stagnant Lid）构造，即整个星球表面是一块完整的硬壳。而地球的岩石圈被分割成数个动力学板块。
	
	水分的关键作用：地球拥有液态海洋。水进入地幔后，能够显著降低岩石的剪切强度和粘度，像“润滑剂”一样让坚硬的岩石圈变得可以滑动。没有水，地球可能也会变成像金星那样死寂的硬壳。
	
	热对流的平衡：板块构造是地球散热最高效的方式，它像一套精密的“冷却系统”，维持了地球数十亿年的宜居温度。
	
	\textbf{板块俯冲带(subduction zones)}
	
	俯冲带是板块构造的核心，也是地球上最复杂的地质“实验室”。
	回收与再生：在俯冲带，密度较大的洋壳（玄武岩）带着冰冷的表层物质和水钻入地幔。这个过程实现了地表物质与深部物质的循环。
	
	通量熔融 (Flux Melting)：这是俯冲带最神奇的地方。被俯冲带入深处的洋壳会“脱水”，释放出的流体进入上方温热的地幔楔。水降低了地幔岩石的熔点，导致其发生部分熔融，产生了大量的岩浆。
	
	地震与火山的摇篮：如你之前看到的图示，俯冲带产生了全球最深的地震（板内地震）和最具爆发力的火山弧（如环太平洋火山带）。
	\subsubsection{大陆、花岗岩的形成}
	地质学界有一句名言：“地球是产生花岗岩的机器。” 在月球、火星上，主要的岩石是玄武岩，而大规模的花岗岩几乎只存在于地球。
	
	从玄武岩到花岗岩：
	
	1)最初的地球表面主要是玄武岩（洋壳）。
	
	2)通过俯冲带的多次重熔和分异：岩浆在上升过程中，较重的成分沉积，较轻的富含硅、铝的成分上升。
	
	3)最终形成了密度轻、化学性质稳定的岩石——花岗岩。
	
	大陆的产生：由于花岗岩密度比地幔和洋壳都轻，它们就像“软木塞”一样漂浮在上方，不会像洋壳那样轻易被俯冲回地幔。随着时间推移，这些花岗岩不断汇聚、生长，形成了我们今天居住的大陆地壳。
	
	\subsection{太阳系的形成与演化}
	\subsubsection{宇宙初始：大爆炸 (The Big Bang)}
	时间点：约138亿年前。
	
	核心过程：宇宙从一个极高温度、极高密度的“奇点”开始急剧膨胀。在最初的几秒钟内，宇宙经历了暴胀、夸克-胶子等离子体阶段。
	
	物理意义：大爆炸奠定了宇宙的时空框架以及物质基础。
	\subsubsection{核素形成：H, He 原子核 (大爆炸核合成)}
	
	过程：在大爆炸后的3分钟内，温度降至约10亿度，质子和中子开始结合。
	
	产物：形成了宇宙中最基础的化学组成,氢 (H) 和 氦 (He)(以及极少量的锂 Li)。
	
	现状：直到今天，这两种元素仍占据宇宙可见物质质量的约98\%以上。
	\subsubsection{恒星的形成与演化 (The Life of Stars)}
	
	第一代恒星：大爆炸后约数亿年，在引力作用下，原始氢氦云气坍缩形成了第一代超大质量恒星（第三星族星）。
	
	演化循环：恒星如同宇宙的“炼金炉”。小质量恒星通过氢核聚变发光；大质量恒星在其生命末期会发生剧烈的超新星爆发，将合成的重元素抛射回星际空间，成为下一代恒星和行星的原料。
	
	\subsubsection{核素形成：恒星中元素合成过程 (Stellar Nucleosynthesis)}
	
	这是地球拥有复杂化学组成的根本原因：
	
	主序星阶段：氢聚变为氦。
	
	后期阶段：氦聚变为碳、氧、氖、镁，直至形成铁 (Fe)。
	
	重元素合成：比铁重的元素（如金、铀等）主要通过超新星爆发（r-过程）或中子星碰撞产生。
	
	结论：你身体里的每一个氧原子，你血液里的铁原子，都曾是某颗恒星核心的一部分。“我们皆为星尘。”
	
	\subsubsection{原始太阳的形成 (Solar Nebula)}
	
	星云坍缩：约46亿年前，银河系旋臂中的一团分子云（太阳星云）在引力波动下（可能受附近超新星爆发的扰动）开始坍缩。
	
	角动量守恒：随着星云收缩，它旋转得越来越快，中心物质堆积形成原始太阳，周围物质受离心力影响扁平化，形成一个旋转的盘状结构。
	
	\subsubsection{原行星盘和原始行星的形成 (Protoplanetary Disk)}
		原行星盘 (Protoplanetary Disk)：这是行星诞生的温床。
		
		温度梯度与雪线：靠近中心的地方温度极高，只有金属和硅酸盐能凝结（形成类地行星：水金火地）；远离中心的地方冷到足以让冰和气体凝结（形成类木行星：木土天海）。
		
		吸积作用 (Accretion)：
		
		微米级的尘埃靠静电力聚集成毫米级颗粒。
		
		颗粒碰撞合并形成微行星 (Planetesimals)（公里级）。
		
		微行星通过引力相互掠夺，形成行星胚胎。
		
		原始地球的诞生：在太阳系内部，无数次的撞击和汇聚最终形成了原始地球。

	\begin{table}[htbp]
		\centering
		\caption{太阳系行星地球物理与地学特征对比汇总}
		\label{tab:geophysics_planets}
		\footnotesize
		\begin{tabularx}{\linewidth}{l X X X X}
			\toprule
			行星 & 内部结构 (壳/幔/核) & 磁场特性 (发电机效应) & 挥发分与水分 (H$_2$O) & 大气特征与地表动力学 \\
			\midrule
			\textbf{水星} & 极薄硅酸盐壳幔；巨型液态铁核 (占直径$>$75\%) & 弱全球性偶极磁场 (地球的1\%)；内核热对流驱动 & 表面极度贫水；极地永久阴影区存在沉积水冰 & 外逸层 (Exosphere)；无板块构造；表面受冷缩形成冲断褶皱 \\
			\addlinespace
			\textbf{金星} & 停滞盖层 (Stagnant Lid)；无分层岩石圈；幔核比例类地 & 无全球性内部磁场；仅由太阳风诱导的弱磁场 & 极度干燥；水分通过氢逃逸流失；地幔无流体弱化 & CO$_2$超稠密大气；失控温室效应；全球性火山重塑事件 \\
			\addlinespace
			\textbf{地球} & 活跃板块构造；双模地壳 (花岗岩/玄武岩)；D''层复杂幔底 & 强偶极磁场；液态外核热-成分对流 (Geodynamo) & 拥有液态海洋；地幔含有显著比例的结构水 & N$_2$-O$_2$大气；活跃碳循环；唯一拥有大陆地壳的行星 \\
			\addlinespace
			\textbf{火星} & 壳层二分性 (南高北低)；厚岩石圈；部分熔融核 & 无活跃全球磁场；仅存强烈的地壳剩余磁化 (化石磁场) & 存在极冠水冰及深层地下冰/卤水 & 稀薄CO$_2$大气；大规模死火山 (奥林帕斯山)；地质历史早期曾有液态水 \\
			\addlinespace
			\textbf{木星} & 无固态表面；金属氢地幔；超高压固态/稀释核 & 极强磁场 (表面为地球14倍)；金属氢对流驱动 & 大气层深部含水；卫星 (如木卫二) 拥有地下海洋 & 氢氦大气；强射电辐射；大红斑长寿气旋 \\
			\addlinespace
			\textbf{土星} & 类似木星；金属氢层相对较薄；可能存在巨大岩石核 & 强磁场；轴对称性极高；氦雨沉降释放引力能 & 大气含水量低于木星；土卫六拥有甲烷/乙烷湖泊 & 氢氦大气；平均密度小于水；显著的环系统动力学 \\
			\addlinespace
			\textbf{天王星} & 氢氦包层；中间层为超离子态水-氨-甲烷冰幔 & 非典型复杂磁场；偏心且倾斜；可能起源于浅层离子对流 & “冰巨星”；内部富含含水化合物的超离子态 & 极冷大气 (约49K)；自转轴极端倾斜 ($98^\circ$)；侧卧运行 \\
			\addlinespace
			\textbf{海王星} & 结构类天王星；内部热流高于天王星 & 偏心偶极磁场；与自转轴夹角达$47^\circ$ & 富含“冰”组分 (挥发分)；内部可能存在氨水海洋 & 强烈大气活动；风速达2100 km/h；大气层富含甲烷呈深蓝色 \\
			\bottomrule
		\end{tabularx}
	\end{table}
	-
	
	-
	\subsection{火星的水与生命现象}
		
	\textbf{蛇纹石化热液体系（Serpentinization）}
	
	蛇纹石化不仅是水的证据，更是能量的证据。
	
	反应过程：当富含橄榄石的地幔岩石与水接触时，会发生如下反应:
	
	橄榄石+水→蛇纹石+磁铁矿+氢气$(H_2)$
	
	生命的摇篮：生成的氢气可以与$CO_2$反应生成甲烷。在地球上（如深海热液喷口“失落之城”），这种不需要阳光的化学反应支撑了庞大的化能合成生态系统。
	
	火星意义：火星早期拥有厚厚的地壳和热液活动，这种体系可能在地下深处维持了数亿年的宜居环境，即便表面已经干涸。
	\begin{table}[htbp]
		\centering
		\caption{火星关键矿物发现及其生物地学意义}
		\begin{tabular}{lll}
			\toprule
			关键矿物/现象 & 化学组成 & 地质/生物学解释 (地球物理视角) \\
			\midrule
			\textbf{蛇纹石} & $(Mg,Fe)_3Si_2O_5(OH)_4$ & 暗示水岩相互作用；释放 $H_2$，支撑化能合成生命。 \\
			\addlinespace
			\textbf{蓝铁矿} & $Fe_3(PO_4)_2 \cdot 8H_2O$ & 典型的缺氧、有机质富集环境指示矿物；磷循环参与。 \\
			\addlinespace
			\textbf{胶硫铁矿} & $Fe_3S_4$ & 磁性矿物；通常与硫酸盐还原细菌的代谢活动相关。 \\
			\addlinespace
			\textbf{“豹纹”斑点} & 铁-镁氧化还原环 & 指示化学能量梯度；可能为有机物降解导致的局部还原。 \\
			\addlinespace
			\textbf{有机大分子} & 芳香族/脂肪族化合物 & 生命的构建基块；但需排除非生物合成 (如费托合成)。 \\
			\bottomrule
		\end{tabular}
	\end{table}
	
	\subsection{陨石}
	\subsubsection{陨石的三大类型}
	\textbf{1. 石陨石 (Stony Meteorites) —— 行星的“皮肤”与“初态”}
	
	这是最常见的类型，占坠落总数的94\%左右。
	
	球粒陨石 (Chondrites)：最原始的陨石。它们从未经历过母体的熔融分异，完整保留了45.6亿年前太阳星云的化学成分。
	
	关键特征：含有“球粒”（毫米级的硅酸盐小球）和CAI（钙铝包裹体），后者是太阳系中已知最古老的固体物质。
	
	碳质球粒陨石：富含水和有机物（氨基酸等），是研究生命起源的宝库。
	
	无球粒陨石 (Achondrites)：源自已经分异的行星或大型小行星的地壳或地幔。它们经历了熔融过程，球粒结构消失。包括来自月球和火星的陨石。
	
	\textbf{2. 铁陨石 (Iron Meteorites) —— 行星的“心脏”}
	
	主要由铁镍合金组成。
	
	来源：源自大型小行星在演化初期通过引力分异形成的金属核心。当这些小行星在剧烈撞击中破碎时，核心碎片进入太空。
	
	科学价值：由于人类无法到达地球内核，铁陨石是研究类地行星核心成分（铁、镍及轻元素）和磁场起源的最佳类比物。
	
	\textbf{3. 石铁陨石 (Stony-Iron Meteorites) —— “过渡带”的标本}
	
	这是最稀有的类型，约占1\%。
	
	橄榄陨铁 (Pallasites)：金黄色的橄榄石晶体嵌入金属基质中。一般认为它们来自小行星的核幔边界 (CMB)，具有极高的科研和审美价值。
	
	中铁陨石 (Mesosiderites)：由破碎的硅酸盐和金属混合而成，通常是天体剧烈碰撞混合后的产物。
	
	\subsubsection{陨石能告诉我们什么信息？}
	\begin{table}[htbp]
		\centering
		\caption{陨石分类及其地球物理/行星科学意义汇总}
		\small
		\begin{tabularx}{\textwidth}{l l l X}
			\toprule
			大类 & 子类 & 母体来源 & 科学意义 (地球物理/地球化学视角) \\
			\midrule
			\textbf{石陨石} & 球粒陨石 & 原始星云、未分异小行星 & 太阳系最原始组分；确定太阳系年龄；提供地球原始丰度参照。 \\
			& 无球粒陨石 & 行星/大型小行星地壳或地幔 & 研究类地行星地质演化；识别火星、月球等地外天体样本。 \\
			\addlinespace
			\textbf{铁陨石} & 八面体铁陨石等 & 分异天体的金属核心 & 研究类地行星核部成分（Fe-Ni合金）；了解早期发电机效应。 \\
			\addlinespace
			\textbf{石铁陨石} & 橄榄陨铁 & 分异天体的核幔边界 (CMB) & 揭示天体内部核幔分异过程及复杂的边界动力学性质。 \\
			& 中铁陨石 & 灾难性撞击混合物 & 记录早期太阳系的天体大碰撞事件及其动力学过程。 \\
			\midrule
			\textbf{特殊类型} & 碳质球粒陨石 & 原始富挥发分小行星 & 含有机大分子与含水矿物；探索地球水起源与前生物化学。 \\
			\bottomrule
		\end{tabularx}
	\end{table}
	
	通过对球粒陨石中CAI的铅-铅同位素定年，科学家确定了太阳系的年龄约为 45.67 亿年。
	
	地球的原始配方：球粒陨石（尤其是顽火辉石球粒陨石）的同位素比例与地球高度相似。地球物理学家通过陨石成分来推测“全地球”的化学组成（Bulk Earth Composition）。
	
	水的来源与生命的种子：碳质球粒陨石含有大量的水（以结合水形式存在）和复杂的有机分子（如核碱基、氨基酸）。这支持了“地球上的水和生命原料可能由陨石撞击带回”的假说。
	
	恒星爆发的前世：陨石中含有“前太阳颗粒”（Pre-solar grains），它们在太阳系形成之前就已存在。通过研究它们，我们可以了解上一代恒星死亡（超新星爆发）的过程。
	
	
	陨石是地质学研究的基准线。没有陨石，我们无法确定地球的年龄，也无法得知地球深处那个不可触及的铁核究竟由什么组成。它们是将地球放入宇宙大背景下审视的关键物证。
	
	\subsection{月球是如何形成的？}
	\textbf{1. 核心理论：大碰撞假说 (The Giant Impact)}
	
	该理论认为，约45亿年前，一颗名为“忒伊亚”（Theia）、大小约如火星的行星胚胎，以一定的角度撞击了原始地球。
	物理过程：撞击释放了巨大的能量，使忒伊亚和原始地球的地幔部分完全熔融并汽化。
	
	物质分配：忒伊亚的金属核沉入地球中心，而撞击抛射出的硅酸盐蒸汽云（Synestia）在地球轨道上冷凝、吸积，最终形成了月球。
	
	\textbf{2. 地球物理与地球化学证据}
	
	\textbf{A. 密度与核心比例 (Density \& Core Size)}
	
	证据：月球的平均密度约为 $3.34g/cm^3$，远低于地球的 5.51g/cm3。月球金属核仅占总质量的 $1\%∼2\%$，而地球核心占约 $30\%$。
	
	解释：这说明月球主要由“贫铁”的硅酸盐物质构成，符合大碰撞中“取自地幔、舍弃金属核”的物理模型。
	
	\textbf{B. 同位素一致性 (Isotopic Consistency) }
	
	—— 最强证据，也是最大谜团
	
	证据：月球岩石与地球地幔在氧 ($^{18}O/^{16}O$)、钛、铬、钨等同位素组成上几乎完全一致。
	解释：在太阳系中，不同位置形成的天体具有不同的同位素“指纹”。这种高度一致性暗示月球和地球拥有完全相同的“原材料”，或者撞击后的物质交换达到了极致的平衡。
	
	\textbf{C. 挥发分亏损 (Depletion of Volatiles)}
	
	证据：相比于地球，月球极度缺乏易挥发的元素（如钾 K、钠 Na、锌 Zn、氯 Cl）。
	
	解释：这表明月球在形成过程中经历过极端的高温过程，导致挥发性成分在吸积成核前就已汽化逃逸。
	
	\textbf{D. 月球岩浆洋 (Lunar Magma Ocean, LMO)}
	
	证据：月球高地主要由斜长岩（Anorthosite）组成。
	
	解释：这证明月球形成初期是完全熔融的。随着冷却，较轻的斜长石向上浮动形成原始月壳，较重的橄榄石和辉石下沉形成月幔。这种全球性的熔融状态是大碰撞产生的巨大热能的直接后果。
	
	\textbf{E. 角动量 (Angular Momentum)}
	
	证据：地月系统的总角动量异常高，且地球自转轴有 $23.5∘$
	的倾角。
	
	解释：大碰撞模型在动力学模拟上能完美解释这种倾斜和自转特征。
	
	\subsection{宜居的行星的关键要素}
	\begin{table}[htbp]
		\centering
		\caption{宜居行星形成的关键地球物理与天体生物学要素}
		\small
		\begin{tabularx}{\textwidth}{l l X}
			\toprule
			关键要素 & 物理/化学机制 & 对生命的意义 \\
			\midrule
			\textbf{适居带轨道} & 热辐射平衡 & 确保行星表面能够长期存在液态水。 \\
			\addlinespace
			\textbf{行星质量} & 引力束缚 & 维持足够的大气压，防止挥发分逃逸到太空。 \\
			\addlinespace
			\textbf{内部发电机} & 偶极磁场 & 阻挡高能宇宙射线及太阳风剥蚀大气层。 \\
			\addlinespace
			\textbf{板块构造} & 深部物质循环 & 实现长期碳循环，起到调节全球气候的作用。 \\
			\addlinespace
			\textbf{水岩作用} & 氧化还原梯度 & 提供能量来源（如氢气），支撑化能合成生物。 \\
			\addlinespace
			\textbf{大行星守护} & 引力屏障 & 如木星可散射外围彗星，降低内行星受撞击频率。 \\
			\bottomrule
		\end{tabularx}
	\end{table}
	\textbf{1. 恒星与天文位置：金发姑娘原则（The Goldilocks Zone）}
	
	行星必须位于一个恒星系统的“适居带”内，即距离恒星远近适中。
	
	液态水：距离太近，水分会因失控温室效应而逃逸（如金星）；距离太远，水分会永久冻结（如火星）。
	
	恒星稳定性：母恒星（如太阳）必须有足够长的寿命（主序星阶段），给生命演化留出数亿至数十亿年的时间。高质量恒星寿命太短，低质量红矮星则容易产生强烈的耀斑辐射。
	
	\textbf{2. 质量与引力：大气层的“守门员”}
	
	行星的质量决定了它能否“留住”生命所需的物质。
	
	大气保持能力：质量太小（如火星或月球），引力不足以束缚浓厚的大气层，导致气压过低且水分流失。
	
	大气组分：需要足够的气压来维持水的液态形式，并提供温室效应（防止全球冻结）以及抵御紫外线的保护层（臭氧层）。
	
	\textbf{3. 内部动力学：全球磁场（The Shield）}
	
	这是地球物理学中最为关键的要素之一。
	
	液态外核与发电机效应：行星必须有一个旋转的液态金属核心来产生全球性磁场。
	
	磁屏蔽作用：磁场像一面盾牌，阻挡高能太阳风剥蚀大气层。火星失去宜居性的核心原因之一就是其内部发电机停止，磁场消失，导致大气被太阳风“吹走”。
	
	\textbf{4. 地质活跃度：板块构造（The Thermostat）}
	
	地球是目前已知唯一拥有活跃板块构造的行星，这不仅是地质现象，更是行星空调。
	
	碳-硅酸盐循环：通过俯冲带将 $CO_2$
	
	带入深部，再通过火山活动释放。这一机制在数亿年的尺度上调节大气温室气体浓度，使气候不至于极热或极冷。
	
	营养物质循环：板块构造不断将地幔深处的矿物质（如磷、钾）带到表层，为生物圈提供养分。
	
	\textbf{5. 关键化学组分：生命的原料（CHNOPS）}
	
	挥发分与有机物：必须富含碳、氢、氮、氧、磷、硫。
	
	水岩相互作用：如之前提到的蛇纹石化，提供化学能。
	
	适度的还原性环境：允许从无机分子向有机单体转换。
	
	\textbf{总结：地球的“稀有”之处}
	
	在宇宙中，满足其中一两个条件的行星可能很多，但同时满足所有条件的行星却极少。例如：
	
	金星：有质量和活跃火山，但离太阳太近，没有板块构造，水分流失。
	
	火星：在适居带边缘，但质量太小，磁场过早消失，变成了寒冷荒漠。
	
	木卫二 (Europa)：有大量水和能量，但没有合适的大气和光照环境（生命可能仅限深海）。
	
	地球的独特性在于它维持了一个完美的动态平衡，将内部动力学（磁场、板块）与外部动力学（太阳能量、轨道）精准地耦合在一起。
	
	\section{矿物与岩石}
	
	\subsection{地质学和地球化学}
	\subsubsection{概念}
	地质学: 研究地球及其它类地行星的自然科学，是地球科学的重要组成部分.
	
	主要内容：地球内部圈层的物质结构、物质组成、地质过程、时空演化，及其与表生大气圈、水圈、生物圈的相互作用.
	
	矿物学、岩石学、矿床学、地球化学、地层与古生物学等.
	
	地球化学: 从化学的角度，定量化探讨地球和行星的起源演化，重大地质事件和地质过程.
	
	\subsubsection{行星地质学研究方法}
	野外观测:地质现象的认识、判别，收集关键地质样品.
	
	实验分析:整理观察结果，开展样品的物理、化学研究分析；高压实验模拟和数值模拟.
	
	将今论古:获得数据资料的归纳总结对比，探讨成因和演化过程.
	
	\subsection{什么叫矿物？}
	1)地质作用或宇宙作用形成的（强调非人工合成）.
	
	2)具有一定化学成分与内部晶体结构，且在一定物理化学条件下相对稳定的化合物或单质 (由此可区分不同矿物种).
	
	3)岩石和矿石的基本组成单位（最基本的物相）
	
	例如：食盐（$NaCl$），石英（$SiO_2$）， 黄铁矿（$FeS$），钾长石（$K[AlSi_3O_8]$）
	
	此外还有：
	
	1)人工矿物（合成矿物）
	
	2)准矿物（强调非晶态）：准矿物有自发转变为矿物的趋势
	\subsubsection{矿物的研究方法}
	\begin{table}[htbp]
		\centering
		\caption{矿物主要研究方法分类汇总}
		\label{tab:mineral_methods}
		\small
		\begin{tabularx}{\textwidth}{llX}
			\toprule
			分析类型 & 研究手段 & 常用仪器/方法 \\
			\midrule
			\textbf{形貌分析} & 显微形貌观测 & 光学显微镜 (OM)、扫描电镜 (SEM) \\
			\addlinespace
			\textbf{结构分析} & 晶体结构与物相 & X射线衍射仪 (XRD)、透射电子显微镜 (TEM)、激光拉曼光谱仪 (Raman)、红外光谱仪 (IR) \\
			\addlinespace
			\textbf{成分分析} & 化学元素与同位素 & 湿化学法、X荧光光谱分析 (XRF)、电子探针 (EPMA)、ICP-MS、离子探针 (SIMS) \\
			\midrule
			\multicolumn{3}{l}{\textbf{常用鉴定组合：} X射线衍射仪 (XRD)、电子探针 (EPMA)、激光拉曼光谱 (Raman) 等} \\
			\bottomrule
		\end{tabularx}
	\end{table}
	\subsubsection{矿物的晶体化学分类体系}
	\begin{table}[htbp]
		\centering
		\caption{矿物的晶体化学分类体系（以钛辉石为例）}
		\label{tab:mineral_classification}
		\small
		\begin{tabularx}{\textwidth}{l p{7cm} l}
			\toprule
			\textbf{级 序} & \textbf{划 分 依 据} & \textbf{举 例（钛辉石）} \\
			\midrule
			大 类 & 化合物类型 & 含氧盐大类 \\
			类 & 阴离子或络阴离子种类 & 硅酸盐类 \\
			(亚类) & 络阴离子结构类型 & 链状硅酸盐亚类 \\
			族 & 晶体结构型和阳离子性质 & 辉石族 \\
			(亚族) & 阳离子种类 & 单斜辉石亚族 \\
			种 & 一定的晶体结构和化学组成 & 普通辉石 \\
			(亚种) & 在完全类质同像中，根据其所含端元组分的比例划分 & -- \\
			(变种或异种) & 化学组成或物性、形态等方面变异 & 钛质普通辉石 (钛辉石) \\
			\bottomrule
		\end{tabularx}
	\end{table}
	
	\subsubsection{矿物的晶体化学分类}
	1)种：具有确定的晶体结构和相对稳定的化学成分。
	
	2)同质多像变体：成分相同，结构不同，物理性质也不同，属不同矿物种。如：金刚石与石墨。
	
	3)多型变体：成分相同，结构差异不大，物理性质也相近，属同一矿物种。如：$2H$石墨与$3R$石墨。
	
	--同质多像化合物为不同矿物种(结构不同)
	
	例：$C, SiO_2,TiO_2$
	
	--多型化合物属于同一矿物种
	
	例：石墨的$2H,3R$多型
	\subsubsection{矿物的命名}
	1)化学成分：钨锰铁矿(Mn,Fe)WO4，自然金
	
	2)物理性质：方解石，孔雀石，天青石，蛇纹石
	
	3)物理性质与化学成分：黄铜矿(CuFeS2)，闪锌矿(ZnS)，磁铁矿(Fe3O4)，方铅矿(PbS)， 辉锑矿(Sb2S3)
	
	4)形态：石榴子石，十字石，菊花石
	
	5)以人名或地名命名：高岭石，包头矿，章氏硼镁石，张衡矿，袁复礼石
	
	6)沿用古名：石英，石膏，辰砂，雄黄
	\subsubsection{矿物的产出}
	1)岩石类型：超基性岩、基性岩、中性岩、酸性岩.
	
	2)温度/压强：不同深度、构造环境.
	
	3)水：改造 vs. 沉淀 （海水、卤水、大气降水）.
	
	\subsection{金刚石}
	\subsubsection{分类}
	依据：氮 (N) 的含量及存在状态。
	
	\begin{forest}
		for tree={
			draw, rounded corners, font=\small\sffamily,
			fill=white!20, % 还原原图黄色框
			edge={thick, -latex},
			l sep=1.2cm, s sep=0.5cm,
			align=center,
		}
		[金刚石
		[I型 (含N)
		[Ia型 (含N较多)
		[IaA型]
		[IaB型]
		]
		[Ib型 (含N较少)]
		]
		[混合型]
		[II型 (缺N)
		[IIa型 (不含B)]
		[IIb型 (含B)]
		]
		]
	\end{forest}
	
	\textbf{I型 (Type I)：主要类型。}
	
	1)Ia型：氮成群聚集。这是绝大多数天然钻石（98\%）的类型。
	
	2)Ib型：氮原子分散存在。天然极少，常见于人造钻石。
	
	\textbf{II型 (Type II)：极度纯净或含硼。}
	
	1)IIa型：几乎不含杂质，热导率极高。
	
	2)IIb型：含硼 (B)，产生半导体性质，外表常呈蓝色。
	
	科学价值：II型金刚石虽然仅占2\%左右，但许多世界名钻（如库里南）均属于此列。
	
	\subsubsection{金刚石的产出特征}
	岩石:橄榄岩、金伯利岩、榴辉岩、沉积岩.其中金刚石的主要赋存岩石是金伯利岩
	
	深度:120-2000多公里, 上地幔、地幔过渡带到下地幔
	
	\subsubsection{砂矿型金刚石}
	形成于内生金刚石矿区附近：风化、剥蚀、搬运、沉积. 
	
	金刚石非常稳定，抗风化.
	
	\subsubsection{金刚石与地学前沿问题}
	地质环境中：石墨和金刚石的温压分布及转化.
	
	石墨→金刚石：石墨经高温高压可转变为金刚石，C原子从石墨的3配位变成金刚石中4配位。
	
	金刚石→石墨: 低温、低压的条件下转变，速率很慢。
	\subsubsection{金刚石示踪地球深部的水含量}
	金刚石来源：巴西Juína
	
	类型：砂矿型
	
	主要特征：包含林伍德石包裹体
	
	红外光谱测定林伍德石含1.4wt\%水，指示地球过渡带某些区域可能富含大量的水。
	
	\subsection{超高压变质矿物}
	\subsubsection{大陆地壳深俯冲的可能性：密度与力的平衡}
	
	虽然大陆地壳的密度（2.6--2.9 g/cm$^3$）显著低于上地幔（3.2--3.6 g/cm$^3$），但在特定的板块动力学环境下，陆壳能够进入地幔深度。其主要地球物理机制包括：
	
	\begin{itemize}[leftmargin=1.5em]
		\item \textbf{板片拉力 (Slab Pull) 的主导作用：} 大陆地壳通常并非独立俯冲，而是紧跟在大洋岩石圈之后。当高密度的洋壳进入深部地幔时，产生的巨大牵引力（Slab Pull）会将连接的陆壳强制拖入俯冲带。
		\item \textbf{高压相变导致的密度增益：} 随着压力增加，大陆地壳物质会发生变质相变。例如，长石转化为柯石英，基性岩转化为\textbf{榴辉岩}（其密度可达 3.4--3.5 g/cm$^3$），这极大地减小了陆壳与地幔之间的浮力差。
		\item \textbf{板块会聚的挤压推力：} 强烈的板块边界聚合运动提供了必要的机械能，辅助克服浮力阻力。
	\end{itemize}
	
	\subsubsection{动力学过程：从深埋到折返的机制}
	
	大陆地壳的“地质过山车”演化过程可分为三个核心动力学阶段：
	
	\textbf{深俯冲阶段 (Deep Subduction)}
	
	在洋壳板片的牵引下，大陆边缘物质被带入离地表 100--150 公里的地幔深度，形成超高压（UHP）变质岩。此时，陆壳与下伏的岩石圈地幔保持物理耦合。
	
	\textbf{板片断离触发 (Slab Break-off)}
	
	当俯冲阻力达到极限时，深部较重的大洋板片会与较轻的大陆板片发生机械断裂（板片断离）。这一事件会导致向下牵引的拉力瞬间消失，是折返发生的关键转折点。
	
	\textbf{浮力驱动的快速折返 (Exhumation)}
	
	失去拉力后，大陆地壳物质由于其固有的低密度，在巨大的正向浮力作用下开始向上移动。主要的折返模式包括：
	\begin{enumerate}
		\item \textbf{通道流模式 (Channel Flow)：} 物质沿低粘度的俯冲通道像流体一样挤回地表。
		\item \textbf{楔形挤出 (Wedge Extrusion)：} 陆壳物质以刚性楔形体的形式在挤压应力场中向上运移。
	\end{enumerate}
	
	\subsubsection{动力学参数汇总}
	
	\begin{table}[htbp]
		\centering
		\caption{大陆地壳俯冲与折返的力学要素总结}
		\label{tab:dynamics}
		\begin{tabularx}{\textwidth}{l l X}
			\toprule
			阶段 & 主导力 & 关键物理现象 \\
			\midrule
			\textbf{俯冲下行} & 板片拉力 + 构造挤压 & 克服浮力阻力，进入 $\sim$100 km 深度 \\
			\textbf{深部停留} & 静水压力 + 相变 & 形成超高压矿物（金刚石、柯石英） \\
			\textbf{折返触发} & 板片断离 & 牵引力消失，势能向浮力能转化 \\
			\textbf{快速上升} & 正向浮力 + 剪切流 & 穿过俯冲通道，最终抬升至地表 \\
			\bottomrule
		\end{tabularx}
	\end{table}
	\subsubsection{矿物相变与压强指示矿物}
	\begin{table}[htbp]
		\centering
		\caption{重要矿物相变点与压强指示意义汇总}
		\label{tab:phase_transitions}
		\small
		\begin{tabularx}{\textwidth}{l l l X}
			\toprule
			系列 & 矿物名称 & 稳定压强/深度区域 & 地质/地球物理意义 \\
			\midrule
			\textbf{碳 ($C$)} & 金刚石 (Diamond) & $> 4.5\text{--}6.0\text{ GPa}$ & 深部碳循环证据；来自 $>150\text{km}$ 深度。 \\
			\addlinespace
			\textbf{二氧化硅} & 柯石英 (Coesite) & $\sim 2.5\text{--}3.0\text{ GPa}$ & 超高压 (UHP) 变质作用标志；陆壳深俯冲证据。 \\
			\textbf{二氧化硅} & 斯石英 (Stishovite) & $> 9.0\text{ GPa}$ & 硅由 4 配位向 6 配位转变；指示极高压环境。 \\
			\addlinespace
			\textbf{橄榄石} & 瓦兹利石 ($\beta$-phase) & $\sim 410\text{ km}$ 界面 & 对应 410km 地震波速度突变；重要的储水相。 \\
			\textbf{橄榄石} & 铃木石 ($\gamma$-phase) & $520\text{--}660\text{ km}$ & 过渡带下部主要矿物。 \\
			& 桥石 (Bridgmanite) & $> 660\text{ km}$ 界面 & 下地幔主成分；标志着硅酸盐向钙钛矿结构转变。 \\
			\addlinespace
			\textbf{铝矿物} & 镁铝榴石 (Pyrope) & 上地幔深部 & 稳定于高压环境；常用于地质压强计计算。 \\
			\bottomrule
		\end{tabularx}
	\end{table}
	
	\textbf{核心要点解析}

		1)\textbf{克拉珀龙斜率 (Clapeyron Slope)：} 相变界面的深度会受温度影响。例如，在冷的俯冲带，410km 界面（正斜率）会抬升，而 660km 界面（负斜率）会下沉。
		
		2)\textbf{配位数演化：} 随着压强增加，矿物结构变得更加紧密，阳离子的配位数普遍增加（如 $Al$ 和 $Si$ 从 4 变为 6）。
		
		3)\textbf{密度突变：} 矿物相变伴随着 5\%--10\% 的密度增加，这是维持地幔对流模式的关键物理因素。
		
	\textbf{专业术语说明}
	
	1)UHP (Ultra-High Pressure)：超高压变质作用，特指柯石英或金刚石稳定的压强区间。
	
	2)Bridgmanite (桥石)：原称硅酸盐钙钛矿（Silicate Perovskite），是地球深部最重要的压强指示矿物。
	
	3)Clapeyron Slope (克拉珀龙斜率)：描述压力-温度变化对相变界面的影响，是地球物理学家分析地幔温度异常的重要工具。
	
	\subsection{锆石$ZrSiO_4$}
	
	\subsubsection{早期地球的见证者：40亿年的时间标尺}
	锆石是地球上最古老的矿物。通过对锆石进行 U-Pb 同位素测年，科学家在澳大利亚西部的 \textbf{Jack Hills} 发现了年龄高达 44 亿年（4.4 Ga）的颗粒。
	
	\begin{itemize}
		\item \textbf{最古老岩石记录：} 锆石不仅存在于沉积岩中，还揭示了地球上最古老的岩石单元（如加拿大 Acasta 片麻岩，约 4.0 Ga）。
		\item \textbf{地壳演化指示：} 锆石中的氧同位素（$\delta^{18}O$）特征暗示了早在 42 亿年前，地球表面可能已经存在液态水和初步的陆壳分异。
	\end{itemize}
	
	\subsubsection{地磁场起源：锆石中的磁性包裹体}
	地磁场是地球宜居性的屏障。探测冥古宙（Hadean）是否存在地磁场，是理解大气保持和生命起源的核心。
	
	\textbf{Jack Hills 锆石的古地磁研究}
	
	研究者利用超高灵敏度的超导量子干涉仪（SQUID）显微镜，对含有微量铁矿物包裹体（如磁铁矿 $\text{Fe}_3\text{O}_4$）的 Jack Hills 锆石进行了测量。
	
	\begin{table}[htbp]
		\centering
		\caption{Jack Hills 锆石年龄与古地磁强度研究概况}
		\label{tab:zircon_mag}
		\begin{tabularx}{\textwidth}{l l X}
			\toprule
			研究指标 & 典型数据范围 & 地球物理意义 \\
			\midrule
			\textbf{锆石结晶年龄} & 3.0 -- 4.2 Ga & 提供了地球早期热演化的绝对时间坐标。 \\
			\addlinespace
			\textbf{古磁场强度} & 1.0 -- 25.0 $\mu$T & 指示早期地发电机效应（Geodynamo）可能在 40 亿年前已开启。 \\
			\addlinespace
			\textbf{包裹体矿物} & 磁铁矿、赤铁矿 & 作为磁性载体，记录了结晶时的热剩磁（TRM）。 \\
			\bottomrule
		\end{tabularx}
	\end{table}
	
	\subsubsection{锆石中磁性矿物包裹体是原生的吗？}
	尽管测量结果显示早期地球拥有磁场，但国际学术界仍存在重大争论，核心在于磁性信号的\textbf{可靠性}。
	
	\begin{enumerate}
		\item \textbf{次生蚀变风险：} 质疑者认为，锆石虽然稳定，但其内部的裂隙可能允许后期流体渗透，导致次生磁性矿物的沉淀。
		\item \textbf{热重置问题：} 如果岩石在后期地质事件中被加热超过磁铁矿的居里点（580$^\circ$C），原始磁信号将被抹除并记录后期信号。
		\item \textbf{原生性判定：} 如何证明锆石内部的 Fe 矿物是在 40 亿年前结晶时被密封进去的（原生），而非通过交代作用进入的（次生），是目前古地磁学的核心挑战。
	\end{enumerate}
	
	\textbf{结论与展望}
	
	如果 Jack Hills 锆石中的磁信号被证实为原生，则意味着地球在 42 亿年前就拥有了足以抵御太阳风剥蚀的强大磁场。这不仅改写了地核演化历史，也为地球生命起源的时间线提供了物理支撑。
	
	\subsection{岩石}
	\subsubsection{概念}
	1)构成岩石圈(地壳和上地幔上部)、软流圈的主要物质, 是天然产出的、具有稳定外形的矿物集合体.
	
	2)主要由一种或数种造岩矿物按一定方式结合而成，少数由天然玻璃或生物遗骸组成.
	
	3)地球发展演化过程中，由各种地质作用形成的产物.
	
	\subsubsection{地球圈层结构的地震波制约}
	
	\textbf{地球内部主要不连续面}
	
	地震波在通过地球内部不同圈层时，由于介质的弹性参数和密度发生改变，其传播速度会产生显著的突变。
	
	\begin{itemize}[leftmargin=1.5em]
		\item \textbf{莫霍面 (Moho Discontinuity)：} 
		位于地壳与地幔的交界处。该界面最显著的特征是纵波速度（$v_p$）由约 $6.7\text{--}7.0 \text{ km/s}$ 突升至 $8.1 \text{ km/s}$ 左右。这一突变反映了物质从硅铝质/硅镁质岩石向超基性岩（如橄榄岩）的转变。
		
		\item \textbf{古登堡面 (Gutenberg Discontinuity / CMB)：} 
		位于地幔与地核的交界处（深度约 $2891 \text{ km}$）。在该界面，纵波速度（$v_p$）由约 $13.7 \text{ km/s}$ 急剧下降至约 $8.1 \text{ km/s}$；最为关键的是，\textbf{横波（$v_s$）完全消失}，这直接证明了外地核处于液态（流体）状态。
		
		\item \textbf{地幔过渡带 (Mantle Transition Zone, MTZ)：} 
		位于上地幔与下地幔之间（深度约 $410\text{--}660 \text{ km}$）。该区域存在两个显著的次级界面：
		\begin{itemize}
			\item \textbf{410 km 界面：} 橄榄石相变为瓦兹利石，导致波速台阶式上升。
			\item \textbf{660 km 界面：} 铃木石分解为桥石和铁方镁石，波速进一步突增。
		\end{itemize}
	\end{itemize}
	
	\textbf{地震波速度突变界面汇总表}
	
	\begin{table}[htbp]
		\centering
		\caption{地球物理主要界面及波速特征}
		\label{tab:discontinuities}
		\begin{tabularx}{\textwidth}{l l l X}
			\toprule
			界面名称 & 界面位置 & 深度 (km) & 地震波响应特征 \\
			\midrule
			\textbf{莫霍面} & 壳幔界面 & $\sim 5\text{--}70$ & $v_p$ 突升；反映波速从地壳向地幔的转变。 \\
			\addlinespace
			\textbf{410 km 界面} & 过渡带顶界 & 410 & $v_p, v_s$ 阶梯式增加；矿物相变引起。 \\
			\addlinespace
			\textbf{660 km 界面} & 过渡带底界 & 660 & 上、下地幔分界；波速显著跳跃。 \\
			\addlinespace
			\textbf{古登堡面} & 核幔界面 (CMB) & $\sim 2891$ & $v_s$ 消失 (入核)；$v_p$ 骤降；液态外核标志。 \\
			\addlinespace
			\textbf{莱曼面} & 内外核界面 & $\sim 5150$ & $v_p$ 再次增加；横波重新出现 (固态内核)。 \\
			\bottomrule
		\end{tabularx}
	\end{table}
	
	\textbf{专业术语解析}
	\begin{description}
		\item[P-wave (纵波):] 能够在固体和液体中传播，速度较快。
		\item[S-wave (横波):] 剪切波，只能在固体中传播，无法通过液态外核。
		\item[Phase Transition (相变):] 地幔过渡带波速跳跃的主要原因，即矿物晶体结构在高压下的重组。
	\end{description}
	
	\subsubsection{地核的组成}
	
	半径超过$3400 km$,以$Fe, Ni, S$等组成, 相当于铁陨石.液态外核旋转产生地磁场.
	
	由3部分组成:
	
	1)外核: 液态, $1742 km$.
	
	2)过渡层: 液态, 固体.
	
	3)内核: 固态.
	
	深内核（最新发现）：固态
	\subsubsection{地幔的组成}
	地幔厚度 $>2800 km$.
	
	上地幔: 以$Fe, Mg, Si, O$组成的超基性岩。软流圈地幔位于上地幔中部.
	
	地幔过渡带: 连接上下地幔的纽带.
	
	下地幔: 成分与上地幔一致，但密度更大 (相当于石铁陨石).
	
	岩石圈：地壳＋上地幔上部，位于软流圈之上，刚性板块.
	
	橄榄石瓦 兹利石 林伍德石 Fe-方镁石 + 钙钛矿.
	
	\subsubsection{地壳的组成}
	厚度7-50 km，平均~30 km.
	
	\textbf{大陆地壳}: 上地壳(硅铝质)和下地壳(硅镁质).
	
	\textbf{大洋地壳}: 硅镁质.
	
	1){硅铝质}: 以Si, O, Al等较轻元素组成, 比重小, 不连续.
	
		酸性岩浆岩+沉积岩+变质岩.
	
	2){硅镁质}: 以Si, O, Fe, Mg等较重元素组成, 比重大, 连续.
	
	地球形成初期发生熔融后, 结晶分异的产物, 最老可以追溯到4.4 Ga （Jack Hills, Australia）.
	
	\subsection{三大类岩石}
	三大岩石类型: 岩浆岩,变质岩,沉积岩.
	
	\subsubsection{岩石类型与地质作用}
	\textbf{岩石循环:}
	
	\includegraphics[width=0.5\textwidth]{3.png} 
	\includegraphics[width=0.5\textwidth]{tupian.png} 
	
	\textbf{岩石分布:}
	
	沉积岩：地表，占陆壳表面积的75\%；地壳越深部，分布越少.
	
	变质岩：地表和中下地壳，少量上地幔.
	
	岩浆岩：地壳、上地幔.
	
	地壳：岩浆岩66\%；变质岩20\%；沉积岩8\%.
	\subsubsection{地质作用}
	\textbf{地质作用}: 自然动力促使岩石圈的物质组成、结构以及地表形态发生变化和发展的作用.
	
	地球上的岩石是各种地质作用形成的产物，不同的地质作用形成不同类型的岩石.
	
	地质作用分为\underline{内动力地质作用}和\underline{外动力地质作用}.
	
	\textbf{内动力地质作用}
	
	1)因地球内部能量活动而引起的地质作用.
	
	2)内部能量：放射性热能、旋转能、重力能等.
	
	3)分类：构造运动（地震等）、岩浆作用、变质作用等.
	
	\textbf{外动力地质作用}
	
	1)因地球外部的能量活动引起的地质作用.
	
	2)外部能量：太阳能、引力能等.
	
	3)分类：风化、剥蚀、搬运、沉积等.
	
	\subsubsection{岩浆岩}
	
	由高温熔融产生的岩浆，在温压条件改变时，在地下凝固或喷发出地表后，冷凝形成的岩石（两种不同的岩浆作用方式）.
	
	岩浆作用形成的常见岩石：如花岗岩、玄武岩等.
	\subsubsection{岩浆的成分与黏性}
	岩浆：熔体+ 结晶矿物
	
	粘性较大的流体，其粘性受许多因素影响，但化学（元素）成分是最重要的因素
	
	元素成分包括$Si, Al, O$以及$Fe, Mg, Ca, Na, K, Ti$等
	
	1)如果岩浆中Si含量相对较低，Fe、Mg等阳离子较多，挥发分少，则粘性小；反之，络阴离子多，则粘性大
	
	2)温度高，粘性小；温度低，粘性大
	
	3)挥发分少，粘性小；挥发分多，粘性大
	
	\subsubsection{岩浆作用的分类}
	喷出作用：熔融的岩浆喷出地表，产生火山现象，及其引起的全部地质过程.
	
	喷出作用形成的岩石称为火山岩.
	
	侵入作用：熔融的岩浆上升，侵入早期形成的岩石，温度逐渐降低而冷却凝固的过程.
	
	侵入作用形成的岩石称为侵入岩.
	\subsubsection{火山作用}
	火山作用是岩浆喷出的一种最主要的表现形式.
	
	害处:地质灾害.
	
	益处:丰富的矿产、肥沃的土壤、丰富的降水.地热资源：火力发电，发展旅游.
	
	形成原因:熔体+ 矿物+气体集聚.
	
	\subsubsection{熔体}
	\textbf{熔体 (Melts) 概述}
	
	熔体是火山喷出物的主体，直接反映了地下岩浆的性质、温度及物理状态。根据 $\text{SiO}_2$ 的质量分数（wt.\%），通常将岩浆/岩石分为以下三类：
	
	\begin{itemize}
		\item \textbf{基性 (Mafic)：} $\text{SiO}_2 < 52\%$
		\item \textbf{中性 (Intermediate)：} $52\% \le \text{SiO}_2 \le 65\%$
		\item \textbf{酸性 (Felsic/Acidic)：} $\text{SiO}_2 > 65\%$
	\end{itemize}
	
	\textbf{典型岩石化学成分表达}
	
	下表列出了酸性代表岩石（花岗岩）与基性代表岩石（玄武岩）的主要氧化物含量对比：
	
	\begin{table}[htbp]
		\centering
		\caption{典型岩石主要氧化物成分表 (wt.\%)}
		\label{tab:rock_composition}
		\begin{tabular}{lrr}
			\toprule
			岩石类型 & 花岗岩 (Granite) & 玄武岩 (Basalt) \\
			\midrule
			\ce{SiO2} (wt.\%) & 73.91 & 46.92 \\
			\ce{TiO2}        & 0.15  & 2.20  \\
			\ce{Al2O3}       & 14.04 & 14.76 \\
			\ce{Fe2O3T}\textsuperscript{*} & 1.23  & 11.98 \\
			\ce{MnO}         & 0.04  & 0.17  \\
			\ce{MgO}         & 0.27  & 8.35  \\
			\ce{CaO}         & 1.18  & 8.91  \\
			\ce{Na2O}        & 3.55  & 3.25  \\
			\ce{K2O}         & 4.07  & 1.06  \\
			\ce{P2O5}        & 0.08  & 0.41  \\
			\midrule
			\textbf{LOI}\textsuperscript{\dag} & 0.82  & 1.84  \\
			\textbf{Total}   & 99.34 & 99.85 \\
			\bottomrule
			\multicolumn{3}{l}{\small * \ce{Fe2O3T} 表示全铁含量。} \\
			\multicolumn{3}{l}{\small \dag \textbf{LOI} (Loss on Ignition)：烧失量，代表岩石中的挥发分 (\ce{H2O} 和 \ce{CO2})。}
		\end{tabular}
	\end{table}
	\subsubsection{基性熔体}
	$SiO_2$含量低，粘度较小，流速快，冷凝后形成玄武岩.
	
	如果表面先结壳，下部熔体继续流动，拖曳表面薄壳，形成波状、绳状熔岩.
	
	如果在火山颈内缓慢冷却，收缩后形成原生节理。若结构均匀，形成无数收缩核心，彼此等距离排列，在垂直于收缩中心连线方向发生破裂，形成横截面为六边形的柱状节理.
	
	如果是水下火山喷发，则形成枕状熔岩.
	
	\subsubsection{酸性熔岩}
	
	SiO2含量高, 粘度大, 难以流动, 挥发分引起上覆硬壳破裂, 形成块状熔岩，如火山角砾岩等.
	
	\subsubsection{固体}
	
	火山喷发时会将破碎的岩块、熔浆喷发到空中, 再以各种形态的碎块降落, 称为火山碎屑.
	
	固体固体：火山块、火山砾、火山灰.
	
	熔体固体：火山弹、熔岩饼.
	
	浮岩、火山渣、黑曜岩等.
	
	\subsubsection{侵入岩}
	
	深部岩浆向上运移并进入上覆岩石，但未达到地表，称为侵入作用.
	
	岩浆在侵入过程中降温冷却，结晶形成的岩石，称为侵入岩.
	
	侵入岩是被周围原有岩石封闭的三维实体，也称为侵入体，包围侵入体的原有岩石称为围岩.
	
	\textbf{侵入方式}
	
	1)浅成侵入
	
	浅成侵入是由于地壳上部近地表处压力小，脆性大，断裂多，岩石间结构较松散，岩浆易于挤入，因此挤入多发生在近地表处，形成浅成侵入体.
	
	地下岩浆处于高温高压状态，向上运移侵入围岩过程中，以巨大的压力沿围岩层面或断裂等薄弱部位挤入围岩中，冷凝形成侵入体.
	
	2)深成侵入
	
	多发生在地壳深部，侵入过程称为深成侵入作用，形成的岩体称为深成侵入体.
	
	侵入过程中以极高的温度熔化围岩，并加以一定的挤压，侵位，冷凝形成的侵入体.
	
	\subsubsection{岩浆岩的分布}
	
	火山和深成岩浆活动受岩石圈的板块构造活动控制.
	
	岩浆岩的分布有一定的规律(洋中脊、俯冲带、热点):环太平洋火山带.地中海、印尼火山带.洋中脊火山带.红海-东非裂谷带.
	
	\subsubsection{岩浆的演化和岩石分类}
	\textbf{岩浆物质来源}
	
	地壳物质熔融（上、下地壳）.
	
	岩石圈地幔物质.
	
	岩石圈地幔+地壳物质.
	
	软流圈地幔物质.
	
	\subsubsection{岩浆演化}
	
	岩浆是具有高温的熔体（>600℃ )，而围岩是相对较冷的岩石.
	
	岩浆结晶分异、岩浆和围岩之间的相互作用，产生多种不同的变化.
	
	\subsubsection{侵入岩主要类型及矿物组成}
	
	花岗岩主要矿物: 石英、钾长石/斜长石、黑云母、角闪石.
	
	闪长岩主要矿物: 角闪石、斜长石、石英、辉石.
	
	辉长岩主要矿物: 辉石、斜长石、橄榄石、角闪石.
	
	橄榄岩主要矿物: 橄榄石、斜方辉石、单斜辉石.
	
	
	\begin{table}[htbp]
		\centering
		\caption{岩浆岩分类简表（依据成分、色率及产状）}
		\label{tab:igneous_rocks}
		\small % 适当缩小字号以适应内容
		\renewcommand{\arraystretch}{1.5} % 增加行高，视觉更清晰
		\begin{tabular}{llcccc}
			\toprule
			\multicolumn{2}{l}{\textbf{岩石类别}} & \textbf{超基性} & \textbf{基性} & \textbf{中性} & \textbf{酸性} \\
			\midrule
			\multicolumn{2}{l}{\textbf{SiO$_2$ 含量}} & $<$ 45\% & 45$\sim$52\% & 53$\sim$65\% & $>$ 65\% \\
			\midrule
			\multicolumn{2}{l}{\textbf{主要矿物组成}} & \makecell[l]{橄榄石、辉石\\角闪石} & \makecell[l]{钙长石、辉石\\角闪石} & \makecell[l]{中长石、角闪石\\碱长石、黑云母} & \makecell[l]{钾长石、钠长石\\石英、黑云母} \\
			\midrule
			\multicolumn{2}{l}{\textbf{色率 (M)}} & $>$ 75 & 75 $\sim$ 35 & 35 $\sim$ 20 & $<$ 20 \\
			\midrule
			\multirow{6}{*}{\textbf{产状}} & \makecell[l]{\textbf{喷出岩}\\\scriptsize (斑状/隐晶质结构\\\scriptsize 气孔/杏仁构造)} & 科马提岩 & 玄武岩 & \makecell{安山岩\\粗面岩} & 流纹岩 \\
			\cmidrule{2-6}
			& \makecell[l]{\textbf{浅成岩}\\\scriptsize (斑状/细粒结构\\\scriptsize 或隐晶质结构)} & 苦橄岩 & 辉绿岩 & \makecell{闪长斑岩\\正长斑岩} & 花岗斑岩 \\
			\cmidrule{2-6}
			& \makecell[l]{\textbf{深成岩}\\\scriptsize (全晶质结构\\\scriptsize 粗粒或似斑状)} & \makecell{橄榄岩\\辉石岩} & 辉长岩 & \makecell{闪长岩\\正长岩} & 花岗岩 \\
			\bottomrule
		\end{tabular}
	\end{table}
	
	\subsubsection{沉积岩}
	\textbf{沉积岩的主要种类}
	
	物理风化产物: 泥岩、砂岩、页岩、角砾岩.
	
	化学风化产物: 石灰岩、白云岩、燧石、条带状铁建造（BIF）、蒸发岩.
	
	原岩（岩浆岩、变质岩或沉积岩）经历风化、剥蚀、搬运和沉积过程形成的岩石.
	
	\subsubsection{沉积岩与环境变化}
	1)沉积环境: 大气、水体等的物理、化学状态和过程.
	
	2)改变大气、海洋的组成：C、O、S、Fe等, 影响地球环境变化.
	\subsubsection{变质岩}
	变质岩与俯冲带地温梯度
	
	1)蓝片岩-榴辉岩：低地温梯度.
	
	2)角闪岩-榴辉岩：中等地温梯度.
	
	3)角闪岩-麻粒岩：高地温梯度.
	
	\subsection{重要科学问题}
	
	\subsubsection{地球磁场的形成?}
	地球磁场是地球宜居性的第一道防线。
	\begin{itemize}
		\item \textbf{形成机制：} 核心在于外地核的\textbf{地发电机效应 (Geodynamo)}。液态外核中的熔融铁在热对流和化学对流的驱动下，结合地球自转产生的科里奥利力，形成复杂的螺旋流，将动能转化为磁能。
		\item \textbf{驱动力：} 内核的缓慢凝固释放潜热和轻元素，为外核对流提供持久动力。
	\end{itemize}
	
	\subsubsection{地球大陆地壳的形成与演化?}
	大陆地壳是地球区别于其他类地行星的关键特征。
	\begin{itemize}
		\item \textbf{早期形成：} 早期陆壳主要通过俯冲带含水玄武岩的部分熔融形成 \textbf{TTG 岩套}。
		\item \textbf{演化特征：} 从太古宙的垂直构造（地壳滴落/地幔柱）向显生宙的水平板块构造演化。陆壳成分逐渐从钠质（Na-rich）向钾质（K-rich）转变，反映了地热梯度的降低。
	\end{itemize}
	
	\subsubsection{地球大气与海洋的成分变化?}
	地球表层系统的化学性质经历了剧变。
	\begin{table}[htbp]
		\centering
		\caption{地球大气与海洋的成分}
		\begin{tabularx}{\textwidth}{l X X}
			\toprule
			时间节点 & 大气/海洋特征 & 关键地质/生物过程 \\
			\midrule
			\textbf{冥古宙} & 富 $CO_2, CH_4$, 还原性环境 & 岩浆洋脱气；原始海洋形成。 \\
			\textbf{太古宙} & 极低氧含量；海洋富铁 & 早期光合作用开始，产氧量低。 \\
			\textbf{元古宙} & 大氧化事件 (GOE) & 自由氧出现；海洋硫化或氧化转变。 \\
			\textbf{显生宙} & 现代水平氧浓度 (21\%) & 陆生植物出现；碳-氧平衡趋于稳定。 \\
			\bottomrule
		\end{tabularx}
	\end{table}
	
	\subsubsection{大陆动力学？}
	大陆动力学研究大陆岩石圈在力的作用下的变形与重组。
	\begin{itemize}
		\item \textbf{核心特征：} 大陆地壳具有高度的非均质性和流变性。碰撞造山作用（如青藏高原）涉及岩石圈缩短、地壳增厚、重力不稳定性导致的下沉（Delamination）以及深部地幔流的相互作用。
		\item \textbf{威尔逊循环：} 描述了大陆裂解、洋盆形成、洋盆关闭到大陆碰撞的旋回过程。
	\end{itemize}
	
	\subsubsection{地球内部运作机制及表生响应?}
	地球是一个深度耦合的整体。
	\begin{itemize}
		\item \textbf{内部驱动：} 地幔对流驱动板块运动，超地幔柱触发大规模火山活动（如 LIPs）。
		\item \textbf{表生响应：} 地震、火山喷发直接改变地表地貌。同时，板块俯冲将碳带入深部，而火山作用将碳释放回大气，构成了百万年尺度的\textbf{碳循环气候调节器}。
	\end{itemize}
	
	\subsubsection{地球宜居性如何形成?}
	地球之所以宜居，是因为建立了以下机制的动态平衡：
	\begin{enumerate}
		\item \textbf{磁场屏蔽：} 防止太阳风剥蚀大气层，保护水分子。
		\item \textbf{板块构造与碳循环：} 通过硅酸盐风化反馈调节温室效应，维持液态水的温度范围。
		\item \textbf{适宜的质量：} 足够的引力束缚大气层，同时也允许轻元素的适度逃逸。
		\item \textbf{月球的稳定作用：} 维持地球自转轴的倾角，确保气候的长期相对稳定。
	\end{enumerate}
	
	\section{地球与天体化学}
	\subsection{定年: 相对年龄}
	相对年龄：确定岩石或地质事件的先后，不需知道它们的确切年龄.
	
	\textbf{地层顺序}
	
	1)叠加原理：未受扰动的沉积地层或岩层，下部层位老于上部层位.
	
	2)水平层理：绝大部分沉积岩石初始为水平的.
	
	3)侧向延续性：沉积盆地中，地层在各个方向总体上是连续的.
	
	4)穿插关系：被切穿的岩石相对古老.
	
	5)包裹关系：如果一个岩石包裹了另外一个，则包裹体更年轻.
	
	6)化石层序：地层中指示时代的独特产出的化石.
	
	\subsection{定年: 绝对年龄}
	
	绝对年龄：放射性同位素定年.
	
	$\alpha$衰变：释放$\alpha$粒子.
	
	$\beta$衰变：释放负电子.
	
	重核裂变：U-Th裂变.
	
	\subsection{放射性同位素衰变方程}
	单位时间内放射性母体同位素衰变减少的量与放射性母体的原子数成正比(一阶反应定律).
	
	$N$为$t$时刻母体同位素的原子数；$dN/dt$为$t$时刻母体同位素的衰变产率.
	
	$$N=N_0e^{-\lambda t}$$
	
	设衰变产生的子体同位素的原子数为$D^*$.
	
	放射成因同位素子体: $$D=D^*+D_0$$
	其中$D_0$是$t=0$时体系中子体同位素的量。
	联立可得：$$D=D_0+N(e^{\lambda t}-1)$$
	
	衰变常数与半衰期
	$$t_{\frac{1}{2}}= \frac{ln2}{\lambda}$$
	
	质谱仪分析同位素比值的精度远高于测定单个同位素绝对含量的精度.
	
	常选取子体同位素中其它稳定同位素（$D_s$）作参照：
	$$\frac{D}{D_s}=\left(\frac{D}{D_s}\right)_{t=0}+\frac{N}{D_s}(e^{\lambda t}-1)$$
	
	\subsubsection{Rb-Sr 同位素体系}
	\begin{itemize}
		\item \textbf{衰变方程：} $^{87}\text{Rb} \rightarrow ^{87}\text{Sr} + \beta^- + \nu + E$ （$\nu$为反中微子，$E$为能量）
		\item \textbf{等时线方程：} 
		$\left( \frac{^{87}\text{Sr}}{^{86}\text{Sr}} \right)_s = \left( \frac{^{87}\text{Sr}}{^{86}\text{Sr}} \right)_0 + \left( \frac{^{87}\text{Rb}}{^{86}\text{Sr}} \right)_s (e^{\lambda t} - 1)$
		\item \textbf{参数含义：} $\lambda$ 为 $^{87}\text{Rb}$ 的衰变常数；参照同位素为 $^{86}\text{Sr}$。
	\end{itemize}
	
	同位素等时线定年三原则：同源、同时、无扰动
	
	\begin{tikzpicture}[>=stealth, thick, scale=1.2]
		% 坐标轴
		\draw[->] (0,0) -- (7,0) node[right] {$x = \frac{^{87}\text{Rb}}{^{86}\text{Sr}}$ (母子比)};
		\draw[->] (0,0) -- (0,5) node[above] {$y = \frac{^{87}\text{Sr}}{^{86}\text{Sr}}$ (子体比)};
		
		% 初始状态 t=0 (水平线)
		\draw[dashed, blue] (0,1.5) -- (6,1.5) node[right] {形成初期 ($t=0$)};
		\fill[blue] (1,1.5) circle (2pt) node[below] {矿物A};
		\fill[blue] (3,1.5) circle (2pt) node[below] {矿物B};
		\fill[blue] (5,1.5) circle (2pt) node[below] {矿物C};
		\node[left, blue] at (0,1.5) {$\left(\frac{^{87}\text{Sr}}{^{86}\text{Sr}}\right)_i$};
		
		% 现今状态 t > 0 (斜线)
		\draw[red, ultra thick] (0,1.5) -- (5.5,4.5) node[above right] {等时线 ($t > 0$)};
		\fill[red] (0.8,1.9) circle (2.5pt);
		\fill[red] (2.4,2.8) circle (2.5pt);
		\fill[red] (4.0,3.7) circle (2.5pt);
		
		% 演化箭头
		\draw[->, gray] (1,1.5) -- (0.8,1.8);
		\draw[->, gray] (3,1.5) -- (2.4,2.7);
		\draw[->, gray] (5,1.5) -- (4,3.6);
		\node[gray, font=\small, rotate=45] at (4.7,2.5) {随时间衰变演化};
		
		% 斜率标注
		\draw (4,3.7) -- (5,3.7) -- (5,4.25);
		\node[right] at (5,4) {斜率 $m = e^{\lambda t} - 1$};
		
		% 说明文字
		\node[align=left, draw, fill=yellow!10] at (3.5,0.8) {
			\small 1. 截距 $\rightarrow$ 物质来源 (初始比值) \\
			\small 2. 斜率 $\rightarrow$ 形成年龄 ($t$)
		};
	\end{tikzpicture}
	\subsubsection{Sm-Nd 同位素体系}
	\begin{itemize}
		\item \textbf{衰变方程：} $^{147}\text{Sm} \rightarrow ^{143}\text{Nd} + \alpha + E$
		\item \textbf{等时线方程：} 
		$\left( \frac{^{143}\text{Nd}}{^{144}\text{Nd}} \right)_s = \left( \frac{^{143}\text{Nd}}{^{144}\text{Nd}} \right)_0 + \left( \frac{^{147}\text{Sm}}{^{144}\text{Nd}} \right)_s (e^{\lambda t} - 1)$
		\item \textbf{主要用途：} 主要用于\textbf{全岩}或\textbf{全岩+矿物}等时线法测年。
	\end{itemize}
	
	\subsubsection{Lu-Hf 同位素体系}
	\begin{itemize}
		\item \textbf{衰变方程：} $^{176}\text{Lu} \rightarrow ^{176}\text{Hf} + \beta^- + \nu + E$
		\item \textbf{等时线方程：} 
		$\left( \frac{^{176}\text{Hf}}{^{177}\text{Hf}} \right)_s = \left( \frac{^{176}\text{Hf}}{^{177}\text{Hf}} \right)_0 + \left( \frac{^{176}\text{Lu}}{^{177}\text{Hf}} \right)_s (e^{\lambda t} - 1)$
		\item \textbf{参数含义：} 参照同位素为 $^{177}\text{Hf}$；$\lambda$ 为 $^{176}\text{Lu}$ 的衰变常数。
	\end{itemize}
	
	\subsubsection{锆石U-Pb 同位素体系 (双衰变体系)}
	\textbf{锆石的优势:}
	
	在中酸性岩石中广泛出现;抗风化，耐熔融;高U，低Pb，易分析;能保存多期环带特征;O-Hf等多种同位素示踪.
	
	U-Pb体系包含两个相互独立的衰变链，可建立两个联立方程及一个 Pb-Pb 方程：
	\begin{enumerate}
		\item $^{238}\text{U} \rightarrow ^{206}\text{Pb} + 8\alpha + 6\beta^- + E$ \quad $\lambda_1$
		\item $^{235}\text{U} \rightarrow ^{207}\text{Pb} + 7\alpha + 4\beta^- + E$ \quad $\lambda_2$
	\end{enumerate}
	\textbf{Pb-Pb 等时线方程：}
	\begin{equation}
		\frac{\left( \frac{^{207}\text{Pb}}{^{204}\text{Pb}} \right)_s - \left( \frac{^{207}\text{Pb}}{^{204}\text{Pb}} \right)_0}{\left( \frac{^{206}\text{Pb}}{^{204}\text{Pb}} \right)_s - \left( \frac{^{206}\text{Pb}}{^{204}\text{Pb}} \right)_0} = \frac{1}{137.88} \cdot \frac{e^{\lambda_2 t} - 1}{e^{\lambda_1 t} - 1}
	\end{equation}
	
	\subsubsection{U-Th-Pb 衰变系列与中间核素}
	铀-钍衰变系列包括多个中间短半衰期核素，常用于晚第四纪定年（非平衡法）。
	
	\textbf{衰变链总结：}
	\begin{itemize}
		\item $^{238}\text{U} \Rightarrow ^{230}\text{Th} \Rightarrow ^{226}\text{Ra} \Rightarrow ^{206}\text{Pb}$
		\item $^{235}\text{U} \Rightarrow ^{231}\text{Pa} \Rightarrow ^{207}\text{Pb}$
		\item $^{232}\text{Th} \Rightarrow \dots \Rightarrow ^{208}\text{Pb}$
	\end{itemize}
	\subsection{放射成因同位素示踪}
	地球不同储库中元素分配
	
	陆壳: $Rb/Sr$高，$Sm/Nd$低，$Lu/Hf$低.
	
	沉积物: $Rb/Sr$高，$Sm/Nd$低，$Lu/Hf$高.
	
	洋中脊玄武岩(MORB) $Rb/Sr$低，$Sm/Nd$高，$Lu/Hf$高.
	
	\subsubsection{硫（S）同位素体系概览}
	硫在自然界中有四个稳定的同位素，其原子序数为 16，质量数如下：
	\begin{itemize}
		\item \textbf{主要同位素：} \ce{^{32}S} (95.0\%), \ce{^{33}S} (0.75\%), \ce{^{34}S} (4.21\%), \ce{^{36}S} (0.02\%)。
		\item \textbf{测量标准：} 通常使用 $\delta^{34}\text{S}$ 来表示 \ce{^{34}S}/\ce{^{32}S} 相对于标准物质（如 VCDT）的千分偏差。
	\end{itemize}
	
	\subsubsection{两种分馏机制：MDF 与 MIF}
	
	\textbf{质量相关分馏 (Mass-Dependent Fractionation, MDF)}
	绝大多数地质、化学和生物过程（如微生物硫酸盐还原）产生的分馏。
	\begin{itemize}
		\item \textbf{物理基础：} 分馏程度严格取决于同位素之间的质量差异。
		\item \textbf{数学关系：} 由于 \ce{^{34}S} 与 \ce{^{32}S} 的质量差约为 \ce{^{33}S} 与 \ce{^{32}S} 质量差的两倍，存在线性关系：
		\begin{equation}
			\delta^{33}\text{S} \approx 0.515 \times \delta^{34}\text{S}
		\end{equation}
		\item \textbf{特征：} 在 $\delta^{33}\text{S}$ vs $\delta^{34}\text{S}$ 图谱上，所有数据点均落在一条斜率约为 0.515 的回归线上。
	\end{itemize}
	
	\textbf{质量不相关分馏 (Mass-Independent Fractionation, MIF)}
	不遵循质量差异比例的异常分馏现象。
	\begin{itemize}
		\item \textbf{物理基础：} 主要由大气高层中 \ce{SO2} 的\textbf{紫外线（U.V.）光化学反应}引起。
		\item \textbf{指示指标 ($\Delta^{33}\text{S}$)：} 定义样点偏离 MDF 线程度的参数：
		\begin{equation}
			\Delta^{33}\text{S} = \delta^{33}\text{S} - 1000 \times \left[ \left( 1 + \frac{\delta^{34}\text{S}}{1000} \right)^{0.515} - 1 \right]
		\end{equation}
		\item \textbf{特征：} $\Delta^{33}\text{S} \neq 0$（显著偏离零点）。
	\end{itemize}
	
	\subsubsection{大气氧含量的示踪原理}
	
	硫同位素的 MIF 信号是判断地球早期大气是否含氧的“金标准”：
	
	\begin{table}[htbp]
		\centering
		\caption{大气环境与硫同位素信号对照}
		\begin{tabular}{p{3cm} p{5cm} p{5cm}}
			\toprule
			\textbf{特征} & \textbf{还原性大气 (太古宙)} & \textbf{氧化性大气(现代/GOE后)} \\
			\midrule
			\textbf{氧气含量} & 极低，无臭氧层保护 & 升高，存在臭氧层 \\
			\textbf{紫外线强度} & U.V. 直达低层大气，诱发 \ce{SO2} 光解 & 臭氧层吸收 U.V.，光解反应受阻 \\
			\textbf{硫的保存} & \ce{S}、\ce{H2S} 迅速沉降并保留 MIF 信号 & 硫被迅速氧化为硫酸盐，信号抹除 \\
			\textbf{$\Delta^{33}\text{S}$ 表现} & \textbf{显著偏离 (MIF signal preserved)} & \textbf{趋于零点 (No MIF signal)} \\
			\bottomrule
		\end{tabular}
	\end{table}
	
	\subsubsection{大氧化事件 (GOE) 的地质记录}
	根据全球沉积岩样本记录（见右侧散点图）：
	\begin{itemize}
		\item \textbf{2.45 Ga 以前（太古宙）：} $\Delta^{33}\text{S}$ 呈现巨大的波动（-4‰ 到 +8‰），证明当时大气缺乏氧气，光化学 MIF 信号能够被沉积物记录。
		\item \textbf{2.4 Ga 以后：} $\Delta^{33}\text{S}$ 迅速收缩至 $\pm 0.1$‰ 范围内。这标志着\textbf{大氧化事件 (Great Oxidation Event)} 的发生，大气氧含量升高导致 MIF 机制失效。
	\end{itemize}
	\subsection{微量元素在矿物-熔体中的分异}
	
	固相-液相之间微量元素分配系数 $D_{sol/liq}$
	
	1)$D>1$, 固相中相容
	
	2)$D<1$, 固相中不相容
	
	微量元素配分的控制因素:离子半径、价态
	
	\subsection{月球最年轻火山活动成因示踪}
	
	\subsubsection{月球的火山活动为什么持续如此之久？}
	月球质量约为地球质量的1.2\%，内动力小.
	
	1)升温（深部热源）$\Rightarrow$ 源区生热元素含量（Th）
	
	2)加水或挥发分（降低熔点）$\Rightarrow$ 源区水（挥发分）含量
	
	3)降压（减压熔融）$\Rightarrow$ 反演岩浆起源压力
	
	\subsection{同位素判断澄阳湖大闸蟹}
	
	利用Sr同位素示踪中华绒螯蟹产地.
	
	1)同一只中华绒螯蟹不同部位的Sr同位素组成一致.
	
	2)来自同一产地的中华绒螯蟹Sr同位素组成一致.
	
	3)不同产地的中华绒螯蟹则具有不同的Sr同位素组成.
	
	\subsection{稳定同位素示踪}
	
	\textbf{稳定同位素}
	
	$O-H-N-S$等,$Mg-Fe-Cu-Zn$等.
	
	\textbf{稳定同位素分馏}
	
	热力学平衡分馏: 测温、示踪; 动力学分馏: 动力学过程.
	\subsubsection{原理}
	不同储库的某种稳定同位素组成差异显著.
	
	特征性的稳定同位素信号.
	
	O同位素为例:流体（大气降水、海水）;地表沉积物、洋壳;幔源玄武岩.
	
	地球样品中S同位素非质量分馏随时间变化,反映大气中氧含量升高.
	
	\subsubsection{基本动力学框架}
	
	\textbf{地球圈层}
	地壳、上地幔、地幔过渡带、下地幔、外核、内核;地壳俯冲;地幔对流;岛弧、大洋中脊、洋岛层圈间物质传输
	
	\subsubsection{揭示前太阳系物质}
	碳质球粒陨石中硅酸盐矿物(橄榄石、辉石、石英/柯石英)具有极其异常的O同位素比值, 指示前太阳系来源.
	
	前太阳系的物质中硅酸盐比非硅酸盐颗粒更广泛.
	
	\subsubsection{碳同位素示踪植被变化与古人类演化}
	\textbf{C3植物}: 小麦，水稻，大豆，苦竹，烟草，棉花，银杏，冷杉，银杉.
	
	光合作用简单，更为原始，大部分乔木、灌木及喜冷的草本植物.
	
	\textbf{C4植物}：玉米，高粱，甘蔗，狗尾草，白茅，小米(粟).
	
	光合作用复杂，效率高，包含一些特定的喜暖草本植物等.
	
	C3植物的稳定碳同位素$\delta^{13}C$值为$-2.2\%\sim-3.0\%$，平均$-2.7\%$.
	
 	C4植物的稳定碳同位素$\delta^{13}C$值为$-0.9\%\sim-1.9\%$，平均$-1.4\%$.
	
	\textbf{地质记录}
	沉积环境中的碳酸盐结核（carbonate nodule）;叶蜡（leaf wax）;哺乳动物牙齿化石.
	
	\subsection{稳定同位素测温}
	\textbf{稳定同位素分馏与温度的关系}
	
	在平衡条件下，两种物相 $A$ 和 $B$ 之间的氧同位素分馏可以用以下公式描述：
	
	\begin{align}
		\Delta ^{18}\text{O}_{\text{A-B}} &= \delta ^{18}\text{O}_{\text{A}} - \delta ^{18}\text{O}_{\text{B}} \nonumber \\
		&= 1000 \ln \alpha_{\text{A-B}} = f(T) \nonumber \\
		&= \frac{A \times 10^3}{T} + \frac{B \times 10^6}{T^2} + C
	\end{align}
	
	\noindent 其中：
	\begin{itemize}
		\item $A$ 和 $B$ 代表平衡共生的两种不同物相（如两种矿物）。
		\item $T$ 为热力学温度（单位：$\text{Kelvin}$）。
		\item $\alpha$ 为分馏系数，$A, B, C$ 是由实验确定的常数。
	\end{itemize}
	
	\textbf{高温简化模型}
	在高温环境（如岩浆作用或深部变质作用）下，上式通常可以简化为：
	\begin{equation}
		\delta ^{18}\text{O}_{\text{A}} - \delta ^{18}\text{O}_{\text{B}} \approx \frac{A \times 10^6}{T^2}
	\end{equation}
	\subsubsection{榴辉岩相变质温度}
	
	榴辉岩中主要矿物: 石英，石榴石，绿辉石，多硅白云母，金红石.
	
	\subsubsection{要用多种方法确定岩石记录的温度}
	
	每种方法的局限;交叉验证;不同测温对象记录的事件、反映的地质意义不同.
	
	\subsubsection{给恐龙测体温}
	
	内温性动物：自主地把自体内温度稳定在相对固定范围，受环境影响不大。典型代表：大部分哺乳动物和鸟类。
	
	外温性：主要依靠外部热量来调节体温，体温多随环境温度变化。典型代表：蜥蜴、鳄鱼和大部分龟鳖类。
	
	中温性：体温调节能力介于两者之间，往往通过自身产的热量维持体温高于环境温度，或者维持在某个区间。但不像内温性生物将体温维持到一个很高且稳定的程度。现存动物的代表很少，包括棱皮龟、单孔类哺乳动物等。
	
	巨温性：体型大，体表表面积相对变小，散热能力相对小体型动物下降，导致它们新陈代谢能力不如典型的内温性动物，体温维持在很高的水平。
	
	\textbf{方法1}： 测量恐龙牙釉质(磷酸盐中碳酸盐).
	
	牙齿的温度，真的代表这种动物的体温么？
	
	---牙齿温度，只能说是口腔温度，不能代表恐龙的体温.
	
	\textbf{方法2}：测量恐龙蛋
	
	恐龙蛋形成于雌性动物体腔内最核心的位置.
	
	只要测到了蛋壳的温度，就能够推断出动物内环境的温度.
	
	\textbf{问题:}
	
	1)取样少；
	
	2)讨论一种生物是否是内温动物的时候，关注的是“保存体温高于环境温度的能力”，而不是“体温的绝对值”.
	
	化石都来自白垩纪中低纬度地区，白垩纪的环境温度本来就偏高，测量出来的较高温度可能不是因为这些恐龙是内温性动物，而是因为环境很热.
	
	3)埋藏和成岩作用影响：生物体死后形成化石，经历漫长的埋藏和成岩作用。如果化石在埋藏和成岩过程中经历了重结晶作用，会重新结晶形成新的碳酸盐，记录的温度可能不能反映当时的体温.
	
	\textbf{方法3}：进一步测试恐龙蛋
	
	1)恐龙蛋化石：高纬度地区.
	
	2)重结晶影响评估：蛋壳形态学、元素示踪和阴极发光显微镜等研究手段.
	
	3)同时对同地层的软体动物和植物用同样的方法测温，来代表环境温度，进而求得环境温度和体温数据的差值.
	
	\section{大陆、生命与地球演化}
	
	\subsection{地球在宇宙中的位置与宜居性}
	\begin{itemize}
		\item 地球是太阳系中目前已知唯一宜居行星。
		\item 太阳系探测历史已超过60年，包括旅行者1号进入星际空间（2012年）和新视野号飞越冥王星（2015年）。
		\item 已确认系外行星数量超过5000颗（截至2022年），2019年诺贝尔物理学奖表彰系外行星的发现。
		\item 韦伯空间望远镜已实现直接拍摄系外行星图像。
	\end{itemize}
	
	\subsection{地球的形成与早期演化（约46亿年前—40亿年前）}
	\begin{itemize}
		\item 地球起源于太阳星云，经历微行星碰撞形成行星胚胎，最终成为类地行星。
		\item 早期地球为“岩浆洋”状态，表面持续遭受小行星撞击。
		\item 月球形成于约45.1亿年前的一次大碰撞：火星大小的天体斜向撞击原始地球。
		\item 嫦娥五号揭示月球岩浆活动持续至约20亿年前，月幔含水量极低（1–5 μg/g）。
		\item 核幔分异发生在约44.5–44.4亿年前，形成液态外核与固态内核，产生地磁场，保护大气与生命。
		\item 最古老岩石为加拿大Acasta片麻岩（约40亿年），最古老矿物为澳大利亚Jack Hills锆石（约44亿年）。
		\item 锆石氧同位素证据表明44亿年前地球表面已有液态水和小型大陆。
		\item 地磁场可能始于42亿年前，40–36亿年前强度显著下降。
		\item 晚期大碰撞（约40–39亿年前）可能造成海洋蒸发与生命灭绝。
	\end{itemize}
	
	\subsection{水的来源与早期生命}
	\begin{itemize}
		\item 地球水主要来自内部，氢同位素证据排除彗星为主要来源。
		\item 地幔过渡带矿物（如林伍德石）含水，表明地球内部为重要水储库。
		\item 最早生命痕迹出现于约37–35亿年前的叠层石（蓝藻光合作用形成）。
		\item “黯淡太阳悖论”：早期太阳辐射较弱，温室气体（$CO_2, CH_4$）维持地球温度，支持液态水与生命。
	\end{itemize}
	
	\subsection{板块构造与大气演化}
	\begin{itemize}
		\item 约25亿年前，地球构造从垂向“滴落构造”转为水平板块构造。
		\item 板块运动促进物质与能量交换，推动地球冷却与宜居性演化。
		\item 大氧化事件（约24亿年前）：蓝藻光合作用使大气氧含量上升，导致厌氧生物大量灭绝。
		\item 条带状铁建造（BIF，25–18亿年）为重要铁矿来源。
		\item 最早“雪球地球”事件发生在24–21亿年前，冰反射反馈导致全球冰冻。
	\end{itemize}
	
	\subsection{生命演化关键节点}
	\begin{itemize}
		\item 真核生物出现于约21–16亿年前，线粒体内共生发生于约20亿年前，叶绿体内共生约15亿年前。
		\item 有性生殖始于约15亿年前，加速生命多样化。
		\item 超大陆（如罗迪尼亚）聚合与裂解影响全球气候与生命分布。
		\item 寒武纪生命大爆发（约5.4亿年前）出现多门类硬体生物，澄江生物群（5.2亿年前）与清江生物群（5.18亿年前）保存完整生态系统。
		\item 植物登陆（约4.8–4.2亿年前）改造大陆地表与沉积系统。
		\item 种子植物出现（约3.9亿年前）提升传播能力。
		\item 动物登陆（约4亿年前），节肢动物率先适应陆地。
		\item 石炭纪（3.6–3.0亿年前）形成大量煤炭，羊膜动物出现使动物脱离水生环境。
	\end{itemize}
	
	\subsection{超级火山、生物大灭绝与气候变迁}
	\begin{itemize}
		\item 中国峨眉山（2.58亿年前）与西伯利亚（2.50亿年前）大火成岩省引发全球环境剧变。
		\item 二叠纪末生物大灭绝（2.52亿年前）导致约95\%海洋物种灭绝，与超级火山活动密切相关。
		\item 多次生物大灭绝与火山喷发、气候变化、海洋酸化与缺氧有关。
		\item 恐龙兴起于三叠纪（约2.3亿年前），部分恐龙具备体温调节能力。
		\item 羽毛演化促进恐龙向鸟类过渡（如中华龙鸟，约1.55亿年前）。
		\item 白垩纪末（6600万年前）恐龙灭绝与德干火山喷发及陨石撞击有关。
	\end{itemize}
	
	\subsection{新生代以来地球演化}
	\begin{itemize}
		\item 印度—欧亚板块碰撞（约7000万年前至今）形成喜马拉雅山脉。
		\item 古新世—始新世极热事件（约5500万年前）导致全球升温与生物更替。
		\item 人类演化：能人（约200万年前）、北京人（约70万年前）、智人（约20万年前）。
		\item 米兰科维奇循环（地球轨道周期）调控冰期—间冰期交替，主导气候变迁。
	\end{itemize}
	
	\subsection{现代人类面临的挑战}
	\begin{itemize}
		\item 当前CO₂浓度达近80万年最高（约400 ppm），全球变暖导致极端天气频发（如2021年郑州暴雨）。
		\item 环境污染（水、土壤、大气）、资源短缺（能源、矿产）与生态健康问题日益严峻。
		\item 碳中和目标（中国：2030年前碳达峰，2060年前碳中和）成为全球共识。
	\end{itemize}
	
	\subsection{地球科学研究方法与前沿问题}
	\begin{itemize}
		\item 研究对象包括陨石、月壤、深部岩石、冰芯、大气等天然与人工样品。
		\item 研究手段涵盖地球物理观测、地球化学测试、高温高压实验与计算模拟。
		\item 前沿科学问题涉及地球形成、生命起源、板块构造、气候变化、灾害预测等。
		\item 地球系统科学强调岩石圈、大气圈、水圈、生物圈、冰冻圈与电磁圈的交叉融合。
	\end{itemize}
	
	\subsection{总结：地球作为宜居星球的独特性}
	\begin{itemize}
		\item 地球具备适宜的能量来源、化学成分、防护体系（磁场、大气）与稳定的地质活动。
		\item 与金星（温室效应强烈）和火星（寒冷稀薄）相比，地球环境适宜生命长期演化。
		\item 地球演化是一部涵盖天体形成、地质变迁与生命进化的壮丽史诗。
	\end{itemize}
	
	\section{地球和类地行星的构造运动}
	
	板块构造运动学说:大陆漂移,海底扩张,地幔对流,板块构造运动.
	\subsection{大陆漂移的证据}
	\subsubsection{形状复杂的板块拟合}
	1)运用大陆架后板块拟合得很好.
	
	2)200百万年前大陆开始分裂.
	\subsubsection{化石}
	\textbf{a) 中龙属Mesosaurus}
	
	石炭二叠纪;淡水爬行动物;能游泳, 但不能越过海洋;南美和非洲都有.
	
	\textbf{b) 水龙兽Lystrosaurus}
	
 	陆地爬行动物;不会游泳;非洲、印度和南极地区都有发现.
 	
	\textbf{c) 犬颌兽属Cynognathus}
	
	陆地爬行动物;不会游泳;非洲和南美.
	
	\textbf{d) 舌羊齿属Glossopteris}
	
	蕨类植物;无花，种子太大，风不能传播;生长区域由植物蔓延的能力决定;非洲, 印度, 南美, 澳大利亚和南极大陆都有发现.
	
	\subsubsection{岩石}
	a) 大西洋两岸的岩石组成、年龄、地质结构相似.
	
	b) 当两边的大陆连在一起，地质现象具有连续性.
	
	\subsubsection{山脉}
	a) 大西洋两岸的山脉具有连续性.
	
	b) 如果北美州和欧洲连在一起，山脉具有连续性 ，岩石具有相似的组成和年龄.
	
	\subsubsection{古气候数据}
	a) 现今的非热带地区有热带气候的证据.
	
	b) 现今的非沙漠地区发现低纬度才有的干旱沙漠证据.
	
	c) 现今的热带气候地区发现冰川的证据.
	
	\subsection{海底扩张}
	
	\subsubsection{核心定义与基本思想}
	海底扩张学说认为，洋中脊（Mid-Ocean Ridge）是地幔物质上升的出口。
	\begin{itemize}
		\item \textbf{物质循环：} 地幔物质在洋中脊处涌出，冷却固结形成新的大洋地壳。
		\item \textbf{传送带机制：} 随着新地壳的不断产生，老的地壳被推向两侧，像传送带一样移动，最终在海沟（Trench）处消亡，回到地幔。
		\item \textbf{动力源：} 核心动力源自地幔对流（Mantle Convection）。
	\end{itemize}
	
	\subsubsection{关键科学证据}
	
	\textbf{磁异常条带 (Magnetic Stripes)}
	这是海底扩张学说最决定性的证据，由 \textbf{Vine-Matthews-Morley (VMM) 假说} 解释：
	\begin{itemize}
		\item 地球磁场会周期性发生倒转。
		\item 洋中脊涌出的岩浆在冷却过居里点时，记录了当时的地磁场方向。
		\item \textbf{特征：} 在洋中脊两侧，磁异常条带呈现出完美的平行对称分布。
	\end{itemize}
	
	\textbf{大洋沉积物特征}
	通过深海钻探（如“格罗玛·挑战者”号）证实：
	\begin{itemize}
		\item \textbf{年龄分布：} 越靠近洋中脊，洋壳年龄越轻；越远离，年龄越老（最老洋壳约 1.8 亿年）。
		\item \textbf{厚度分布：} 洋中脊处几乎无沉积物，随距离增加，沉积物厚度稳步增加。
	\end{itemize}
	
	\textbf{热流值与地震分布}
	\begin{itemize}
		\item \textbf{高热流：} 洋中脊处具有极高的热流值，反映了深部岩浆的上升。
		\item \textbf{地震特征：} 洋中脊处多为浅源地震，反映了拉张应力环境。
	\end{itemize}
	
	\subsubsection{海底扩张速率计算}
	扩张速率可以通过磁异常条带的宽度与其对应的地磁极性倒转时间标尺计算：
	\[	v = \frac{L}{t} \]
	其中 $v$ 为扩张速率，$L$ 为样点到洋中脊的距离，$t$ 为该样点磁记录对应的年龄。
	\textit{注：通常区分“半速度”（单侧移动速度）和“全速度”（两侧背离速度）。}
	
	\subsubsection{洋中脊类型对比}
	
	\begin{table}[htbp]
		\centering
		\caption{快速扩张与慢速扩张洋中脊特征对比}
		\label{tab:ridge_compare}
		\begin{tabularx}{\textwidth}{l l l}
			\toprule
			特征指标 & 慢速扩张 (如大西洋中脊) & 快速扩张 (如东太平洋海隆) \\
			\midrule
			\textbf{全速率} & 1--5 cm/yr & 9--18 cm/yr \\
			\textbf{形貌} & 具有深邃的中央裂谷 & 地形平坦，呈穹窿状 \\
			\textbf{岩浆供应} & 间歇性、供应不足 & 持续性、供应充裕 \\
			\textbf{断裂特征} & 发育大量正断层 & 断裂作用较不显著 \\
			\bottomrule
		\end{tabularx}
	\end{table}
	
	\subsubsection{总结}
	海底扩张学说不仅解决了地壳更新的机制问题，还证明了大洋地壳是年轻的且处于动态循环中。它完美连接了大陆漂移说（过去）与板块构造学说（现代），使地球科学进入了全新的动力学时代。
	
	\subsection{地幔对流}
	\textbf{地幔对流}是指地幔物质由于温度分布不均导致的密度差异，在引力场中产生的缓慢爬行流动。
	\begin{itemize}
		\item \textbf{热源：} 主要来源于地球形成时的余热、放射性元素（U, Th, K）衰变释放的热量以及内核凝固释放的潜热。
		\item \textbf{流变学前提：} 虽然地幔在地震波时间尺度下表现为固体，但在地质时间尺度（百万年）下表现为高粘度的流体（粘度约为 $10^{21} \text{--} 10^{23} \text{ Pa}\cdot\text{s}$）。
	\end{itemize}
	
	\subsubsection{驱动板块运动的力学机制}
	地幔对流通过以下几种力的耦合驱动板块运动：
	
	\begin{enumerate}
		\item \textbf{板片拉力 (Slab Pull)：} 俯冲板片因温度低、密度大，在重力作用下向深部下沉，产生向下的拉力。这是公认的\textbf{最主要驱动力}。
		\item \textbf{脊推力 (Ridge Push)：} 洋中脊位置较高，扩张后的洋壳随温度降低而增厚、下沉，在重力作用下产生向两侧的推力。
		\item \textbf{底部拖拽力 (Basal Drag)：} 对流的地幔物质通过粘性耦合作用在岩石圈底部产生的摩擦力。
	\end{enumerate}
	
	
	
	\subsubsection{地球物理观测证据}
	\begin{itemize}
		\item \textbf{地震层析成像 (Seismic Tomography)：} 揭示地幔内部的波速异常（冷/热异常），勾勒出俯冲板片和上升地幔柱的形态。
		\item \textbf{大地水准面异常：} 反映了地幔内部密度不均匀性导致的引力场变化。
		\item \textbf{热流测量：} 验证了能量从内部向表层传递的速率。
	\end{itemize}
	\subsubsection{对流模式的争议与模型}
	如下表所示:
	\begin{table}[htbp]
		\centering
		\caption{地幔对流主要模型对比}
		\label{tab:convection_models}
		\begin{tabularx}{\textwidth}{l X X}
			\toprule
			模型名称 & 核心观点 & 地球物理证据 \\
			\midrule
			\textbf{全地幔对流} & 对流单元贯穿 660 km 界面，板块可直接俯冲至核幔边界。 & 地震层析成像显示俯冲板片进入下地幔。 \\
			\addlinespace
			\textbf{分层对流} & 以 660 km 不连续面为界，上下地幔各自形成独立的对流循环。 & 解释了上下地幔之间明显的化学组成差异。 \\
			\addlinespace
			\textbf{混合对流} & 存在间歇性的贯穿，部分物质越过界面，部分被阻挡。 & 现代主流观点，兼顾了地震学与地球化学证据。 \\
			\bottomrule
		\end{tabularx}
	\end{table}
	
	\subsection{板块构造的基本单元及特征}
	
	板块构造学说认为地球表层是由若干坚硬的板块组成，它们在塑性较强的软流圈上进行大规模运动。以下是该系统的核心单元与关键特征：
	
	\subsubsection{岩石圈 (Lithosphere)}
	\begin{itemize}
		\item \textbf{定义：} 由地壳（Crust）和上地幔最上部（Uppermost Mantle）组成的坚硬外壳。
		\item \textbf{流变学特征：} 表现为刚性、脆性。其底部为\textbf{软流圈 (Asthenosphere)}，界面通常以温度（约 1300$^\circ$C）或地震波速度的显著降低（低速带, LVZ）为标志。
		\item \textbf{厚度：} 洋壳岩石圈较薄（约 5--100 km），随年龄增加而增厚；陆壳岩石圈较厚（可达 150--250 km）。
	\end{itemize}
	
	\subsubsection{洋中脊系统 (Ocean Ridge System)}
	\begin{itemize}
		\item \textbf{性质：} \textbf{离散型边界 (Divergent Boundary)}，海底扩张的核心区域。
		\item \textbf{地学特征：} 
		\begin{itemize}
			\item 产生新的大洋地壳（扩张中心）。
			\item 具有极高的热流值。
			\item 发育浅源地震和玄武岩质岩浆作用。
			\item 磁异常条带呈对称分布。
		\end{itemize}
	\end{itemize}
	
	\subsubsection{消减带 (Subduction Zone)}
	\begin{itemize}
		\item \textbf{性质：} \textbf{汇聚型边界 (Convergent Boundary)}，洋壳俯冲进入地幔的区域。
		\item \textbf{构造组合：} 包含海沟（Trench）、增生楔、弧前盆地、火山岛弧/边缘山脉。
		\item \textbf{地震学特征：} 发育\textbf{和达-贝尼奥夫带 (Wadati-Benioff Zone)}，地震震源深度可从地表延伸至 700 km（深源地震）。
	\end{itemize}
	
	\subsubsection{大陆碰撞带 (Continent Collision Zone)}
	\begin{itemize}
		\item \textbf{性质：} 两个大陆板块之间的汇聚边界（如喜马拉雅造山带）。
		\item \textbf{特征：} 
		\begin{itemize}
			\item 由于陆壳密度小、浮力大，难以俯冲。
			\item 导致剧烈的岩石圈缩短、地壳增厚（可达 70--80 km）。
			\item 发育大规模的逆冲断层、褶皱及高级变质作用。
		\end{itemize}
	\end{itemize}
	
	\subsubsection{转换断层 (Transform Fault)}
	\begin{itemize}
		\item \textbf{性质：} \textbf{守恒型边界 (Conservative Boundary)}，两板块相互水平滑动。
		\item \textbf{特征：} 断层两侧板块既不产生也不消亡。连接洋中脊段，或洋中脊与消减带。
		\item \textbf{经典实例：} 圣安德烈斯断层 (San Andreas Fault)。
	\end{itemize}
	
	\subsubsection{三联点 (Triple Junction Point)}
	\begin{itemize}
		\item \textbf{定义：} 三个板块边界的交汇点。
		\item \textbf{分类：} 根据边界类型（R-脊, T-沟, F-转换断层）组合，如 RRR 三联点。
		\item \textbf{动力学：} 并非所有三联点在几何上都是稳定的。
	\end{itemize}
	
	\subsubsection{热点 (Hotspot) 与地幔热柱 (Mantle Plume)}
	\begin{itemize}
		\item \textbf{热点：} 位于板块内部（非边界）的长期活跃火山活动中心（如夏威夷）。
		\item \textbf{地幔热柱：} 热点的物理根源，源自地幔深部（甚至核幔边界 D'' 层）上升的狭窄圆柱状热物质流。
		\item \textbf{意义：} 提供了板块运动的绝对参考系。
	\end{itemize}
	
	\subsubsection{地幔对流 (Mantle Convection)}
	\begin{itemize}
		\item \textbf{作用：} 板块构造的驱动力（引擎）。
		\item \textbf{机制：} 
		\begin{itemize}
			\item \textbf{板片拉力 (Slab Pull)：} 俯冲板片因冷却变重而在引力作用下向深部拖拽板块（主导力量）。
			\item \textbf{脊推力 (Ridge Push)：} 洋中脊的高势能向两侧推动板块。
			\item \textbf{全地幔对流与分层对流：} 现代地球物理倾向于认为存在穿过 660 km 界面的全地幔循环。
		\end{itemize}
	\end{itemize}
	\subsection{青藏高原的形成}
	
	青藏高原是世界上最年轻的一个高原。2.4亿年前，由于板块运动，印度板块以较快的速度开始向北向亚洲板块移动、挤压，其北部发生了强烈的褶皱断裂和抬升，促使昆仑山和青海可可西里隆升为陆地。
	
	到了距今8000万年前，印度板块继续向北漂移，再次引起了强烈的构造运动。随着印度板块不断向北推进，并不断向亚洲板块下插入，青藏高原在对此上升阶段中形成。
	
	青藏高原的形成并不是一次就完成的，其抬升过程不是一次性的猛增，也不是匀速的运动，而是经历了几个不同的上升阶段。
	
	\section{地震学和行星的内部结构}
	
	\subsection{用地震波透视行星内部}
	\subsubsection{地震波动力学基础}
	地震波是行星内部弹性应变能释放后产生的机械波，分为体波和面波。
	
	\textbf{体波 (Body Waves)}
	在行星内部三维空间传播的波。
	\begin{itemize}
		\item \textbf{P波 (Primary wave, 纵波)：} 质点振动方向与波传播方向平行。其速度 $v_p$ 为：
		\begin{equation}
			v_p = \sqrt{\frac{K + \frac{4}{3}\mu}{\rho}}
		\end{equation}
		其中 $K$ 为体积模量，$\mu$ 为剪切模量，$\rho$ 为密度。P波可在固体和液体中传播。
		
		\item \textbf{S波 (Secondary wave, 横波)：} 质点振动方向与波传播方向垂直。其速度 $v_s$ 为：
		\begin{equation}
			v_s = \sqrt{\frac{\mu}{\rho}}
		\end{equation}
		由于液体的 $\mu = 0$，\textbf{S波不能在液态外核中传播}。
		\begin{itemize}
			\item \textbf{SV分量：} 在垂直平面内振动的横波。
			\item \textbf{SH分量：} 在水平平面内振动的横波。
		\end{itemize}
	\end{itemize}
	
	\textbf{面波 (Surface Waves)}
	沿行星表面或分界面传播的波，振幅随深度增加而迅速衰减。在地球内部的传播速度大约在$ 2\sim5km/s $的范围.
	\begin{itemize}
		\item \textbf{勒夫波 (Love wave)：} 质点在水平面内做垂直于传播方向的振动（类似SH波）。比瑞利波速度快,
		在地球内部的传播速度大约在 $3\sim6km/s$的范围.
		\item \textbf{瑞利波 (Rayleigh wave)：} 质点在垂直平面内做逆向椭圆运动（地滚波）。
	\end{itemize}
	
	\subsubsection{地震仪与地震观测}
	地震仪利用惯性原理记录地面的运动。
	\begin{itemize}
		\item \textbf{三分量记录：} 现代地震观测站通常同时记录垂直（Z）、南北（N）和东西（E）三个方向的震动，用于重建质点运动轨迹。
		\item \textbf{宽频带观测：} 能够捕捉从高频（数赫兹）到超长周期（数千秒，如地球自由震荡）的完整信号。
	\end{itemize}
	
	\subsubsection{地震震相与内部结构}
	\textbf{震相 (Seismic Phases)} 是指在地震图上具有特征性振幅、频率和视速度的波段，代表了波经过的不同路径。
	
	\begin{table}[htbp]
		\centering
		\caption{常见地震震相及其物理意义}
		\begin{tabularx}{\textwidth}{l l X}
			\toprule
			震相符号 & 传播路径描述 & 结构指示意义 \\
			\midrule
			\textbf{P / S} & 直接穿过地幔的纵波/横波 & 探测地幔速度结构的基本手段。 \\
			\textbf{PP / SS} & 在地表反射一次后的纵波/横波 & 研究地幔中段的反射面。 \\
			\textbf{ScS} & 在核幔边界 (CMB) 反射的横波 & 测定地幔底部高度及 CMB 性质。 \\
			\textbf{PKP} & 进入外核并再次穿出的纵波 & 证实液态外核的存在（S波消失，P波折射）。 \\
			\textbf{PKIKP} & 穿过液态外核进入固态内核的纵波 & 探测内核的大小及各向异性。 \\
			\bottomrule
		\end{tabularx}
	\end{table}
	
	\subsubsection{地震层析成像 (Seismic Tomography)}
	地震层析成像利用海量的地震波走时数据（Travel-time），通过反演算法重建行星内部的三维波速分布，类似于医学中的 CT 扫描。
	
	\begin{itemize}
		\item \textbf{基本原理：} $\Delta t = \int \Delta (1/v) \, ds$。观测走时与参考模型走时的差值（残差）反映了路径上波速的异常。
		\item \textbf{物理应用：} 
		\begin{itemize}
			\item \textbf{蓝色高速度区：} 通常代表冷的、高密度的下沉板块（如俯冲带板片）。
			\item \textbf{红色低速度区：} 通常代表温热、低密度的上升物质（如地幔热柱）。
		\end{itemize}
	\end{itemize}
	\subsection{地震波的反射与折射}
	\subsubsection{地震波的传播规律：射线理论与 Snell's Law}
	
	\textbf{1. 地震波的反射与折射}
	当平衡状态下的地震波（如纵波 P）入射至两种不同速度/密度介质的交界面时，会遵循波动力学中的 \textbf{Snell 定律}。由于地球是固体介质，在界面上通常会发生\textbf{波型转换 (Mode Conversion)}，即入射 P 波可以激发出反射 P 波、反射 S 波、折射 P 波以及折射 S 波。
	
	\textbf{Snell's Law 数学表达：}
	\begin{equation}
		\frac{\sin i_1}{V_1} = \frac{\sin i'_2}{V_2} = \frac{\sin r_3}{V_3} = p
	\end{equation}
	
	\noindent \textbf{参数定义：}
	\begin{itemize}
		\item $i_1, i'_2, r_3$：分别代表入射角、反射角和折射角。
		\item $V_1, V_2, V_3$：代表相应波在各层中的传播速度（P 波速 $\alpha$ 或 S 波速 $\beta$）。
		\item $p$：称为\textbf{射线参数 (Ray Parameter)}，对于同一条射线，该值在整个传播路径中保持不变。
	\end{itemize}
	
	\includegraphics[width=0.4\textwidth]{fanshen.png} 
	\includegraphics[width=0.5\textwidth]{4.png} 
	\textbf{2. 分层均匀介质中的传播特征}
	在分层均匀介质模型中（水平层状模型），地震射线经过每一个界面时都遵循波速与角度的正弦比例关系。
	
	通用射线方程：
	\begin{equation}
		\frac{\sin i_1}{V_1} = \frac{\sin i_2}{V_2} = \dots = \frac{\sin i_n}{V_n} = p = \text{Constant}
	\end{equation}
	
	\noindent \textbf{物理推论：}
	\begin{enumerate}
		\item \textbf{波速与角度关系：} 随着深度 $Z$ 的增加，如果地层速度 $V$ 增大（$V_1 < V_2 < \dots < V_n$），则折射角 $i$ 也会相应增大。这意味着地震射线会不断向远离法线的方向偏折。
		\item \textbf{射线路径：} 在速度随深度连续增加的行星内部，地震射线的轨迹通常表现为向下弯曲的弧线。
		\item \textbf{临界角：} 当折射角 $i_n = 90^\circ$ 时，正弦值为 1，射线参数 $p = 1/V_n$。此时波沿界面传播，称为临界折射（首波形成的前提）。
	\end{enumerate}
	
	\textbf{3. 专业术语汇总}
	\begin{table}[htbp]
		\centering
		\begin{tabular}{ll}
			\toprule
			术语 & 物理/地学含义 \\
			\midrule
			Ray Path & \textbf{地震射线}：垂直于等时面的波前传播方向。 \\
			Snell's Law & \textbf{斯内尔定律}：描述跨越界面时速度与角度的守恒关系。 \\
			Ray Parameter ($p$) & \textbf{射线参数}：表征单条射线几何特征的物理量，单位通常为 s/km。 \\
			Velocity Gradient & \textbf{速度梯度}：介质速度随深度变化的速率，决定射线的曲率。 \\
			\bottomrule
		\end{tabular}
	\end{table}
	\subsection{液态地核和固态内核}
	1906年，英国地震学家R. D. Oldham发现地核是液态的.
	
	1913年，美国地震学家B. Gutenberg根据地震资料计算出核－幔边界的深度为2900 km。
	
	1936年，丹麦地震学家Inge Lehmann(1888-1993)发现了固态内核，以及后来的220 km间断面。
	
	\textbf{地球内核存在东西半球差异}和西半球(180E－40E)相比，东半球(40E –180E) 地震波波速快($\sim1\%$)，衰减强。
	
	液核的存在使得地球表面存在接收不到地震波的地方—地震波影区.
	
	\includegraphics[width=0.4\textwidth]{6.png}
	\includegraphics[width=0.5\textwidth]{5.png}
	
	\subsection{地球的圈层结构（PREM模型）}
	
	\includegraphics[width=0.5\textwidth]{7.png}
	\includegraphics[width=0.5\textwidth]{9.png}
	
	\includegraphics[width=0.9\textwidth]{8.png}
	
	\section{磁场、重力与行星探测}
	
	\subsection{地球的磁场特点}
	
	1)地球磁场的主要部分相当于一巨大磁铁(磁偶极子)
	
	2)地球磁场的北极指向地理南极但存在约11.5度的偏角
	
	3)地磁偶极子场：赤道0.307 高斯, 两极0.6高斯；磁矩 $7.94X10^{23}Gs\text{·}cm3$
	
	4)地球磁场除基本磁场外还有一变化磁场.
	
	5)地球磁场在空间可以延伸到很远的空间.
	
	6)地球磁场分来自内部的内源场和来自空间的外源场.
	
	\subsection{地磁三要素}
	\textbf{磁偏角、磁倾角、磁场强度}
	\begin{itemize}

		\item 关系式：
		\[
		F^2 = X^2 + Y^2 + Z^2, \quad H^2 = X^2 + Y^2
		\]
		\[
		Y = H \sin D, \quad X = H \cos D, \quad Z = H \tan I
		\]
		\item 地磁场包括来自地球内部的“主磁场”和来自太空的“外源场”。
	\end{itemize}
	
	\subsubsection{地磁场测量与特征}
	\begin{itemize}
		\item 测量方法包括陆地观测、空中观测、卫星观测和古地磁测量。
		\item 古地磁测量需在“零磁空间”中进行，使用旋转磁力仪、超导磁力仪等设备。
		\item 全球磁场分布：赤道磁场强度约0.307高斯，两极约0.6高斯。
		\item 地磁偶极矩：\(7.94 \times 10^{25}\)高斯·厘米\(^3\)。
		\item 地磁极随时间移动，历史上发生过多次磁极倒转。
	\end{itemize}
	
	\subsubsection{地球磁场的起源——地球发电机理论}
	\begin{itemize}
		\item 地磁场源于外核液态铁镍流体的对流运动（地球发电机效应）。
		\item 内核形成于地球早期，外核为液态，内核为固态铁镍合金。
		\item 地磁场保护地球大气层和生命免受太阳风等宇宙辐射伤害。
		\item 磁极倒转过程模型显示，内核磁化方向与主磁场相反，倒转时间尺度受外核磁场向内核扩散控制。
	\end{itemize}
	
	\subsubsection{磁异常与古地磁学应用}
	\begin{itemize}
		\item 磁异常：测量总磁场强度后减去主磁场强度，得到异常值。
		\item 感应磁化：材料在外磁场 \(H\) 中产生磁化强度 \(I = kH\)，\(k\) 为磁化率。
		\item 正异常：异常体磁场加强地球磁场；负异常：异常体磁场减弱地球磁场。
		\item 海底磁异常条带为板块构造和海底扩张提供了关键证据。
		\item 古地磁学可用于重建板块运动历史（如150百万年前板块位置重建）。
	\end{itemize}
	
	\subsubsection{磁场几何关系}
	\begin{itemize}
		\item 磁倾角与纬度关系：\(\tan I = 2 \cot \theta_m = 2 \tan \phi_m\)。
		\item 磁场分量：
		\[
		B_H = B \cos I, \quad B_V = B \sin I
		\]
		\[
		B_{HN} = B \cos I \cos D, \quad B_{HE} = B \cos I \sin D
		\]
		\item 偶极子磁场公式：
		\[
		B_r = \frac{\mu m}{2\pi r^3} \cos \theta_m, \quad B_\theta = \frac{\mu m}{4\pi r^3} \sin \theta_m
		\]
	\end{itemize}
	
	\includegraphics[width=0.4\textwidth]{10.png}
	\includegraphics[width=0.5\textwidth]{11.png}
	\subsection{地球的重力、均衡与潮汐}
	
	\subsubsection{重力场基本概念}
	\begin{itemize}
		\item 重力 \(g\) 是引力 \(f\) 与离心力 \(q\) 的矢量和：
		\[
		g(r_p) = f(r_p) + q(r_p)
		\]
		\item 重力位 \(W\) 是引力位 \(V\) 与离心力位 \(Q\) 之和：
		\[
		W(x,y,z) = V(x,y,z) + Q(x,y,z) = G \int \frac{dm}{r} + \frac{\omega^2}{2}(x^2 + y^2)
		\]
		\item 重力单位：
		\begin{itemize}
			\item SI制：m/s\(^2\)，常用单位：g.u. (gravity unit)，1 g.u. = \(10^{-6}\) m/s\(^2\)
			\item CGS制：伽（gal），1 gal = 1 cm/s\(^2\) = \(10^3\) mgal = \(10^8\) μgal
		\end{itemize}
	\end{itemize}
	
	\subsubsection{大地水准面与重力异常}
	\begin{itemize}
		\item 重力等位面（水准面）：重力位为常数的曲面，与重力矢量正交。
		\item 大地水准面：平均海平面延伸至陆地的重力等位面，是海拔高程的基准。
		\item 扰动重力场：实际重力场与参考椭球体重力场之差。
		\item 重力异常 \(\Delta g\) 的计算（以球体为例）：
		\[
		\Delta g = \frac{4\pi G R^3 \Delta \rho}{3} \cdot \frac{b}{(x^2 + b^2)^{3/2}}
		\]
	\end{itemize}
	
	\subsubsection{重力均衡}
	\begin{itemize}
		\item 阿基米德原理：浮力等于排开流体的重量。
		\item 地壳漂浮在地幔上，遵循均衡原理。
		\item 艾黎（Airy）均衡模型：
		\begin{itemize}
			\item 陆地：山越高，山根越深
			\item 海洋：水越深，反山根越高
		\end{itemize}
		\item 山根厚度计算（以5 km高山为例，地壳密度2670 kg/m\(^3\)，地幔密度3300 kg/m\(^3\)）：
		\[
		t = h \cdot \frac{\rho_c}{\rho_m - \rho_c}
		\]
	\end{itemize}
	\subsubsection{Airy 均衡模型 (Airy Isostasy)}
	
	Airy 均衡模型基于“浮力原理”。它假设地壳由密度均匀的块体组成，通过改变厚度（形成山根）在密度较大的地幔中达到平衡。
	
	\textbf{图解说明}
	
	\begin{figure}[h]
		\centering
		\begin{tikzpicture}[scale=0.8]
			% 绘制地幔
			\fill[orange!20] (-2,-5) rectangle (8,-1);
			\node at (7,-0.7) {地幔 $\rho_m$};
			
			% 绘制标准地壳
			\fill[gray!30] (-2,-3) rectangle (0,0);
			\draw (-2,0) -- (0,0);
			\draw (-2,-3) -- (0,-3);
			\node at (-1,-1.5) {标准地壳};
			\draw[<->] (-1.8,0) -- node[left] {$T$} (-1.8,-3);
			
			% 绘制山脉地壳
			\fill[gray!50] (1,-4.5) rectangle (5,1.5);
			\draw (1,1.5) -- (5,1.5);
			\draw (1,-4.5) -- (5,-4.5);
			\node at (3,-1.5) {山脉地壳 $\rho_c$};
			
			% 标注 H, T, r
			\draw[<->] (5.2,0) -- node[right] {$H$ (海拔)} (5.2,1.5);
			\draw[<->] (5.2,-3) -- node[right] {$T$ (标准厚度)} (5.2,0);
			\draw[<->] (5.2,-4.5) -- node[right] {$r$ (山根)} (5.2,-3);
			
			% 海平面和补偿面
			\draw[dashed, blue] (-3,0) -- (8,0) node[right] {海平面};
			\draw[dashed, red, thick] (-3,-4.5) -- (8,-4.5) node[right] {补偿面 (Depth of Compensation)};
			
			% 标注密度
			\node at (3,0.7) {$\rho_c$};
		\end{tikzpicture}
		\caption{Airy 均衡模型示意图}
	\end{figure}
	
	\textbf{公式推导}
	
	在红色虚线所示的\textbf{补偿面}处，静水压力必须相等。
	
	设标准地壳下方的压力为 $P_0$，山脉下方的压力为 $P_1$：
	

		$$P_0 = \rho_c g T + \rho_m g r$$

	
	
	$$	P_1 = \rho_c g (H + T + r)$$

	
	根据均衡假设 $P_0 = P_1$，得：
	
	$$	\rho_c g T + \rho_m g r = \rho_c g H + \rho_c g T + \rho_c g r$$

	
	约去共同项 $g$ 和 $\rho_c T$：
		$$\rho_m r = \rho_c H + \rho_c r$$
	
	解出山根深度 $r$：
		$$\mathbf{r = \frac{\rho_c}{\rho_m - \rho_c} H}$$
	\textbf{物理意义}
	\begin{itemize}
		\item \textbf{补偿机制}：地表正向的地形起伏（山脉）由地底负向的密度亏损（山根取代了地幔）来补偿。
		\item \textbf{海洋情况}：对于海洋，地表是负地形（深度为 $d$），则下方会有“反山根”（Anti-root），即高密度的地幔会上涌。
		\item \textbf{应用}：该理论解释了为什么喜马拉雅山脉下方的莫霍面（Moho）深度比平原地区深得多。
	\end{itemize}
	
	\subsubsection{地球转动与潮汐}
	\begin{itemize}
		\item 地球自转与公转导致岁差、章动、极移、日长变化等现象。
		\item 潮汐：日月引力引起地球形变，包括海潮、固体潮、大气潮和磁层潮。
		\item 起潮力计算：地球与月球绕共同质心运动产生的离心力与引力的合力。
		\item 起潮力公式（以月球为例）：
		\[
		t(A) = \frac{2Gm_m R}{r_m^3} \approx 104.7 \, \text{μgal}
		\]
		\end{itemize}
	\subsubsection{潮汐摩擦}
	如果地球是完全弹性体，那么潮汐椭园的长轴将一直在天体和地球的连心线上。但是观测表明对于月－地系统，潮汐椭园迟后月—地联心线5度左右。产生这一现象的原因可以用地球非完全弹性，或具有一定的粘滞性来解释。而这一现象的后果是，由于潮汐椭园的滞后，月球作用于椭球凸起部分的引力合将迫使地球自转速度不断减慢。这一现象称之为潮汐摩擦.
	\includegraphics[width=0.8\textwidth]{12.png}
	
	然而，由于当我们将月－地视为一个系统时，其总的角动量守恒。因此，月球离开地球的距离也会相应增加。不过使得地球自转减慢主要的因素是海潮运动相应的潮汐摩擦，其作用远远超出了固体潮的影响。有估计显示，由于月球产生的潮汐磨擦地球转速变慢，大约在9亿年前地球 18小时/天，那时月球比现在更接近地球。现代月球后退速率 4cm/年，如果太阳不崩溃的话，此关系将维持 150 亿年。
	\subsubsection{重力测量技术}
	\begin{itemize}
		\item 地面重力仪：弹簧重力仪、超导重力仪、原子重力仪等。
		\item 卫星重力测量：如CHAMP、GRACE、GOCE等任务，提供全球重力场模型。
		\item 海洋重力测量：卫星测高（如TOPEX/Poseidon）反演海洋重力异常。
	\end{itemize}
	
	\subsection{行星探测}
	
	\subsubsection{行星重力场}
	\begin{itemize}
		\item 行星重力场是研究行星内部结构的重要手段。
		\item 类地行星（火星、金星、月球）重力场可通过轨道器探测获得。
		\item 重力异常揭示行星内部质量分布、地壳厚度、均衡状态等。
	\end{itemize}
	
	\subsubsection{行星磁场}
	\begin{itemize}
		\item 行星磁场强度与行星内部动力学状态相关。
		\item 地球、木星、土星等具有强磁场；火星、金星磁场较弱。
		\item 行星磁场保护其大气层，并影响空间环境。
	\end{itemize}
	
	\subsubsection{行星地震探测}
	\begin{itemize}
		\item 通过地震波研究行星内部结构（如月球、火星）。
		\item 地震数据可推断行星内部层圈结构、物质组成与物理状态。
	\end{itemize}
	
	\section{行星大气环境与深空探测}
	
	\subsection{太阳系与地球}
	
	八大行星按照离太阳的距离从近到远，依次为水星、金星、地球、火星、木星、土星、天王星、海王星。
	
	八大行星自转方向多数和公转方向一致。只有金星和天王星两个例外。金星自转方向与公转方向相反。
	
	在2006年8月24日冥王星被划为矮行星，从太阳系九大行星中被除名。\\
	\includegraphics[width=1.0\textwidth]{13.png}
	
	\subsubsection{类地行星}
	
	类地行星是以硅酸盐石作为主要成分的行星。它们跟类木行星有很大的分别，因为那些气体行星主要是由氢、氦和水等组成，而不一定有固体的表面。
	
	类地行星的结构大致相同：一个主要是铁的金属中心，外层则被硅酸盐地幔所包围。它们的表面一般都有峡谷、陨石坑、山和火山。
	
	类地行星的大气层都是再生大气层，有别于类木行星直接来自于太阳星云的原生大气层。
	
	\subsubsection{地球有什么特别之处？}
	
	恰当的位置, 恰当的日地距离 (1个天文单位), 恰当的大小和质量.
	
	\subsubsection{地球磁场}
	
	地球磁场属于电磁场，是通过外核的电子随地球自转的电流效应 (即电生磁) 产生的磁场。
	
	地球磁场强度大约是0.5高斯，也就是0.00005特斯拉（$50\sim60 \mu T$）。
	
	\subsubsection{气辉 (airglow)}
	地球高层大气吸收了太阳电磁辐射产生的微弱光辐射。
	
	20世纪初，天文学家在研究夜晚天空的亮度和总的星光时，确认天空的亮度除来源于星光外，还有来源于地球大气的部分，当时称作地球光。这种光中总有一条在极光光谱中出现的\underline{绿色光谱}，故又称永极光或非极区极光。直至1950年才被称作气辉。
	
	\subsubsection{磁层被破坏后会引起极光}
	
	大多数极光出现于地球上空90-130公里处，宽度可达数千公里。但有些极光要高得多，高度可达560-1000公里以上。地球的极光是来自地球磁层或太阳的高能带电粒子流 (太阳风) 使高层大气分子或原子激发 (或电离) 而产生。
	
	\subsection{太阳辐射常数}
	太阳常数是指在日地平均距离 (D=15000万公里) 上，大气顶界垂直于太阳光线的单位面积每秒钟接受的太阳辐射。太阳常数要在地球大气层之外，垂直于入射光的平面上测量。
	
	因为太阳表面常有黑子等太阳活动，太阳常数并不是固定不变的，一年中的变化幅度在1\%左右，且呈现太阳活动的11年变化周期。
	
	\subsubsection{地球大气层有多厚？}
	
	大气层是地球最外部的气体圈层，包围着海洋和陆地，大气层的厚度大约在1000千米以上，但没有明显的界限。
	
	这一层气体主要由氮气、氧气以及少量的其他气体 (如二氧化碳、氩气、水蒸气等) 组成，对地球生物的生存和地球气候的形成起着至关重要的作用。
	
	地球大气层没有明显的界限，其范围十分广阔，理论上可延展至数千公里，但气体密度随高度增加而急剧减小，大部分质量集中在离地表约30千米以下，98.2\%的质量集中在30千米以下。
	
	\subsubsection{太阳辐射能量分布和可见光谱}
	
	太阳辐射能量分布覆盖紫外区 ($<0.4 \mu m$)、可见光区($0.4 \sim 0.76 \mu m$)和红外区 ($>0.76 \mu m$)，其中可见光区占总能量的50\%，是人眼可以感知的部分.
	
	太阳光谱能量分布的最大能量出现在波长约0.475微米(475纳米) 的蓝绿光区域。此外，红外光谱约占43\%的
	能量，而紫外光谱约占7\%.
	
	\subsubsection{地球大气窗区}
	大气窗区指电磁波穿透地球大气的特定波段，按波长可分为可见光窗区、红外窗区和射电窗区。大气窗区是发展遥感探测与通信技术的技术基础。
	
	\subsection{黑体辐射三大定律}
	
	黑体辐射是量子力学诞生的摇篮，描述了处于热平衡状态的黑体所发射出的电磁辐射特性。其核心由以下三大定律构成：
	
	\subsubsection{普朗克黑体辐射定律 (Planck's Law)}
	普朗克定律描述了在任意温度 $T$ 下，从一个黑体中发射出的电磁辐射的辐射率与频率（或波长）之间的关系。它是量子物理的基础。
	
	\begin{itemize}
		\item \textbf{频率形式：} 光谱辐射亮度 $B_\nu(\nu, T)$ 表示为：
		\begin{equation}
			B_\nu(\nu, T) = \frac{2h\nu^3}{c^2} \frac{1}{e^{\frac{h\nu}{k_B T}} - 1}
		\end{equation}
		
		\item \textbf{波长形式：} 能量密度频谱也可写成波长的函数 $u(\lambda, T)$：
		\begin{equation}
			u(\lambda, T) = \frac{8\pi hc}{\lambda^5} \frac{1}{e^{\frac{hc}{\lambda kT}} - 1}
		\end{equation}
	\end{itemize}
	其中：$h$ 为普朗克常数，$c$ 为光速，$k_B$ 为玻尔兹曼常数，$T$ 为热力学温度。\\
	\includegraphics[width=0.4\textwidth]{14.png}
	\includegraphics[width=0.5\textwidth]{15.png}
	
	\subsubsection{维恩位移定律 (Wien's Displacement Law)}
	维恩位移定律描述了黑体电磁辐射光谱辐射度的\textbf{峰值波长}与自身\textbf{温度}之间成反比关系的规律。
	
	\begin{itemize}
		\item \textbf{公式：}
		\begin{equation}
			\lambda_{\text{max}} = \frac{b}{T}
		\end{equation}
		\item \textbf{维恩位移常数：} $b \approx 2.8977729 \times 10^{-3} \, \text{m} \cdot \text{K}$。
		\item \textbf{物理意义：} 随着温度升高，辐射能量的峰值向短波方向（高频方向）移动。例如，物体加热时颜色从暗红变为橙黄，再变为青白。
	\end{itemize}
	
	\subsubsection{斯特藩-玻尔兹曼定律 (Stefan-Boltzmann Law)}
	该定律描述了黑体表面单位面积在单位时间内辐射出的总功率（称为\textbf{辐照度}或\textbf{能量通量密度} $E$ 或 $j^*$）与黑体本身的热力学温度 $T$ 的四次方成正比。
	
	\begin{itemize}
		\item \textbf{公式：}
		\begin{equation}
			E = \sigma T^4
		\end{equation}
		\item \textbf{比例系数 $\sigma$：} 称为斯特藩-玻尔兹曼常数（Stefan-Boltzmann constant）。
		\item \textbf{物理意义：} 总辐射能量对温度极其敏感。当温度翻倍时，辐射出的能量将增加到原来的 16 倍。
	\end{itemize}
	
	\subsubsection{总结与联系}
	\begin{itemize}
		\item \textbf{普朗克定律}是普适公式，维恩位移定律和斯特藩-玻尔兹曼定律均可以从普朗克定律推导得出。
		\item 维恩定律对应于普朗克公式极值点的轨迹。
		\item 斯特藩-玻尔兹曼定律对应于普朗克公式在全波段上的积分（即光谱曲线下方的总面积）。
	\end{itemize}
	
	\subsection{地球的大气环境}
	
	地球大气环境以氮气和氧气为主，包含少量其他气体的混合物，其通过吸收紫外线、保护地表免受小行星撞击、调节气候以及提供生命必需的氧气来维持地球生态系统的平衡和生命的存在。
	
	地球大气是指地球外围因为由于地球引力作用而环绕地球的空气层，是地球自然环境的重要组成部分之一，与人类的生存息息相关。
	
	通 常 把 从 地 面 到1000-1400千米高度内的气层作为地球大气层的厚度，其中大气总质量的98.2\%集中在30千米以下。
	
	\textbf{主要成分} 为氮、氧、氩、二氧化碳、水等，组成比率因时、地不同而有所差异，其中以二氧化碳变动率最大。
	
	\textbf{大气密度} 不是均匀的，以海平面的密度最大，往上密度渐小，大气约50\%集中在海拔5.6公里内。
	
	\subsubsection{地球大气的卡门线}
	
	一般认为，在海拔100公里的位置，就是地球大气层与外太空的分界线。在这个高度，气体分子已经非常稀少，大约每立方米才100至120个左右的气体分子。
	
	海拔100公里的交界线也被称为“卡门线”，这是国际航空联合会所接受的现行大气层和外太空的界线。这个概念由匈牙利裔美国工程师西奥多·冯·卡门提出。普通的航空器一旦超过这个“卡门线”，燃油发动机将因为氧气分子过少而无法启动。
	
	\subsubsection{地球大气的分层结构}	
	
	大气层又称大气圈，是因重力关系而围绕地球的一层混合气体，是地球最外部的气体圈层，包围着海洋和陆地。从下到上大气层可以分为五层：对流层、平流层、中间层、热层和散逸层。平时说的臭氧层是平流层中臭氧浓度高的部分。
	\subsubsection{地球大气温度和密度等的高度变化}
	
	地球大气受引力束缚聚集在地球周围。随高度增加，大气的温度、密度、压力及声速均呈现出复杂的分布规律。
	
	\begin{enumerate}
		\item \textbf{大气质量与密度的分布规律}
		\begin{itemize}
			\item \textbf{高度集中性：} 地球大气密度随高度增加而显著减小。约 90\% 的大气质量集中在海平面以上 16 km 以内的低层大气中。
			\item \textbf{极端稀薄性：} 随高度升高，大气变得极其稀薄。在海拔 240 km 处，大气密度仅为海平面的百万分之一 ($10^{-6}$)。到数千公里高空时，其密度已接近星际空间。
			\item \textbf{影响因素：} 空气密度同时受海拔和气温影响。海拔升高或温度升高均会导致空气密度降低。
		\end{itemize}
		
		\item \textbf{大气层级结构与温度特性 (Atmospheric Layering)}
		依据温度随高度的变化趋势（气温垂直递减率），大气可分为以下主要层次：
		\begin{itemize}
			\item \textbf{对流层 (Troposphere, 0--11 km)：} 温度随高度增加而降低。人类活动及天气现象主要发生在此层。顶部为对流层顶 (Tropopause)。
			\item \textbf{平流层 (Stratosphere, 11--50 km)：} 由于臭氧层 (Ozone Layer) 吸收紫外线，温度随高度增加而上升。顶部为平流层顶 (Stratopause)。
			\item \textbf{中间层 (Mesosphere, 50--85 km)：} 温度随高度增加而再次下降，是大气中最冷的区域。顶部为中间层顶 (Mesopause)。
			\item \textbf{热层 (Thermosphere, 85--600 km)：} 受太阳高能辐射影响，温度随高度剧烈上升（中性温度可达 1000K 以上）。此层包含电离层 D 层 (Ionosphere D layer) 和极光 (Aurora) 现象。
			\item \textbf{散逸层 (Exosphere, > 600 km)：} 大气的最外层，气体分子极少，逐渐过渡到外层空间。
		\end{itemize}
		
		\item \textbf{其他关键物理参数}
		\begin{itemize}
			\item \textbf{气压 (Pressure)：} 随高度增加呈指数级下降。
			\item \textbf{声速 (Speed of Sound)：} 声速的变化规律与温度曲线高度同步。在对流层随高度下降，在平流层随高度上升。
		\end{itemize}
	\end{enumerate}
	
	\subsubsection{大气层分层结构的热力学特征}
	
	大气层按热力学特征主要分为对流层、平流层、中间层、热层和散逸层五层，各层的主要特征包括温度随高度的变化趋势、大气扰动程度以及与太阳辐射和臭氧等因素的关系。
	
	\subsubsection{地球大气层的温室效应}
	来自太阳的可见光、不可见的紫外线和红外线，可以穿透覆盖地球的大气层。大约70\%的太阳光被地球上的海洋、陆地和大气吸收，剩下的30\%会被反射回太空。
	
	温室效应，又称“花房效应”，是大气保温效应的俗称。大气中的一些气体能对气候造成一定程度的影响，这些气体擅长吸收长波辐射，因其作用类似于栽培农作物的温室，故名温室效应。
	\subsubsection{温室气体增加导致的全球变暖}
	全球气候变暖是一种和自然有关的现象，是由于温室效应不断积累，导致地气系统吸收与发射的能量不平衡，能量不断在地气系统累积，从而导致温度上升，造成全球气候变暖.
	\subsection{臭氧层}
	
	臭氧层是大气层的平流层中臭氧浓度高的层次。 浓度最大的部分位于20-25公里的高度处。
	
	臭氧含量随纬度、季节和天气等变化而不同。 紫外辐射在高空被臭氧吸收，对大气有增温作用，同时保护了地球上的生物免受远紫外辐射的伤害，透过的少量紫外辐射，有杀菌作用，对生物大有裨益。
	
	\subsubsection{地球大气臭氧层的作用}
	
	\textbf{保护作用} 臭氧层能够吸收太阳光中波长306.3微米以下的紫外线（部分UV-B和全部UV-C），保护地球上的人类和动植物免遭短波紫外线的伤害。因为臭氧层能让太阳光中的可见光通过，且能吸收掉99\%以上的有害紫外辐射，所以有人称臭氧层为地球生命的保护神、保护伞或臭氧屏障。
	
	\textbf{加热作用}臭氧吸收太阳光中的紫外线并将其转换为热能加热大气，因此大气温度结构在高度50千米左右存在一个峰值。正是由于臭氧才有平流层的存在。
	
	地球以外的星球因不存在臭氧和氧气，所以也就不存在平流层。
	
	\subsubsection{地球臭氧的空间分布特征}
	
	大气层的臭氧主要是通过氧原子和氧气结合形成的，而氧原子主要来自紫外线照射氧气的分解。因为紫外线会一直照射
	大气层中的氧气，所以在大气层中臭氧是可以不断生成的。
	
	距地面15~50千米高度的大气平流层，集中了地球上约90\%的臭氧，这就是通常所说的“臭氧层”。其主要作用是吸收短波紫外线，是包围在地球外围空间的保护层。
	
	\subsection{地球的海洋}
	
	地球上海洋总面积约为3.6亿平方千米，约占地球表面积的71\%，平均水深约3795米。
	
	\subsubsection{不同视角看固体地球}
	\textbf{地球的真实形状-不规则的椭圆形}
	
	经过测量，发现地球不是规则的圆形，两极收缩变平，中间隆起变胖。这只是个大概的样子。地球上的水也只是说像地球的软壳一样，地球更不规则。大面积隆起和大面积凹陷，像发育畸形的山核桃。
	
	\subsubsection{地球上的水和水资源}
	
	地球上的总水量有13.86亿立方千米,这些水以液态、固态和气态三种形式存在于地球表面、地下及大气中。
	
	由于海水难以直接利用，因而所说的水资源主要指陆地上的淡水资源 (大约约3500万立方千米)。
	
	\subsubsection{全球海洋观测系统 (Global Ocean Observing System)}
	
	全球海洋观测系统，由4个国际机构发起并组织实施：政府间海洋学委员会 (IOC)、世界气象组织、国际科学联合会理事会和联合国环境规划署。致力于获得与分发有关海洋环境现状与未来状态的可靠评估和预报资料，以便有效、安全和持续利用海洋环境；为气候变化预报做出贡献，以便使广大用户获益；为海洋科学各学科的研究、开发和培训指明方向。海洋观测由4个方面组成，即志愿观测船 (VOS)，海洋高空探测 (ASAP)，漂浮浮标和系留浮标。
	
	\subsection{地球系统中的主要循环过程}
	
	地球系统的循环过程主要包括水循环、碳循环和营养元素 (生物地球化学) 循环，它们描述了水、碳和营养元素等物质在地球大气圈、水圈、岩石圈和生物圈之间的持续运动和转化。这些循环对维持地球生态系统的平衡和生命活动至关重要。
	
	此外，岩石循环也是地球系统中一个重要的长期地质过程。
	
	\subsubsection{地球大气层中的水循环}
	
	水循环是指地球上不同的地方上的水，通过吸收太阳的能量，改变状态到地球上另外一个地方。
	
	形成水循环的外因是太阳辐射和重力作用，其为水循环提供了水的物理状态变化和运动能量；形成水循环的内因是水在通常环境条件下气态、液态、固态三种形态容易相互转化的特性。
	
	\subsubsection{地球系统中的碳循环}
	
	地球系统的碳循环是指碳元素通过生物圈、大气圈、水圈和岩石圈之间的相互作用，不断进行储存和释放的动态过程，主要包括植物的光合作用固碳、生物体的呼吸作用和分解作用释放碳、以及化石燃料的燃烧等过程，这些过程维持了地球碳平衡并对气候变化至关重要。
	
	地球表层的碳库总储量约为40多万亿吨(4×104 Gt)，其中38万亿吨在海洋里。这
	些是在人类时间尺度上进行碳循环的组成部分。
	
	\subsubsection{地球系统中的氮循环}
	
	地球系统氮循环是在大气圈、生物圈、水圈和地壳中，氮元素以不同化学形式运动和转化的自然过程，主要由氮固定、硝化作用、硝酸盐转化和反硝化作用等环节驱动。氮循环为生命提供必需的反应性氮，而人类活动大幅增加了反应性氮，导致水体富营养化和温室气体排放等环境问题，对地球生态平衡造成严重影响。
	
	\subsection{其他行星的大气和海洋}
	
	太阳系中有大气层的行星包括金星、地球、火星、木星、土星、天王星和海王星。
	
	海洋在太阳系中非常罕见，地球是唯一拥有大范围液态水的行星，而木星和土星拥有液态核心，天王星和海王星则拥
	有巨大的水冰海洋，但水并非以液态水的形式存在。
	
	\subsubsection{水星 (Mercury)}
	
	太阳系八大行星中最小的行星, 基本没有大气, 由大约70\%的金属和30\%的硅酸盐组成.
	
	\textbf{水星凌日 (transit of Mercury)}
	当水星走到太阳和地球之间时，在太阳圆面上会看到一个小黑点穿过，这种现象称为水星凌日。其道理和日食类似，不同的是水星比月亮离地球远，视直径仅为太阳的190万分之一。
	
	水星挡住太阳的面积太小，不足以使太阳亮度减弱。所以，肉眼是看不到水星凌日的，只能通过望远镜进行投影观测。
	
	\subsubsection{金星 (Venus)}
	金星在夜空中的亮度仅次于月球，又称启明星、太白金星.
	
	金星的自转周期 (243天) 比金星绕太阳公转周期 (224.7天) 长.
	
	金星在过去可能拥有海洋，但是随着失控的温室效应导致温度上升而全部蒸发掉，是太阳系中最热的行星.
	
	2020年9月15日，科学家在金星大气层中侦测到磷化氢存在，这可能是地外生命存在的迹象.
	
	\textbf{金星大气层} 比地球大气层更为厚重与浓密，其表面温度较高，而气压则为93大气压，主要为二氧化碳所构成。金星的大气层中有硫酸形成的不透明云，因此在地球或金星环绕探测器上不可能以可见光观测金星表面。
	
	\textbf{金星表面的地形} 是以雷达成像的方式探测得知。金星大气层主要由二氧化碳和氮组成，以及少许痕量气体。
	
	\textbf{金星上的神秘暗斑} 不透明的高反射云永久地覆盖了这星球。金星的上层云很奇怪的是，那些云中有神秘的暗斑，被称为”未知吸收物”，它们构成微小颗粒吸收大部分的紫外线和来自太阳的一些可见光，从而影响金星的反照率。
	
	2007年底，欧洲航天局的“金星快车”号探测器首次记录了金星上的独特闪电。这是除地球之外，银河系中第三个被发现有闪电的行星，前两个分别为土星和木星。但是科学家也证实，\underline{金星上的闪电是已知的唯一一种与水云无关的闪电}，金星的闪电与硫酸云有关。
	
	\textbf{金星地表特征} 金星大部分表面相对平坦，主要分为三种地形单元：低地、高地和平原。
	
	金星被一层厚厚的云层所覆盖，阻挡了对行星表面的观测。然而，根据一项最新研究，这些云或可能研究云层下方炙热昏暗的金星表面。科学家们发现云层下方的气候模式会因地表的地形样式而发生改变，观察这些云层或可帮助理解金星表面。
	
	\subsubsection{火星 (Mars)}
	
	火星很有变成殖民地的价值，证明地球不只是宇宙中唯一生命生存地，人类也不是孤独的。
	
	1)火星和太阳的距离和地球与太阳距离相差不大（即温度和辐射差别不大）.
	
	2)火星存在大气层，即可能发展成像地球一样能避免紫外线伤害和陨石伤害的星球.
	
	3)火星可能存在水资源，而在有液态水的地方，是有可能存在着生命的（尽管可能存在的只是像单细胞之类的微生物）.
	
	\textbf{火星地表特征} 火星地表沙丘、砾石遍布，没有稳定的液态水体。二氧化碳为主的大气既稀薄又寒冷，沙尘悬浮其中，每年常有沙尘暴 (持续可达数天) 发生。火星上有太阳系最高的山：奥林帕斯山。火星上有太阳系最大的峡谷：水手号峡谷。2013年12月公布的探测资料表明，火星上存在已干涸的远古湖泊，但无法证明古代曾出现过流水。
	
	火星奥林帕斯山是太阳系最高的山.
	
	\textbf{火星大气层} 按高度划分为低层大气、中层大气、上层大气和外气层等层结，其主要成分是二氧化碳，相对地球稀薄得多，平均气压仅为地球的0.6\%。
	
	火星大气层受季节变化影响显著，冬季时大气中大量二氧化碳会凝华成固态干冰并沉积在极冠，夏季则会升华回大气，这种季节性的干冰沉积和升华过程会运输大量的尘埃和水蒸气。
	
	\subsection{太阳系外行星 (Exoplanet)}
	
	太阳系外行星，或简称系外行星，是位于太阳系之外的行星。截至2021年10月20日，已经被确认的系外行星总共有4531颗，当中约有四分之三是通过凌日现象发现的。这些行星分属3364个行星系，其中有778个多行星系。开普勒任务 (Kepler Mission) 已经检测到18,000颗行星候选者，包括262颗位于潜在适居带。
	
	\subsubsection{潜在适居带 (potential habitable zone)}
	
	适居带 (或称宜居带) 指的是恒星周围，行星系中温度适合液态水存在、适合生命存在的区域，根据恒星类型和大小而变化。适居带可能像地球一样出现高等生命。
	
	地球相似度指数 (Earth Similarity Index, ESI) 根据行星半径、密度、质量、逃逸速度、表面温度、处在宜居带的位置等，通过公式量化打分，取值0$\sim$1之间，0代表完全不同，1代表完全相同。
	
	\subsubsection{于开普勒太空望远镜}
	
	开普勒太空望远镜 (KEPLER) 是世界首个用于探测太阳系外类地行星的飞行器，于美国东部时间2009年3月6日发射。在为期至少3年半的任务期内，KEPLER将对天鹅座和天琴座中大约10万个恒星系统展开观测，以寻找类地行星和生命存在的迹象。KEPLER携带的光度计装备有直径为95厘米的透镜，通过观测行星的“凌日”现象搜寻太阳系外类地行星。
	
	迄今KEPLER经历的重要事件有：发现首颗系外行星；发现质量和直径都最小的系外行星；发现首个6行星系统；发现首个围绕两个太阳运行的行星；发现位于宜居带中，围绕一颗类太阳恒星运行的最小行星等等。
	
	\subsection{行星大气中的放电现象}
	
	闪电不是地球的专利！系外行星上也有闪电！
	
	具有闪电的星体可以分为6类：类地行星、水行星、没有液态表面的岩石行星、类金星行星、巨大的气体行星，以及褐矮星。
	
	火山活动较活跃的行星很可能具有最高的闪电发生频率，这其中就包括了岩石行星Kepler-10b和55 Cancrie，前者被形容为一颗“酷热的、潮汐锁定的岩石超级地球”。
	
	\subsubsection{闪电可能在生命形成中扮演重要作用}
	
	原始的地球，因为火山的爆发，大气的成分以二氧化碳、氢气、氨气、甲烷、二氧化硫等为主。
	
	此时全球已经遍布海洋，水循环和大气循环已经开始，地面和空气中的电荷形成了雷电 (物理)，雷电击中火山附近高温的海洋，在雷电和温度的催化下，水中的有些无机物质开始产生了化学反应，合成一些以碳原子为骨架的有机物 (化学)。
	
	原始海洋中开始到处都是这些有机小分子，因此原始海洋被形象地称为“有机汤”。到处都是有机物的海洋，有火山源源不断地热能和大气中狂暴的雷电，有机物们相互反应，形成生命只是一个概率问题 (数学).
	
	\subsubsection{米勒实验 (Miller-Urey experiment)}
	
	米勒模拟实验（Miller–Urey experiment）一种模拟在原始地球还原性大气中进行雷鸣闪电能产生有机物 (特别是氨基酸)，以论证生命起源化学进化过程的实验。1953年由美国芝加哥大学研究生米勒 (S. L. Miller) 在其导师尤利 (H. C. Urey) 指导下完成。
	
	1、将水注入左下方的500毫升烧瓶内；
	
	2、先将玻璃仪器中的空气抽去，然后打开左方的活塞，泵入CH4、NH3和H2的混合气体 (模拟还原性大气)；
	
	3、再将500毫升烧瓶内的水煮沸，使水蒸汽 (H2O) 和混合气体同在密闭的玻璃管道内不断循环；
	
	4、并在另一容量为5升的大烧瓶中，经受火花放电 (模拟雷鸣闪电) 一周，最后生成的有机物，经过冷却后，积聚在仪器底部烧杯的溶液内 (模拟原始大气中生成的有机物被雨水冲淋到原始海洋中)。
	
	\includegraphics[width=0.5\textwidth]{16.png}
	
	\textbf{结果推测}
	
	原始大气在高温、紫外线及雷电等自然条件的长期作用下，形成了许多简单的有机物，这是一个化学过程，这说明在原始地球条件下，有无机物到有机物是完全可能的。
	
	\textbf{质疑}
	
	（1）米勒试验提供持续的电能，但是原始时代的地球不一定。
	
	（2）不能完全确定米勒试验各物质浓度的配比。
	
	（3）氨基酸很可能是宇宙流星和彗星在撞击地球的时候带出的，因为当时这种现象十分普遍，科学证明氨基酸可以在宇宙的恶劣环境中存在。
	
	（4）由于不能确定烧瓶中的水的无菌性，且取样器与烧瓶为一联通器，因此米勒实验所产生的有机物不排除是微生物高温分解得来的。
	
	
	\subsubsection{火星上的沙尘暴闪电}
	
	火星沙尘暴会显著影响火星气候、大气环流及水、二氧化碳的交换，改变地表反照率，并对火星表面的探测器构成严重威胁，降低太阳能转换效率并可能导致其失联。全球性沙尘暴还能影响电离层高度，而高层大气中的沙尘颗粒则可能形成稀薄的云层。
	
	火星沙尘暴内部摩擦产生的电荷分离可能引起一些小尺度放电，并导致了类似于舒曼共振的辐射信号 (Rufet al., 2009GRL)
	
	目前尚未探测到长尺度放电通道的无线电信号 (尽快在观测期期间有两次比较大的沙尘暴和很多小的沙尘暴) (Gurnett et al., 2010GRL)
	
	
	\subsection{引力弹弓效应}
	引力弹弓效应：利用行星重力场给太空探测船加速，即把行星当作“引力助推器”。
	
	引力弹弓就是利用行星的重力场来给太空探测船加速，将它甩向下一个目标，也就是把行星当作“引力助推器”。
	
	利用引力弹弓能够探测冥王星以内的所有行星。在航天动力学和宇宙空间动力学中，所谓的引力助推 (也被称为引力弹弓效应或绕行星变轨)是利用行星或其他天体的相对运动和引力改变飞行器的轨道和速度，以此来节省燃料、时间和计划成本。
	
	引力助推既可用于加速飞行器，也能用于降低飞行器速度。\\
	\includegraphics[width=1\textwidth]{17.png}
	
	\section{大气探测和人工影响天气}
	
	\subsection{大气探测的内容}
	大气探测是对表征大气状况的气象要素、天气现象及其变化过程进行个别或系统的连续的观察和测定，并对获得的记录进行整理。随着科学技术的发展，大气探测的要素量和空间范围越来越大。
	
	大气探测分为近地面层大气探测、高空大气层探测和专业性大气探测。
	
	近几十年来，作为\underline{主动遥感的各种气象雷达探测}和\underline{作为被动遥感的气象卫星探测}，以及\underline{地面微波辐射探测}等获得较多信息的大气探测方法，正在逐步进入常规大气探测领域。这些现代大气探测技术应用于大气科学的研究领域，极大的丰富了大气探测的内容。
	
	\subsection{气象观测场}
	\subsubsection{观测要素}
	\textbf{基础气象要素}：如温度、湿度、气压、风速、降水等。
	
	\textbf{大气成分}：包括温室气体 (如CO₂)、气溶胶、反应性气体、臭氧等。
	
	\textbf{特殊现象}：如云、雾、降水、湍流等天气现象的宏观与微观特性。
	
	\subsubsection{探测方式}
	\textbf{直接探测}：通过地面观测站、探空仪、飞机等设备直接测量大气参数。
	
	\textbf{间接探测 (遥感)}：
	
	\textbf{1)主动遥感}：如雷达 (测云雨雷达、风廓线雷达) 发射电
	磁波并接收回波信号。
	
	\textbf{2)被动遥感}：如卫星接收大气辐射信号 (如紫外、红外波
	段监测臭氧)。
	
	\subsection{地表大气状态测量的六要素}
	
	传感器分为模拟传感器、数字传感器和智能传感器三大类，常用的六要素自动气象站传感器有：
	
	\textbf{气压}——振筒式气压传感器、膜盒式电容气压传感器；
	
	\textbf{温度}——铂电阻温度传感器；
	
	\textbf{湿度}——湿敏电容湿度传感器；
	
	\textbf{风向}——单翼风向传感器；
	
	\textbf{风速}——风杯风速传感器；
	
	\textbf{雨量}——翻斗式雨量传感器、容栅式雨量传感器；
	
	数据采集器是地面气象观测站的核心，其主要功能是数据采样、数据处理、数据存储及数据传输.
	
	\textbf{自动气象站} 是一种能自动地观测和存储气象观测数据的设备，主要由传感器、采集器、系统电源、通信接口、计算机和采集软件、业务软件等组成，能自动采集气压、温度、湿度、风向、风速、雨量、蒸发量、日照、辐射、地温等全部或部分气象要素。
	
	\subsubsection{温度测量}
	温度是表示物体冷热程度的物理量。
	
	在热力学和气象学研究中，常用的温标包括摄氏温标 ($^{\circ}\text{C}$)、华氏温标 ($^{\circ}\text{F}$) 和热力学温标（开尔文，$\text{K}$）。以下是三者之间的主要转换公式：
	
			\textbf{摄氏度转华氏度：}
			$$F = C \times \frac{9}{5} + 32$$
			
			\textbf{摄氏度转开尔文：}
			$$	K = C + 273.15$$
			
			\textbf{华氏度转开尔文：}
			$$	K = (F - 32) \times \frac{5}{9} + 273.15$$
	\begin{enumerate}
		\item \textbf{Steadman 通用公式}
		1984年，罗伯特·史特德曼 (Robert G. Steadman) 发表了《体感温度的通用公式》(A universal scale of apparent temperature)，其数学表达式如下：
		\begin{equation}
			AT = 1.07T + 0.2e - 0.65V - 2.7
		\end{equation}
		其中，水汽压 $e$ 的计算公式为：
		\begin{equation}
			e = \frac{RH}{100} \times 6.105 \times \exp \left( \frac{17.27T}{237.7 + T} \right)
		\end{equation}
		
		\item \textbf{参数定义与单位}
		\begin{itemize}
			\item $AT$：体感温度 ($^{\circ}\text{C}$)
			\item $T$：气温 ($^{\circ}\text{C}$)
			\item $e$：水汽压 ($\text{hPa}$)
			\item $V$：风速 ($\text{m/s}$)
			\item $RH$：相对湿度 ($\%$)
		\end{itemize}
		
		\item \textbf{物理意义与应用环境}
		该公式主要针对\textbf{暑热天气}设定，揭示了湿度和风速对人体散热的影响：
		\begin{itemize}
			\item \textbf{湿度影响：} 相对湿度 ($RH$) 越大，水汽压 ($e$) 越高，人体汗液蒸发散热受阻，导致体感温度升高（感觉更热）。
			\item \textbf{风速影响：} 风速 ($V$) 越大，人体表面的对流散热越快，体感温度越低（感觉较凉爽）。
			\item \textbf{结论：} 在高温环境下，\textbf{相对湿度越大、风速越小}，体感温度显著高于实际气温。
		\end{itemize}
	\end{enumerate}
	
	\subsubsection{气压的测量}
	
	气压是大气压强的简称，是在任何表面的单位面积上，空气分子运动所产生的压力。物理学上压强的国际单位为帕斯卡，气象上常用百帕（hPa）、毫巴（mb）和毫米汞柱（mmHg）衡量气压。标准条件下，单位面积上承受760mmHg的压力，将其作为一个标准大气压力。
	\textbf{换算公式:}
	$$1hPa=0.750069mmHg\approx 3/4mmHg$$
	$$1mmHg=1.333238hPa\approx 4/3hPa$$
	
	\subsubsection{湿度的测量}
	
	表示大气中水汽含量程度的物理量叫做湿度。
	
	表征大气中水汽含量的物理参数：
	
	1）混合比r：湿空气中水汽质量与干空气质量之比
	
	2）比湿q：湿空气中水汽质量与空气总质量之比
	
	3）绝对湿度ρ：单位容积湿空气中所含的水汽质量
	
	4）水汽的摩尔分数
	
	5）水汽压e
	
	6）饱和水汽压E
	
	7）相对湿度u
	
	8）露点温度Td和霜点温度Tf
	
	\subsubsection{风向和风速的测量}
	
	\textbf{风向标} 是测量风向最通用的仪器，一般由尾翼、指向杆、平衡锤及旋转主轴组成。
	
	\textbf{风杯} 是测定风速最常用的传感器，杯型风速器的主要优点是它与风速无关。
	
	\textbf{压管风速计} 利用伯努利原理，测量流体的全压力和静压力之差，进而测量风速。
	
	\subsubsection{降水的测量}
	
	降水量是指降落在地面上未经蒸发、渗透和流失的液态或固态降水的积水量；降水量以积水深度表示。
	
	常用的降水量测量仪器：雨量器;虹吸式雨量器;翻斗式雨量器.
	
	当有液态降水时，降水从承水器经漏斗、进水管引入浮子室；降水使浮子上升，带动自记笔画出曲线；当浮子室的水上升到虹吸管顶端，虹吸管开始排水，重新开始记录。
	
	当有液态降水时，降水从承水器经漏斗进入上翻斗；当雨水积到一定量时，翻斗翻转，水进入汇集漏斗；而后降水进入计量翻斗和计数翻斗，记录计数翻斗翻转的次数，进而换算降水量。
	
	\subsection{大气测量技术}
	\subsubsection{遥感测量技术-气象雷达}
	气象雷达是专门用于大气探测的雷达，属于主动式微波大气遥感设备；它以云、雨、风等气象现象作为主要观测目标，提取它们的信息，根据其回波来识别判断和预报此时或未来的天气。
	\subsubsection{激光雷达 (LIDAR-LIght Detection And Ranging)}
	基本原理：激光定向发射到大气中，大气示踪物 (大气分子、原子或者气溶胶)对激光产生散射现象，其中后向散射的部分 (与发射激光同一路径) 被接收望远镜收集。通过光电倍增管将弱的光子信号加强，在被采集卡所采样，最后得到
	散射收集的光子计数。根据相应的反演算法，得到需要的物理参数。
	
	\subsection{卫星遥感探测技术}
	
	卫星气象观测是通过人造地球卫星搭载可见光、红外、微波等遥感器探测地球大气物理现象的技术，具有观测范围广、时效快、观测频次高且不受自然地理条件及国界限制的特点，数据传输采用国际电信联盟分配的专用频段。该技术起源于20世纪60年代。
	
	\subsubsection{中国空间站“三大黑科技”}
	1)整齐合理的核心舱内部管线布局.
	
	2)首次使用全新的推进技术-霍尔推进器, 减少推进剂存储空间，更有效的利用航天员有限活动空间.
	
	3)利用机械臂进行自主对接和外部检测.
	
	\subsubsection{霍尔效应 (Hall Effect)}
	
	霍尔效应是电磁效应的一种，是美国物理学家霍尔 (E. H. Hall，1855-1938) 于1879年在研究金属的导电机制时发现的。
	
	当电流垂直于外磁场通过半导体时，载流子发生偏转，垂直于电流和磁场的方向会产生一附加电场，从而在半导体的两端产生电势差，这种现象就是霍尔效应，这个电势差也被称为霍尔电势差。霍尔效应使用左手定则判断。
	
	\subsubsection{霍尔推进器 (Hall thruster)}
	
	交叉电磁场捕获从阴极发射的电子，电子绕磁力线旋转并在放电区内作角向漂移，形成的电子电流称为霍尔电流。而角向漂移是交叉的径向磁场与轴向电场作用的结果(即霍尔效应)，因此称为霍尔推进器。
	
	角向漂移电子与通过阳极进入环形放电室的推进剂分子发生碰撞后电离，形成等离子体，其中离子在电磁场的作用下沿轴向加速，并高速喷出，从而产生推力。
	
	\subsubsection{气象卫星}
	
	气象卫星是从太空对地球及其大气层进行气象观测的人造地球卫星。
	
	\subsubsection{风云三号E星“黎明星”}
	
	作为全球首颗民用晨昏轨道气象卫星， FY-3E卫星综合观测能力有四大特色：
	
	1、高精度光学微波组合大气温度湿度垂直分布探测能力；
	
	2、主动遥感仪器风场精确探测能力；
	
	3、全球气候变化监测与预测的支撑能力；
	
	4、太阳和空间环境综合探测能力。
	
	\subsubsection{掩星技术和GNSS气象学}
	
	受电离层、平流层和对流层大气的影响，GPS或北斗等GNSS卫星导航系统的信号在穿过地球大气层时会发生传输路径的改变而产生延迟量。利用GNSS卫星信号延迟量，可以反演大气的水汽、温度和压力等气象要素。
	
	GNSS探测几乎不受天气状况的影响，具有高垂直分辨率、高覆盖率、高精度、全天候和长期稳定的优点，能够有效弥补目前常规探空资料的不足，可以对对流层顶和平流层下部等常规探测难以达到的区域进行探测。
	
	\subsection{雷电探测技术和应用}
	
	\subsubsection{雷电探测要素的多样化}
	光 (光学成像、光谱测量及高能射线探测)
	
	电\&磁 （闪电本质上就是电荷的运动）
	
	声 (雷声thunder的测量，可用于定位雷声源)
	
	热 (比如：闪电通道的温度有多高？)
	
	\subsubsection{雷电探测平台的多样化 (多为被动探测)}
	地面场地试验
	
	近地探空平台（飞机、无人机和探空气球等）
	
	卫星 (极轨、地球同步和行星际)
	\subsubsection{闪电熔岩 (Fulgurite)}
	
	闪电熔岩是天然造成的玻璃长管。当闪电击中泥土或沙, 就可能令它们迅间熔化, 然后又凝固, 便会形成闪电熔岩, 因此闪电熔岩的形状多是长条状, 和闪电的路径相近。闪电熔岩长可达数米。颜色由形成的泥土和沙决定, 有黑色、绿色、白色等。内部光滑, 可能有小气泡, 外部多数为粗糙的沙粒。
	
	\subsubsection{雷电的功绩}
	第一，可制造化肥。雷电过程中闪电温度极高，并能造成高电压。在高温、高电压条件下，
	空气分子会电离；重新结合时，其中的氮和氧会化合为亚硝酸盐和硝酸盐分子，并溶解在雨
	水中降落地面，成为天然氮肥。据测算，全球每年因雷电落到地面的氮肥有四亿吨。
	
	第二，雷电能促进植物生长。雷电在发生时，地面和天空间电场强度可达到每厘米万伏以上。
	受这样强大的电位差的影响，植物的光合作用和呼吸作用增强，因此，雷雨后一两天内植物
	新陈代谢特别旺盛。例如，有人用闪电刺激作物，发现豌豆提早分枝，而且分枝数目增多。
	
	第三，雷电能制造负氧离子。负氧离子又被称为空气维生素，可以起到消毒杀菌、净化空气
	的作用。在雷雨后，空气中高浓度的负氧离子，使得空气格外清新，人们感觉心旷神怡。
	
	第四，雷电还有巨大的能量。地球上平均每秒有100次闪电，一次闪电释放约8000瓦小时的
	电能，因此，每年全世界的雷电约放出250亿千瓦小时的能量。遗憾的是，人类目前还无法
	对它加以利用。
	\subsection{人工影响天气的基本原理}
	人工影响天气，指用人为手段使天气现象朝着人们预定的方向转化。
	
	人工影响天气业务是指为避免和减轻气象灾害，合理利用气候资源，在适当条件下通过科技手段对局部大气的物理、
	化学过程进行人工影响，实现增雨、防雹、消雨、消雾、森林防（灭）火、消除公共污染等目的业务活动。
	
	\subsubsection{什么是人工影响天气？}
	
	人工影响天气是通过科技手段对局部大气的物理、化学过程进行人工影响来减轻或者避免气象灾害，在合理条件下，利用气候资源，达到增雨雪、防雹、消雨、消雾、防霜等目的。也有利用小型飞机、高射炮等运载工具进行人工增雨、防雹作业，向云中播撒催化剂；在一些农田进行人工防霜，以及在机场、航道、高速公路等进行人工消雾等。
	
	\subsubsection{发展人工影响天气技术的意义}
	
	\textbf{趋利避害 减灾增收}
	自1958年我国开展人工影响天气工作以来，在各级政府的大力支持下, 人工影响天气业务体系日渐完善，开发云水资源能力不断增强，在农业抗旱、防智减灾、缓解水资源短缺、改善生态环境等方面发挥了显著效益.
	
	\subsubsection{人工增雨的原理和方法}
	对于冷云：人为地增加云内冰晶的浓度，如撒干冰或碘化银；
	
	对于暖云：撒播吸湿性物质（凝结核），如食盐或氯化钙。
	
	\textbf{成云降雨的三个条件} 一定的的疑结核和冰核, 上升的气流, 充足的水汽.
	
	人工影响天气是改变冰核（冷云催化）和凝结核（暖云催化）.人工增雨需要的催化剂分为成冰剂、致冷剂和吸湿剂三类。
	
	\subsubsection{人工消雹的原理和方法}
	
	1)撒入大量冰核，使过冷水滴形成很多小冰雹，不能形成大冰雹。
	2)在云底的上升气流中撒入大量吸湿性核，使云内过冷却水降落，弱化冰雹增长的条件。
	3)扰乱云中的上升气流，破坏冰雹增长的条件。
	
	\subsection{天气预报的方法}
	天气预报技术是现代气象学的核心应用，其制作流程融合了数据采集、数值计算、人工分析和多平台发布等多个环节。
	
	天气预报（测）或气象预报（测）是使用现代科学技术对未来某一地点地球大气层的状态进行预测。从史前人类就已经开始对天气进行预测来相应地安排其工作与生活 (比如农业生产、军事行动等)。今天的天气预报主要是使用收集大量的数据 (气温、湿度、风向和风速、气压等)，然后使用目前对大气过程的认识 (气象学) 来确定未来空气变化。由于大气过程的混乱以及今天科学并没有最终透彻地了解大气过程，因此天气预报总是有一定误差的。
	
	\subsubsection{按时间范围划分}
	即按天气预报的时效长短，可分为：
	
	1．短时预报。根据雷达、卫星探测资料，对局地强风暴系统进行实况监测预报未来1-6小时的动向。
	
	2．短期预报。预报未来24-48小时天气情况。
	
	3．中期预报。对未来3-15天的预报。
	
	4．长期预报。指1个月到1年的预报。
	
	5．超长期预报。预报时效1-5年的。
	
	6．气候展望。10年以上的。
	
	\subsubsection{现状}
	
	短时强降水预报准确率不足25\%.
	
	冰雹大风预报准确率仅为4.5-5\%.
	
	极值报不出、落区不准确、时段偏差大.
	
	\textbf{灾害天气机理复杂} 多尺度气象系统相互作用;触发和发展机制不明;尺度小，结构变化剧烈.
	
	\section{气象学和空间圈层物理}
	
	\subsection{气候、气象与天气的区别}
	
	\textbf{气候} 是指整个地球或其中某一个地区一年或一段时期的气象状况的多年特点 (或平均天气状况)；通常以几十年为统计周期。它反映了该地区相对稳定的大气状态和风、降水等现象的总体趋势。
	
	\textbf{气象} 是指发生在天空中的风、云、雨、雪、霜、露、虹、晕、闪电、打雷等一切大气中的物理现象。
	
	\textbf{天气} 是指某一地点或区域在短时间范围内 (或在影响人类活动瞬间) 的气象特点综合状况，如温度、湿度、风速、云量等。这些因素可能在短时间内发生显著变化，通常按小时或天来预报。
	
	气候是气象的背景，天气是气象对人的影响。
	
	\subsection{地球大气中的环流}
	
	大气环流是指具有全球规模的、大范围的大气运行现象。
	
	\subsubsection{大气环流形成原因}
	1、太阳辐射是地球上大气运动能量的来源，由于地球的自转和公转，地球表面接受太阳辐射能量是不均匀的。热带地区多，而极区少，从而形成大气的热力环流。
	
	2、地球自转：在地球表面运动的大气都会受地转偏向力作用而发生偏转。
	
	3、地球表面海陆分布不均匀，以及大气内部南北之间热量、动量的相互交换。
	
	\subsubsection{热带大气对流活动}
	热带西太平洋暖池是全球海洋最大的暖水区域，可谓全球“气候心脏”。 其高海温促进了大气强烈的对流运动，是驱动热带大气环流的主要热源，也是东亚季风和厄尔尼诺现象的策源地。
	
	\subsubsection{沃克环流 (Walker circulation)}
	沃克环流是赤道海洋表面因水温的东西部差异而产生的一种纬圈热力环流。它是热带太平洋上空大气循环的主要动力之一。
	\subsection{云物理学和灾害性天气}
	云物理学属于大气科学分支，核心内容包括云粒子微物理过程和云体宏观动力结构的演变规律。
	
	云物理学以大气热力学和动力学为基础，研究云、雾及降水的形成、发展规律及其对天气、气候的影响。其研究对象尺度从微米级云滴延伸至天气系统级别，核心机制涉及水汽相变与云雨演变的矛盾关系。
	
	\subsubsection{研究云和降水的重要性}
	1)主要的灾害性天气如暴雨、雷暴、冰雹、台风、龙卷风等均与云雨过程相联系.
	
	2)云成年覆盖了地球一半的面积，对地球大气系统的辐射传输过程影响极大.
	
	3)云和降水可以为地面提供新鲜的水源.
	
	4)云以释放潜热的形式向大气输送热量，是大气运动的重要能量来源.
	
	5)云和降水影响光和电磁波在大气中的传播.\\
	\includegraphics[width=1.0\textwidth]{18.png}
	
	\subsubsection{云的成因及条件}
	充足的水汽和上升气流是云生成的必要条件.
	
	不同的上升运动，可以生成不同种类的云.
	
	\subsubsection{云的种类}
	
	\textbf{高云} 卷云、卷层云、卷积云
	\textbf{中云} 高层云、高积云
	\textbf{低云} 层云、层积云、雨层云
	\textbf{直展云系} 积云、积雨云
	
	\subsubsection{雨滴的形成}
	小的云滴粒子在气流托举下抬升，不断碰并，形成较大的云滴粒子，开始下坠，且小粒子不断加入，最终形成雨滴。
	
	\subsubsection{雷暴(thunderstorm)}
	雷暴是热带和温带地区可见的局地性强对流天气。雷暴发生时可伴随有闪电、强风和强降水，例如雨或冰雹。雷暴可发生于春季和夏季，常见的例子是夏季午后，但也可能在冬季随暴风雪发生，被称为雷雪（thundersnow）。
	
	雷暴的持续时间通常不超过2小时，其生命周期包括积云阶段(cumulus stage)、成熟阶段(mature stage)和消散阶段(dissipating stage)。
	
	\subsubsection{雷暴发展的条件}
	按中尺度环流结构，雷暴可分为4类：单体雷暴(single cell thunderstorm)、多单体雷暴(multi-cell thunderstorm)、飑线(squall line)和超级单体雷暴(super cell thunderstorm)。
	
	按触发机制，雷暴可分为热雷暴(thermal thunderstorm)、锋雷暴(frontal thunderstorm)和地形雷暴(orographic thunderstorm)。
	
	\subsubsection{穿透性对流(overshootingtop)}
	垂直伸展高度高, 上升运动速度大.
	
	穿透性对流带入平流层的大量水汽，影响平流层的一些化学过程及辐射平衡.
	\subsubsection{龙卷风及其强度分档}
	龙卷风是从积状云悬垂并向下伸展至与地面接触的旋转空气柱，是大气中最猛烈的涡旋，在近地面常表现为漏斗云、旋转的沙尘碎屑等现象，在短时间内易造成局地重大人员伤亡和财产损失。
	
	由于龙卷风空间尺度小，具有突发性、持续时间相对较短等特征，其监测和预报预警仍是世界难题。
	
	\subsubsection{龙卷风怎么形成的？}
	(A) 当不同高度的风速不同时，旋转的空气可能形成风切变。
	
	(B) 如果旋转的空气被卷入上升的风流（上升气流）中，旋转会加快并收紧，形成漏斗云。
	
	(C) 雷暴中的雨和冰雹会使漏斗向下弯曲。如果它触碰地面，就形成龙卷风。
	\subsubsection{干雷暴 (dry storm)}
	干雷暴，通俗的理解就是“干打雷不下雨”的电闪雷鸣天气现象，是在小范围内的一次短暂强对流天气。在中国东北地区，春季气候干燥，到每年5-7月份春夏之交，林区极易发生由干雷暴引发的森林雷击火灾。
	
	干雷暴引发的火灾有地点不确定、扑救难度大等特点，往往需要人工进行扑救。因此，预防和扑救干雷暴引发的森林大
	火是世界性难题。
	
	没有降水的闪电、雷鸣现象，称干雷暴，是一种多发于澳洲的特殊天气情况。干雷暴是由于天气炎热干燥，上层空气的云层遇到冷空气降雨。但是雨还没落到地面，由于下层地表高温的缘故，雨马上被蒸发变为潮湿的热空气再次上升到空中。由于下雨时会有雷电产生，并伴有大风，极易引起森林大火。
	
	\subsection{大气污染和气象效}
	大气污染物，指由于人类活动或自然过程排入大气的并对人和环境产生有害影响的物质。
	
	世界卫生组织列出了六种“典型”空气污染物：一氧化碳 (CO)、铅 (Pb)、二氧化氮 (NO2)、悬浮颗粒物 (SPM) (包括尘土、烟灰、烟雾和烟尘)、二氧化硫 (SO2)和对流层臭氧。
	
	\subsubsection{大气污染物来源}
	大气污染物主要来源包括工业排放、交通尾气、生活燃煤等，自然源如火山喷发也会释放污染物。
	
	大气污染物按存在状态可分为气态污染物(如二氧化硫、氮氧化物) 和颗粒物 (如PM2.5、PM10)
	
	按形成过程分为：
	
	\textbf{一次污染物}：直接排放的原始物质 (如汽车尾气中的CO)
	
	\textbf{二次污染物}：经化学反应生成的新物质 (如臭氧、光化学烟雾)
	
	\textbf{复合污染}：多种污染物相互作用形成的复杂污染
	\subsubsection{酸雨及其危害}
	酸雨泛指PH值小于5.6的雨雪或其它形式的大气降水，是大气受污染的一种表现，最早引起注意的是酸性的降雨，所以习惯上统称为酸雨。
	
	\textbf{酸雨可导致：}
	
	1. 土壤酸化，并能使土壤中的铝从稳定态中释放出来，严重地抑制林木的生长；
	
	2. 加速土壤矿物质营养元素的流失，改变土壤结构，导致土壤贫瘠化，影响植物正常发育；
	
	3. 诱发植物病虫害，使农作物大幅度减产；
	
	4. 抑制某些土壤微生物的繁殖；
	
	5. 使非金属建筑材料表面硬化水泥溶解，出现空洞和裂缝，导致强度降低，从而损坏建筑物；
	
	6. 使儿童免疫功能下降，慢性咽炎、支气管哮喘发病率增加，同时可使老人眼部、呼吸道患病率增加。
	
	\subsubsection{臭氧污染与光化学烟雾}
	所谓的臭氧污染永远和光化学烟雾分不开，日常所说的臭氧污染其实就是指的光化学烟雾。
	
	地表臭氧变化与人类健康和生态可持续发展息息相关.
	
	\subsubsection{天然氧吧与空气负 (氧) 离子}
	
	氧吧是指人们吸氧的地方。天然氧吧是就指自然条件下形成的氧吧，多指植被茂密、氧气含量大的地方。
	
	负氧离子，是指获得1个或1个以上的电子带负电荷的氧气离子。被誉为“空气维生素”的负氧离子有利于人体的身心健康。它主要是通过人的神经系统及血液循环对人的机体生理活动产生影响。空气中负氧离子浓度多少，是空气清新与否的标志。
	
	负氧离子对改善生态、生活环境确实可以起到积极作用，负氧离子中的小离子可以捕捉漂浮微尘，使其凝聚而沉淀，从而净化空气，当浓度达到20000个每立方厘米时，空气中的飘尘会减少98\%以上。
	
	\subsection{电离层物理}
	
	电离层 (Ionosphere)，是地球大气的一个电离区域。电离层受太阳高能辐射及宇宙线的激励而电离的大气高层。60千米以上的整个地球大气层都处于部分电离或完全电离的状态，电离层是部分电离的大气区域，完全电离的大气区域称磁层。
	
	除地球外，金星、火星和木星都有电离层。
	
	\subsubsection{全球大气电路}
	
	全球大气电路的等效模型:
	
	地球表面的电荷总量：500,000 库伦 (负电荷).
	
	电离层同地球表面之间的电流：2,000 安培.
	
	\textbf{总电子浓度 (Total Electron Content, TEC)} 电离层总电子浓度（TEC）是衡量特定路
	径上电子总量的参数，单位通常为TECU($1 TECU=10^{16}$个电子/平方米)。
	
	\subsubsection{南大西洋异常区 (South Atlantic Anomaly)}
	南半球有一大块区域 (覆盖范围遍及南美洲南部及南大西洋海域)的 磁 场 强 度 比 其 他 地 区 弱30~50\%，弱化速度也比其他地区快10倍，被称做南大西洋异常区 (South Atlantic Anomaly,SAA)。南大西洋异常区的磁场甚至出现磁场翻转，产生一个局部朝南的磁北极。
	
	由于该区地磁较弱，阻挡太阳粒子的范艾伦辐射带在该区域上空形成一凹陷部分，让粒子可以到达更接近地球的位置，导致穿越该区域上空的人造卫星受粒子影响而出现运作异常。
	
	\subsubsection{临近空间}
	
	临近空间（Near Space）通常指海拔 20 km 至 100 km 之间的区域。这一区域位于传统航空器（约 20 km 以下）的最高升限之上，以及轨道航天器（约 100 km 以上，即卡门线）的最低轨道高度之下，被称为“航空航天领域的空白地带”。
	
	\begin{enumerate}
		\item \textbf{传统飞行器难以涉足的原因（技术挑战）}
		\begin{itemize}
			\item \textbf{空气过于稀薄（动力瓶颈）：} 
			在此高度，空气密度仅为海平面的几分之一到万分之一。对于传统航空器（飞机），稀薄的空气无法产生足够的\textbf{气动升力}；同时，由于含氧量极低，依靠吸气式发动机（如涡喷、涡扇）的飞行器无法获得足够的\textbf{氧化剂}供燃料燃烧。
			\item \textbf{空气过于浓密（轨道瓶颈）：} 
			对于航天器（卫星），虽然此处空气稀薄，但产生的\textbf{气动阻力}仍足以让卫星在极短时间内丧失速度并坠毁。因此，这一高度无法维持稳定的近地轨道飞行。
			\item \textbf{环境极端恶劣：} 
			该区域不仅存在强紫外线辐射和高能粒子，还跨越了臭氧层产生的增温区和中间层的极寒区（$-90^{\circ}\text{C}$ 以下），对飞行器的材料耐受性和热控系统提出了严苛要求。
		\end{itemize}
		
		\item \textbf{临近空间的科学与应用价值}
		尽管环境恶劣，临近空间却是连接地球与太空的枢纽，具有独特的科研意义：
		\begin{itemize}
			\item \textbf{大气科学与全球气候研究：} 
			该层包含臭氧层，是研究全球变暖、平流层与对流层能量交换的核心地带。通过监测此处的化学成分变化，可以更准确地预测全球气候走势。
			\item \textbf{空间天气监测：} 
			临近空间与电离层紧密相邻，是研究太阳风、极光现象以及地磁暴对无线电通信和 GPS 信号干扰的最佳观察点。
			\item \textbf{新型观测平台（驻空优势）：} 
			相比卫星，临近空间平台（如高空科学气球、太阳能无人机）距离地面更近，可实现\textbf{亚米级的高分辨率对地观测}，且成本较低，可长期驻留于特定区域上方。
			\item \textbf{深空探测模拟：} 
			临近空间的大气密度与压力环境与\textbf{火星表面}非常相似。因此，它常被用作火星探测器（如降落伞测试、着陆器实验）的地球模拟试验场。
		\end{itemize}
		
		\item \textbf{总结}
		临近空间目前已成为大国科技竞争的战略制高点。随着高超声速推进技术和新型复合材料的发展，人类正在填补这一“空白地带”，使其从“难以逾越的障碍”转变为“通向太空的新通道”。
	\end{enumerate}
	\subsubsection{GNSS 系统导航定位原理 (Principles of GNSS Navigation and Positioning)}
	
	GNSS（如 GPS、北斗等）定位的核心原理是通过多颗卫星的已知空间位置和信号传播时间，利用空间几何交会的方法确定接收机的位置和时间。
	
	\begin{enumerate}
		\item \textbf{基本定位原理}
		\begin{itemize}
			\item \textbf{已知量：} 至少 4 颗导航卫星的三维空间坐标 $(x_i, y_i, z_i)$ 以及信号发射的时间 $t_i$。
			\item \textbf{观测量：} 导航卫星与接收机之间的距离（伪距）$d_i$。
			\item \textbf{待求量（4个未知数）：} 接收机的三维坐标 $(x, y, z)$ 以及接收机相对于卫星系统的时间偏差 $t$。
		\end{itemize}
		
		\item \textbf{数学模型：四个距离方程}
		为了求解 4 个未知数（三维空间坐标和 1 个时间参数），必须同时观测到至少 4 颗卫星。其距离方程组如下（其中 $c$ 为光速）：
		\[
			\begin{cases}
				\sqrt{(x_1 - x)^2 + (y_1 - y)^2 + (z_1 - z)^2} + c(t_1 - t) = d_1 \\
				\sqrt{(x_2 - x)^2 + (y_2 - y)^2 + (z_2 - z)^2} + c(t_2 - t) = d_2 \\
				\sqrt{(x_3 - x)^2 + (y_3 - y)^2 + (z_3 - z)^2} + c(t_3 - t) = d_3 \\
				\sqrt{(x_4 - x)^2 + (y_4 - y)^2 + (z_4 - z)^2} + c(t_4 - t) = d_4
			\end{cases}
		\]
		
		\item \textbf{物理意义与流程}
		\begin{itemize}
			\item \textbf{空间交会：} 每一颗卫星对应一个球面方程，三个球面相交可以确定空间中的一点。
			\item \textbf{时间偏差：} 由于接收机时钟通常不如卫星上的原子钟精确，因此必须引入时间参数 $t$ 作为一个独立未知数进行解算。
			\item \textbf{结论：} 通过解这四个非线性方程组，接收机即可实时获得其精确的\textbf{三维坐标和授时信息 $(x, y, z, t)$}。
		\end{itemize}
	\end{enumerate}
	\subsection{不同圈层间的物质和能量交换}
	地球的四大圈层 (大气圈、水圈、岩石圈、生物圈) 通过物质循环和能量流动相互作用，形成复杂的生态系统。
	
	中高层大气研究是20世纪70年代以来持续得到大气科学界和日地物理学界共同关注的重要研究领域，是当前地球系统科学研究领域的重点方向。
	
	\subsubsection{对流层同平流层之间的物质输送}
	对流层与平流层之间的物质输送主要通过垂直运动和水平输送两种方式实现，涉及臭氧、水汽等关键物质的交换，对大气成分和气候系统有重要影响。
	\subsubsection{火山喷发如何影响气候？}
	\includegraphics[width=1.0\textwidth]{20.png}
	
	\section{太阳物理和行星际环境}
	
	\subsection{宇宙大爆炸（Big Bang）}
	
	\textit{物质和能量，哪个更本质？}
	
	从现代物理学的角度来看，能量更为本质，因为物质是能量的一种表现形式。
	
	\textit{宇宙的主要成分是什么？}
	
	1)暗能量： 68.3\% 
	
	2)暗物质： 26.8\% 
	
	3)发光物质：0.4\%（恒星及其行星，发光气体，辐射）
	
	4)不可见的普通物质： 4.5\%（星系际气体，中微子，黑洞）
	
	宇宙平均密度:约0.01电子质量/立方厘米
	
	\subsubsection{奥博斯佯谬 (Olbers' Paradox)}
	\textit{夜空为什么是黑色的？}
	
	若宇宙是稳恒态而且无限的，则晚上应该是光亮而不是黑暗的。漆黑一片的夜空印证了宇宙并非稳恒态的，是大爆炸理论的证据之一。
	\subsection{恒星演化与太阳基本情况}
	
	恒星演化是恒星在生命过程中所经历急遽变化的序列。恒星依据质量，一生的范围从质量最大的恒星只有几百万年，到质量最小的恒星比宇宙年龄还要长的数兆亿年。
	
	一般而言，恒星的演化分为三个阶段：主要由恒星引力收缩提供能量的主序前 (pre-main sequence) 阶段、由恒星核心处的氢到氦的核聚变反应提供能量的主序 (main sequence) 阶段、以及待恒星核心处的氢消耗殆尽后，由氦、碳或更重元素的燃烧提供能量的主序后 (post-main sequence) 阶段。
	
	\subsubsection{赫罗图 (H-R diagram)：恒星演化的研究工具}
	赫罗图 (HRD) 是以恒星的绝对星等或光度相对于光谱类型或有效温度绘制的散布图，是了解恒星演化的重要工具。
	\subsubsection{太阳能量来源-热核聚变}
	太阳能量巨大，太阳核心每秒有数万吨氢化为氦，核聚变过程中超过5百万吨质量转化为能量。
	
	太阳每秒发出的能量超过地球有史以来耗用的总量。
	\subsubsection{氦闪 (helium flash)}
	氦闪是在中等质量恒星 (0.8到2.0个太阳质量) 的核心，或是白矮星表面堆积的氦突然开始的短暂失控核聚变 。
	\subsubsection{天文单位}
	日地平均距离 (149597870700米)，又称为一个天文单位(astronomical unit, AU)
	
	\subsubsection{太阳爆发活动的影响}
	来自太阳的耀斑和日冕物质抛射等爆发现象，会影响地球的磁场和电离层，可能会干扰人造卫星、影响导航定位精度、导致无线电通信中断、甚至可能引起电网故障而导致大范围停电。
	
	此外，在载人航天器飞行、运行，以及航天员出舱等活动中，空间天气都会对其安全产生影响。
	
	因此，需要时刻关注太阳活动，捕捉太阳爆发过程，为更准确的空间天气预报提供重要科学数据。
	
	\subsection{磁层}
	地球磁层是地球磁场与太阳风相互作用形成的一个彗星状保护屏障，旨在保护地球上的生命和环境免受有害的太阳和宇宙辐射的影响。
	
	地球磁层由地球核心液态金属的对流产生的主磁场延伸到太空中形成。在太阳风 (来自太阳的带电粒子流) 的持续吹袭下，磁层形态受到影响。
	
	\textbf{向阳面}-正对太阳风的方向，磁场被压缩，形成弓状的激波面和磁层顶，距离地表约5至10个地球半径。
	
	\textbf{背阳面}-背离太阳风的方向，磁场被拉伸，形成一条长达数百万公里的长尾巴，称为磁尾。
	
	\textbf{范艾伦辐射带}-磁层内部存在两个轮胎状的区域，捕获了大量来自太阳风的高能带电粒子，即内、外辐射带。
	\subsubsection{磁重联 (magnetic reconnection)}
	
	磁重联，即磁力线“断开”(break) 再“重新连接”(reconnect) 的物理过程。
	
	磁重联, 或磁力线重联(magnetic field line reconnection)，又称磁场湮灭，是天体物理中一种非常重要的快速能量释放过程，也是磁能转化为粒子的动能、热能和辐射能的过程。普遍认为太阳上的能量释放就是磁重联导致的。
	\subsubsection{范艾伦辐射带 (Van Allen radiation belt )}
	范艾伦辐射带是在地球附近的近层宇宙空间中包围着地球的大量带电粒子聚集而成的轮胎状辐射层，主要来源是被地球磁场俘获的太阳风粒子，这些带电粒子在范艾伦带两转折点间来回运动。
	
	2012年8月，NASA发射了一对范艾伦探测器，同时从不同位置和角度对辐射带进行观测。观测结果认为范艾伦辐射带负责加速粒子，而不是收集来自其他地方的高能粒子。
	
	当太阳发生磁暴时，地球磁层受扰动变形，而局限在范艾伦带的高能带电粒子大量泄出，并随磁力线于地球的极区进入大气层，激发空气分子产生美丽的极光。
	
	\subsection{极光与空间天气}
	
	空间天气是指由太阳活动引起的近地空间环境的短时变化现象，主要发生在距地面30公里以上的磁层、电离层及中高层大气等区域。
	
	它会影响到卫星运行、通信导航、电力系统等，表现为地磁暴、极光等，并且会通过电磁辐射、高能粒子流等形式影响地球空间。
	\subsubsection{极光 (aurora)}
		极光（Aurora），是一种绚丽多彩的等离子体现象，其发生是由于太阳带电粒子流
		（太阳风）进入地球磁场，在地球南北两极附近地区的高空，夜间出现的灿烂美丽
		的光辉。在南极被称为南极光，在北极被称为北极光。地球的极光是来自地球磁层
		或太阳的高能带电粒子流（太阳风）使高层大气分子或原子激发（或电离）而产生。
	
	\subsubsection{极光背后的灾害-太阳活动攻击}
	
	\includegraphics[width=1.0\textwidth]{21.png}
	
	\subsubsection{地磁感应电流}
	地磁感应电流 (Geomagnetically induced current,
	GIC) 是重要的空间天气灾害现象，强烈的GIC会危害地
	面长距离导体设备的运行，扰乱电网的正常运行，加速
	油气管道的腐蚀，甚至造成电网瘫痪、油气泄漏等重大
	事故，造成严重的经济损失和社会影响。
	
	\subsection{太阳系的气态行星}
	太阳系有Atmospheric 四颗气态行星circulation ，统称为类木行星 (Jovian planets)或巨行星。它们都位于小
	行星带以外的外太阳系。
	
	\textbf{木星}：太阳系中最大的行星，主要由氢和氦组成。
	
	\textbf{土星}：以其壮观的行星环系统而闻名，主要由氢和氦构成。
	
	\textbf{天王星}：常被归类为冰巨星 (而非典型的气态巨星)，含有较多“冰”物质 (如水、甲烷、氨)。
	
	\textbf{海王星}：也是一颗冰巨星，外观呈蓝色，有微弱的行星环。
	
	\subsubsection{木星大红斑}
	木星大气层中的大红斑是一团激烈的沿逆时针方向运动的下沉气流 (高气压风暴)，这个气流物质中含有大量的红磷化物，所以呈深褐色。
	
	木星的大红斑位于南纬23°处，长2万千米，宽1.1万千米。在大红斑中心部分有个小颗粒，是大红斑的核，其大小约几百千米。大红斑的寿命很长，可维持几百年或更久 (观测自1830年开始已经持续了193年)。
	
	大红斑是逆时针旋转的，周期大约是6个地球日，或14个木星日。2004年初，大红斑在经度的方向上只有一个世纪前的一半大小，而之前它的直径是40,000公里。
	
	\subsubsection{土星光环}
	
	土星光环的成因尚不清楚，推测可能是由彗星、小行星与较大的卫星相撞后产生的碎片组成的。
	
	土星的光环可以分成几个不同的部分，最明亮、最宽阔的是A环和B环，较暗的是C环。较宽的光环其实是由许多狭窄的小环组成的。
	
	\subsubsection{天王星 (Uranus)}
	
	天王星是第一颗使用望远镜发现的行星，几乎横躺着围绕太阳公转。
	
	天王星大气主要成分是氢、氦、甲烷和氘。据推测，其内部可能含有丰富的重元素。地幔由甲烷和氨的冰组成，可能含水。内核由冰和岩石组成。
	
	天王星是太阳系内大气层最冷的行星，最低温度为49K (-224℃)。
	
	\subsubsection{海王星 (Neptune)}
	
	海王星在1846年9月23日被发现，是仅有的利用数学预测而非观测意外发现的行星。天文学家利用天王星轨道的摄动推测出海王星的存在与可能的位置。
	
	海王星的大气层的化学组成以氢分子和氦为主。海王星有14颗已知的天然卫星。
	
	\subsection{行星际航行}
	
	行星际航行是指航天器脱离地球引力后在太阳系范围内开展的航行活动，涵盖环绕行星运行、飞越探测和着陆探测等任务形式。
	
	目前和未来行星际航行的关键技术包括：
	
	\textbf{轨道动力学}：利用霍曼转移轨道等技术计算和执行不同行星之间的最经济路径。
	
	\textbf{重力助推}：利用行星或卫星的引力改变航天器的速度和方向，从而节省燃料。
	
	\textbf{推进系统}：
	
	1) 化学火箭：目前主流推进方式，效率有限，行程耗时较长；
	
	2) 电力推进/离子推进器：提供持续的较小推力，具有更高的燃料效率，可缩短旅行时间；
	
	3) 核热火箭：利用核能加热推进剂产生推力，目前在研发中 (如NASA的DRACO项目)。
	
	\textbf{导航与制动}：精准的导航定位技术和大气制动 (利用行星大气减速)是成功到达目的地的关键。
	
	\subsubsection{火星的地球化方法}
	
	1) 在大气中增加适合适量的气体 (包括温室气体和适合生物生存的气体)，增加地表温度与气压，主要是为了液态水，其次是植物、动物。
	
	2) 在太空中架设巨大反射 (或折射) 镜群，增加照射到火星表面的太阳光强度。
	
	3) 大量融解地下冻土层，再把水引到地表。虽然一开始会结冰，但随着工程进行，冰层进而融化形成水圈。
	
	4) 在冰上 (包括两极) 培植深色藻类或散布煤灰等深色物质增加吸热进而加速融化。
	
	5) 散布固沙菌类、植物，防止沙暴的发生，进而生成土壤，扩大居住地。
	
	6) 建立行星推进器 (引擎)，改变火星运行轨道。
	
	\subsubsection{拉格朗日点}
	又称平动点，在天体力学中是限制性三体问题的五个特解。一个小物体在两个大物体的引力作用下在空间中的一点，在该点处，小物体相对于两大物体基本保持静止。
	
	\section{行星气候与全球变化}
	
	\subsection{行星气候的影响要素}
	地球气候的形成是\textbf{太阳辐射、大气环流、海陆分布、洋流、地形和人类活动}等多因素共同作用的结果。
	
	\textbf{太阳辐射}是根本驱动力，由于地球是球体，不同纬度接收的太阳辐射量不同，形成了从赤道到两极的热量差异，这是全球大气环流和气候带形成的原动力。
	
	\textbf{大气环流}将热量和水分从盈余地区输送到亏损地区，塑造了全球的气候格局。
	
	\textbf{全球六大风带} (极地东风带、中纬西风带和低纬信风带，南、北半球各一套) 的形成源于太阳辐射差异、地球自转和大气运动的复杂相互作用。
	
	\subsubsection{影响地表气温的最基本因素-太阳辐射}
	
	纬度不同，获得的太阳辐射量不同.从低纬到高纬依次形成了热带、亚热带、温带、亚寒带、寒带等不同的气候带。
	\subsubsection{大气环流和洋流-自然界的“空调机”}
	促进了高低纬度之间、海陆之间热量和水分的交换，调节了全球热量和水汽的分布，显著的形象各 地气候。控制各地的气压带、风带不同，形成的气候不同。
	
	暖流对沿岸地区的气候具有增温增湿作用；寒流对沿岸地区具有减温减湿的作用。
	
	\subsubsection{影响地球气候的地表因素}
	\textbf{海陆分布}的热力性质差异导致了季风的形成：夏季风从海洋吹向陆地，带来丰沛降水；冬季风从陆地吹向海洋，寒冷干燥。
	
	\textbf{洋流}如同海洋中的“暖气”和“冷气”管道，暖流使其流经的沿岸地区增温增湿；寒流则使其降温减湿。
	
	\textbf{地形}能阻挡气流运动，导致迎风坡降水丰沛，背风坡干燥少雨，形成“雨影区”。
	
	\textbf{人类活动}通过燃烧化石燃料、改变土地利用方式等，大幅增加了大气中温室气体的浓度，正在以前所未有的速度和规模改变着全球气候系统，导致全球变暖。
	
	\subsubsection{地球反照率 (albedo)}
	
	地球反照率是指地球表面对入射太阳光照射后反射太阳光的能力。这个能力取决于地球表面的物质组成和太阳光照射的角度，地球反照率的平均值大约为30\%。
	
	对太阳辐射的反射率不同，使地面获得的热量有多有少；比热容不同，气温变化有快有慢。
	
	\subsubsection{影响行星气候的主要因素}
	
	\textbf{1. 恒星辐射}：最主要的影响因素。行星接收到的主恒星的能量决定其基本温度。恒星亮度、行星到恒星的距离 (轨道偏心率也会造成季节性变化)，以及行星的自转和倾角，会影响能量在行星表面的分布。
	
	\textbf{2. 大气成分与温室效应}：大气层的存在对气候起着至关重要的作用。某些气体，如水蒸气、二氧化碳、
	甲烷等 (即温室气体)，能够捕获热量，形成温室效应，使行星表面保持温暖。大气压力和密度也影响热
	量的传递和储存。
	
	\textbf{3. 行星反照率}：行星表面和大气反射光的能力 (反照率) 会影响其吸收的能量。例如，冰和云层具有较
	高的反照率，反射更多光线，有冷却作用；而岩石和水吸收更多光线，有加热作用。
	
	\textbf{4. 行星大小和质量}：行星质量决定了其引力，从而影响其能否保持住大气层以及大气层密度。质量和大
	小影响行星内部的热量产生和地质活动 (如火山喷发)，这些活动会将气体释放到大气中，改变其成分。
	
	\textbf{5. 水的存在 (及相态)}：水在气候系统中扮演关键角色。液态水调节温度，水循环 (蒸发、凝结、降水)
	在全球范围内重新分配热量和水分。
	
	\textbf{6. 地质活动和板块构造}：火山活动释放温室气体，影响大气成分；板块构造改变大陆位置和海洋循环模
	式，从而影响长期气候模式。
	
	\textbf{7. 行星自转}：行星的自转速度影响昼夜周期和大气环流模式，例如科里奥利力对风向和洋流的影响。
	
	\subsection{极地气候环境研究}
	极地气候环境研究对于理解全球气候系统、预测气候变化趋势以及制定环境保护政策具有至关重要的意义。极地地区是全球气候变化最敏感、响应最显著的区域。
	
	极地变暖速度远超全球平均水平 (北极约为4倍)，即“极地放大效应”，监测极地冰雪融化、海冰变化等是评估全球变暖程度的关键指标。
	
	\subsubsection{北极是否也有闪电？}
	闪电通常是在暖空气的帮助下形成的，这就是闪电在北极历史上很罕见的原因。
	
	美国地球物理联合会的一次虚拟会议上，研究人员报告了一项新发现，指出北极出现闪电的次数是十年前的数倍，而且这一频率很快就会翻倍。
	\subsubsection{臭氧层破坏的危害}
	到达地面紫外辐射的增加，会破坏人体抗病能力，诱发皮肤癌、麻风、天花等疾病并危害呼吸器官和眼睛。
	
	到达地面紫外辐射增加会引起海洋生物大量死亡，造成某些生物灭绝。
	
	臭氧层中臭氧含量严重减少会使高层和低层大气对太阳辐射能的吸收量发生变化，而太阳辐射能又是大气运动的主要能源，因此臭氧层中臭氧含量的变化对气候也会产生扰动。
	
	\subsubsection{青藏高原隆起的气候效应}
	
	青藏高原的隆起彻底重塑了亚洲气候格局，其核心效应是催生了强大的东亚季风系统，并深刻影响了区域乃至全球的大气环流。
	
	青藏高原是我国重要的生态安全屏障和战略资源基地，其隆起的高海拔大地形被称为世界屋脊，在我国气候变化、水资源供应、生物多样性保护、碳收支平衡等方面具有重要的安全屏障作用。
	
	青藏高原的巨变，造就了今天的江南水乡.
	
	\subsection{古气候研究}
	
	古气候学是对地球过去气候的研究，它利用地球的自然记录来重建百万年来的气候和环境变化。这项研究对于理解当前的气候系统和预测未来趋势至关重要。
	
	\textbf{研究目的}
	
	\textbf{理解自然变率}：区分自然气候周期 (如冰期和间冰期) 和人为造成的变化。
	
	\textbf{测试气候模型}：利用过去的数据验证现代气候模型的准确性。
	
	\textbf{预测未来变化}：深入了解过去的气候系统如何应对温室气体浓度等因素，有助于预测未来的气候变化。
	
	\subsubsection{古气候研究的方法}
	
	古气候学研究主要依靠“代用指标”(Proxy Data)，即保存了过去气候信息的地质、化学和生物材料。
	
	\textbf{冰芯}：在格陵兰岛和南极洲钻取的冰芯包含被困的古代大气气泡。分析这些气泡可以得出过去数十万年二氧化碳、甲烷的浓度和温度记录。
	
	\textbf{海洋和湖泊沉积物}：海洋和湖底的沉积物岩心包含了生物遗骸 (如浮游生物、花粉) 和化学信号，可以揭示过去的海水温度、盐度和植被类型。
	
	\textbf{树木年轮 (树轮学)}： 树木的年轮宽度和密度反映了生长季节的气候条件 (如温度、降水)。
	
	\textbf{珊瑚}：珊瑚骨骼中的化学成分随水温和盐度的变化而变化，是热带地区气候记录的良好来源。洞穴沉积物 (石笋和钟乳石)：洞穴中的碳酸钙沉积物（石笋）包含的氧同位素比例可以反映过去陆地上的降雨量和温度变化。
	
	\subsubsection{北纬30°线}
	北纬30度线是一条贯穿全球、汇集了众多自然与人文奇观的“神奇纬线”，串联起四大文明古国(古埃及、古巴比伦、古印度、中国)与神秘的百慕大三角、金字塔、珠穆朗玛峰、马里亚纳海沟等，同时也是“死亡三角区”和“外星人基地”的传说集中地，蕴含了地球复杂的地质活动、多样的文明信息和未解之谜，被誉为“地球的脐带”。
	\subsubsection{火山活动的主要表现形式}
	全球火山活动的主要表现形式为：熔岩流、火山灰扩散、火山弹抛射、火山气体逸出、火山地震、热异常等。
	
	火山活动有可能造成火山灾害。
	
	\subsection{双碳与全球变化}
	全球变化是由人类活动和自然过程共同驱动，导致地球气候、生态系统等发生的一系列深刻变化。它不仅仅是气候变暖，还包括海平面上升、生物多样性减少、极端天气增多等。
	
	主要驱动力包括自然因素 (如太阳活动、火山喷发) 和人为因素 (如温室气体排放、森林砍伐)。
	
	\subsubsection{地球大气中的甲醛}
	甲醛由于具有大气寿命短、源自非甲烷挥发性有机物（NMVOCs）氧化过程的产率高和大气柱浓度可被卫星探测的特点，其垂直柱浓度被广泛用于表征来自人类活动、自然植被、野火的NMVOCs排放。
	
	但甲醛同时因为其贡献源复杂，反应过程受大气氧化性影响明显、以及卫星反演信号受干扰多等问题，用甲醛表征NMVOCs排放依然具有很大挑战性。
	
	\subsubsection{全球变化是太阳活动引起的吗？}
	太阳常数是指在日地平均距离上，大气顶界垂直于太阳光线单位面积每秒接受的太阳辐射。太阳常数在地球大气层之外，垂直于入射光的平面上测量。人造卫星测得的数值是每平方米大约1367瓦特[$1353(\pm 21) W/m^2$(1976年, NASA)]。
	
	\subsubsection{什么是“碳”？}
	
	“碳”指以二氧化碳为代表的若干种主要的温室气体，具体包括二氧化碳 ($CO_2$)、甲烷($CH_4$)，氧化亚氮($N_2O$)、氢氟碳化物 ($HFCs$)、全氟化碳 ($PFCs$) 和六氟化硫 ($SF_6$)。
	
	这些温室气体的排放将会导致温室效应，对地球的生存环境造成严重影响。
	
	\subsubsection{自然界的碳循环与温室气体的变化趋势}
	
	碳循环是碳元素在地球生物圈、岩石圈、水圈和大气圈之间不断交换和循环的过程。它主要通过
	生物循环和地质循环两个环节进行。
	
	\textbf{生物循环}：绿色植物通过光合作用吸收大气中的二氧化碳，将其转化为有机物储存起来。动物通过食物链摄
	取有机物，并通过呼吸作用将碳以二氧化碳形式释放回大气。动植物残骸被微生物分解，同样释放二氧化碳。
	
	\textbf{地质循环}：部分动植物的残骸经过漫长的地质作用，形成煤、石油等化石燃料。这些燃料被人类开采并燃烧
	后，其中的碳又以二氧化碳形式返回大气。此外，岩石中的碳也会通过风化等作用进入水圈和大气。
	
	\subsubsection{温室气体增加如何毁灭地球环境？}
	碳循环是维持地球生态平衡的关键。然而，人类活动，如大量燃烧化石燃料和破坏
	森林，打破了这一平衡，导致大气中二氧化碳浓度急剧上升，加剧了全球变暖。
	
	\subsubsection{如何理解碳中和？}
	企业、团体或个人测算在一定时间内，直接或间接产生的温室气体排放借由植树造林、节能减排等形式进行抵消，使得最终核算的二氧化碳当量为零，即为实现“碳中和”。
	
	实现这一目标，需要能源转型 (发展风电、光伏等清洁能源)、产业升级(改造钢铁、建材等高耗能行业)、技术创新 (攻关储能、氢能、碳捕集技术) 以及每个人的参与 (如空调调高1度、随手关灯)。
	\subsubsection{全球变暖的应对}
	应对全球变暖，需要从个人行动到国家政策的多层面努力。核心在于减少温室气体排放、增强自然吸碳能力、推动国际合作，并适应已发生的气候变化。
	
	\subsubsection{中国碳中和目标的挑战与前景}
	回顾中国40年改革开放成就，比如说单位GDP的碳排放强度奇迹般的持续下降，离不开以下因素：
	
	在政府领导下，全民一致发展经济；
	
	开放的国际环境和管理、创新技术基础，如太阳能、风能、高铁；
	
	中国人勤劳与智慧的高度发挥。
	
	展望将来，这些条件缺一不可。
	\subsubsection{清洁能源}
	清洁能源，即绿色能源，是指不排放污染物、能够直接用于生产和生活的能源，它包括核能和“可再生能源”。
	
	\subsubsection{中国碳卫星}
	
	碳卫星搭载了一台高空间分辨率的高光谱温室气体探测仪，其工作原理是在可见光和近红外谱段利用分子吸收谱线探测二氧化碳等温室气体浓度。
	
	碳卫星还搭载多谱段的云与气溶胶偏振成像仪。
	
	\subsection{什么是大气科学？}
	大气科学是研究大气的各种现象 (包括人类活动对它的影响)，这些现象的演变规律，以及如何利用这些规律为人类服务的一门学科。
	
	大气科学的分支学科主要有大气探测、气候学、天气学、动力气象学、大气物理学、大气化学、人工影响天气、应用气象学等。
	
	\section{第六章：地震灾害与地震预报预警}
	
	本章重点探讨地震作为一种高破坏性自然灾害的物理本质、度量标准及现代科技在减灾中的预警应用。
	
	\subsection{地震灾害与地震成因}
	
	\subsubsection{1. 地震灾害的破坏现象}
	地震灾害具有突发性强、破坏力巨大的特点。其损失主要来源于以下几个方面：
	\begin{itemize}
		\item \textbf{直接灾害：} 地表破裂、建筑倒塌、桥梁断裂、地基沉降。
		\item \textbf{次生灾害：} 海啸（水下地震）、滑坡、泥石流、火灾（管线破裂）、瘟疫及核泄漏等。
		\item \textbf{损失评估：} 除了直接的人员伤亡和财产损失，还包括长期心理创伤和区域经济发展的停滞。
	\end{itemize}
	
	\subsubsection{2. 地震大小的表示：震级与烈度}
	这是地震学中两个极易混淆但本质不同的概念。
	\begin{enumerate}
		\item \textbf{震级 (Magnitude, $M$):} 
		表示地震释放\textbf{能量}的大小。
		\begin{itemize}
			\item \textbf{定量性：} 一次地震只有一个震级。常用里氏震级或矩震级。
			\item \textbf{能量关系：} 震级每增加 1 级，能量约增加 32 倍；增加 2 级，能量约增加 1000 倍。
			\item \textbf{经验公式：} $\log_{10} E = 4.8 + 1.5M$ (单位：焦耳)。
		\end{itemize}
		
		\item \textbf{烈度 (Intensity, $I$):} 
		表示地震对特定地面造成的\textbf{破坏程度及影响}。
		\begin{itemize}
			\item \textbf{多变性：} 一次地震有多个烈度分布。离震中越近，通常烈度越高。
			\item \textbf{影响因素：} 与震级、震源深度、震中距、地形地貌及建筑抗震性能密切相关。
			\item \textbf{标准：} 中国采用 12 等级烈度表 (I-XII)，通常以人的感觉、物体反应和建筑破坏为依据。
		\end{itemize}
	\end{enumerate}
	\subsubsection{附录：中国地震烈度表 (China Seismic Intensity Scale)}
	
	为了定量评估地震对地面及建筑的影响，我国采用《中国地震烈度表》。该表将地震影响分为 12 个等级（用罗马数字 I--XII 表示）：
	
	\begin{table}[htbp]
		\centering
		\small
		\renewcommand{\arraystretch}{1.3} % 调节行高
		\begin{tabularx}{\textwidth}{|c|X|X|X|}
			\hline
			\textbf{烈度} & \textbf{在地面上人的感觉} & \textbf{房屋震害程度} & \textbf{其他震害现象} \\ \hline
			\textbf{I} & 无感 & & \\ \hline
			\textbf{II} & 室内个别静止中人有感觉 & & \\ \hline
			\textbf{III} & 室内少数静止中人有感觉 & 门、窗轻微作响 & 悬挂物微动 \\ \hline
			\textbf{IV} & 室内多数人、室外少数人有感觉，少数人梦中惊醒 & 门、窗作响 & 悬挂物明显摆动，器皿作响 \\ \hline
			\textbf{V} & 室内普遍、室外多数人有感觉，多数人梦中惊醒 & 门窗、屋顶、屋架颤动应作响，抹灰出现微细裂缝，个别烟囱掉砖 & 不稳定器物摇动或翻倒 \\ \hline
			\textbf{VI} & 多数人站立不稳，少数人惊逃户外 & 损坏——墙体出现裂缝，檐瓦掉落，少数屋顶烟囱裂缝、掉落 & 河岸和松软土出现裂缝，饱和砂层出现喷砂冒水 \\ \hline
			\textbf{VII} & 大多数人惊逃户外，骑自行车的人有感觉，行驶中的汽车驾驶员有感觉 & 轻度破坏——局部破坏，开裂，小修或不需要修理可继续使用 & 河岸出现坍方；饱和砂层常见喷砂冒水；大多数独立砖烟囱中等破坏 \\ \hline
			\textbf{VIII} & 多数人摇晃颠簸，行走困难 & 中等破坏——结构破坏，需要修复才能使用 & 干硬土上亦出现裂缝；房屋破坏导致人畜伤亡 \\ \hline
			\textbf{IX} & 行动的人摔倒 & 严重破坏——结构严重破坏，局部倒塌，修复困难 & 干硬土上出现地方有裂缝；滑坡坍方；砖烟囱倒塌 \\ \hline
			\textbf{X} & 骑自行车的摔倒，有抛起感 & 大多数倒塌 & 山崩和地震断裂出现；拱桥破坏 \\ \hline
			\textbf{XI} & & 普遍倒塌 & 地震断裂延续很长；大量山崩滑坡 \\ \hline
			\textbf{XII} & & & 地面剧烈变化，山河改观 \\ \hline
		\end{tabularx}
		\caption{中国地震烈度简表（根据图像内容整理）}
	\end{table}
	
	\subsubsection{地震成因深究：弹性回跳模型 (The Elastic Rebound Theory)}
	
	弹性回跳模型是解释构造地震发生机制最经典、最受认可的物理模型。它将地震视为地壳中能量积累与瞬间释放的循环过程。
	
	\begin{enumerate}
		\item \textbf{背景与提出}
		\begin{itemize}
			\item \textbf{提出者：} 哈里·菲尔丁·里德 (Harry Fielding Reid)。
			\item \textbf{观测依据：} 通过对 1906 年旧金山大地震前后圣安德烈斯断层两侧大地测量数据的对比，发现地表发生了显著的水平错动。
		\end{itemize}
		
		\item \textbf{物理过程详解}
		该模型将地震过程分为四个连续的物理阶段：
		\begin{itemize}
			\item \textbf{应力积累阶段 (Stress Accumulation)：} 由于板块运动，地壳深部的构造应力持续作用于断层两侧的岩石。
			\item \textbf{弹性形变阶段 (Elastic Deformation)：} 断层线受到摩擦力的“锁死”作用。岩石在应力下不再相对滑动，而是像巨大的弹簧一样发生弯曲和变形，此时电磁能和机械能转化为\textbf{弹性势能}储存在岩层中。
			\item \textbf{断裂启动阶段 (Rupture Initiation)：} 当应力不断增加，最终超过了断层面上岩石的摩擦阻力或岩石自身的抗剪强度时，岩石在最脆弱处（震源）发生断裂。
			\item \textbf{弹性回跳阶段 (Rebound)：} 断层两侧已经发生形变的岩石像被剪断的皮筋一样，迅速向相反方向弹回，恢复到原来的无应变状态（或较低应变状态）。
		\end{itemize}
		
		\item \textbf{能量转换关系}
		在回跳的瞬间，储存的巨大弹性势能发生剧烈转换：
		\begin{equation}
			E_{\text{potential}} \rightarrow E_{\text{seismic}} (\text{地震波动能}) + E_{\text{heat}} (\text{摩擦生热}) + E_{\text{fracture}} (\text{破碎能})
		\end{equation}
		其中，$E_{\text{seismic}}$ 以纵波、横波和面波的形式向四周传播，造成地表震动。
		
		\item \textbf{模型的科学意义}
		\begin{itemize}
			\item \textbf{地震循环模型：} 该理论确立了“应力积累—破裂—再积累”的地震循环概念，是长期地震预测的理论基础。
			\item \textbf{解释震级差异：} 断层面积越大、岩石弹性越好、锁死时间越长，积累的能量就越多，释放时的震级就越高。
			\item \textbf{局限性：} 虽然该模型完美解释了地表附近的浅源构造地震，但对于数十公里深处的深源地震（高温高压下岩石表现为塑性）以及流体诱发地震，其解释力仍存在争议。
		\end{itemize}
		
		\item \textbf{形象比喻：}
		可以将弹性回跳想象成**折断一把直尺**的过程：双手缓慢用力使直尺弯曲（积累势能），直尺形变越来越大（弹性形变），当力超过直尺极限时，直尺“咔嚓”一声断掉（断裂），断开的两截迅速弹回原位并伴随震动和声响（能量释放）。
	\end{enumerate}
	
	\subsubsection{1. 地震三要素}
	描述一次地震的基本参数包括：
	\begin{itemize}
		\item \textbf{地震时刻：} 地震发生的具体时刻。
		\item \textbf{震中位置：} 包括经纬度及震源深度。
		\item \textbf{震级：} 释放能量的大小。
	\end{itemize}
	
	\subsubsection{2. 地震预报 (Earthquake Prediction)}
	地震预报是指在地震发生\textbf{之前}，对未来地震发生的时间、地点和震级的预测。
	\begin{itemize}
		\item \textbf{分类：} 长期（十年尺度）、中期（一至两年）、短期（三个月内）和临震预报（几天至几小时内）。
		\item \textbf{现状：} 这是一个世界级科学难题。由于地壳深部的不可入性、地震过程的高度非线性和复杂性，目前尚无法做到精准的临震预报。
	\end{itemize}
	
	\subsubsection{3. 地震预警 (Earthquake Early Warning, EEW)}
	地震预警是指在地震发生\textbf{之后}，利用电波比地震波快的原理，在破坏性地震波到达目标区之前发出的警报。
	\begin{itemize}
		\item \textbf{物理原理：} 
		\begin{itemize}
			\item 纵波 ($P$ 波) 传播速度快（约 $6 \text{ km/s}$）但破坏小；横波 ($S$ 波) 速度慢（约 $3.5 \text{ km/s}$）但破坏大。
			\item 电波传播速度为光速 ($3 \times 10^5 \text{ km/s}$)。
		\end{itemize}
		\item \textbf{预警时间 ($t$) 计算公式：}
		\begin{equation}
			t = D \left( \frac{1}{V_s} - \frac{1}{V_p} \right) - t_{\text{processing}}
		\end{equation}
		其中 $D$ 为震中距，$V_s, V_p$ 为波速，$t_{\text{processing}}$ 为系统反应时间。
		\item \textbf{意义：} 虽然只有几秒到几十秒，但足以让高速列车减速、关闭燃气阀门、切断电源、或人员紧急避险，从而大幅减少伤亡。
	\end{itemize}
	
	\subsubsection{4. 总结：预报与预警的区别}
	\begin{itemize}
		\item \textbf{时间差：} 预报是在“事发前”，预警是在“事发后，受灾前”。
		\item \textbf{实现难度：} 预报极难，处于探索阶段；预警技术成熟，已在全球多个地震带（如中国四川、日本）广泛部署。
	\end{itemize}
	
	
	
	
	
	
\end{document}